% ============================================================================
% MÉMOIRE — Allocation de capital sous contraintes réglementaires
% Document maître : compiler ce fichier uniquement.
% ============================================================================
\documentclass[12pt, a4paper, twoside, openright]{report}

% ─── Encodage & langue ────────────────────────────────────
\usepackage[utf8]{inputenc}
\usepackage[T1]{fontenc}
\usepackage[french]{babel}
\usepackage{csquotes}

% ─── Police ───────────────────────────────────────────────
\usepackage{mathpazo}          % Palatino (serif, lisible, académique)
\usepackage[scaled=0.85]{helvet}
\usepackage{microtype}         % Améliorations typographiques fines

% ─── Couleurs ────────────────────────────────────────────
\usepackage{xcolor}
\definecolor{bleutitre}{RGB}{0, 51, 102}     % bleu marine sobre
\definecolor{bleulien}{RGB}{0, 70, 130}      % bleu légèrement plus clair

% ─── Géométrie ────────────────────────────────────────────
\usepackage{geometry}
\geometry{
  top=2.5cm,
  bottom=2.5cm,
  inner=3cm,
  outer=2.5cm,
  headheight=15pt,
}
\usepackage{setspace}
\onehalfspacing

% ─── Mathématiques ────────────────────────────────────────
\usepackage{amsmath, amssymb, amsthm}
\usepackage{mathtools}

% ─── Tableaux & listes ───────────────────────────────────
\usepackage{booktabs}
\usepackage{multirow}
\usepackage{array}
\usepackage{longtable}
\usepackage{enumitem}
\usepackage{textcomp}
\usepackage[most]{tcolorbox}

% ─── Graphiques ───────────────────────────────────────────
\usepackage{graphicx}
\usepackage{caption}
\usepackage{subcaption}

% ─── Algorithmes ──────────────────────────────────────────
\usepackage{algorithm}
\usepackage{algpseudocode}

% ─── En-têtes et pieds de page ────────────────────────────
\usepackage{fancyhdr}
\pagestyle{fancy}
\fancyhf{}
\fancyhead[LE]{\small\itshape\leftmark}
\fancyhead[RO]{\small\itshape\rightmark}
\fancyfoot[C]{\thepage}
\renewcommand{\headrulewidth}{0.4pt}

\fancypagestyle{plain}{%
  \fancyhf{}
  \fancyfoot[C]{\thepage}
  \renewcommand{\headrulewidth}{0pt}
}

% ─── Liens ────────────────────────────────────────────────
\usepackage{hyperref}
\hypersetup{
  colorlinks=true,
  linkcolor=bleutitre,
  citecolor=bleulien,
  urlcolor=bleulien,
  pdftitle={Allocation de capital sous contraintes réglementaires},
  pdfauthor={},
}

% ─── Bibliographie ────────────────────────────────────────
\usepackage[authoryear, round]{natbib}

% ─── Titres colorés ─────────────────────────────────────
\usepackage{titlesec}
\titleformat{\part}[display]
  {\centering\Huge\bfseries\color{bleutitre}}{\partname~\thepart}{20pt}{\Huge}
\titleformat{\chapter}[display]
  {\bfseries\Large\color{bleutitre}}{\chaptertitlename~\thechapter}{15pt}{\LARGE}
\titleformat{\section}
  {\bfseries\large\color{bleutitre}}{\thesection}{1em}{}
\titleformat{\subsection}
  {\bfseries\normalsize\color{bleutitre}}{\thesubsection}{1em}{}

% ─── Théorèmes & définitions ─────────────────────────────
\theoremstyle{definition}
\newtheorem{definition}{Définition}[chapter]
\newtheorem{proposition}{Proposition}[chapter]
\newtheorem{remarque}{Remarque}[chapter]

% ============================================================================
\begin{document}

% ── Page de titre ─────────────────────────────────────────
\begin{titlepage}
\centering
\vspace*{3cm}

{\Huge\bfseries Allocation de capital\\[0.4em]
sous contraintes réglementaires}

\vspace{1.5cm}

{\LARGE IFRS~9, IFRS~13, Bâle~III\,/\,CRR3}

\vspace{3cm}

\rule{0.5\textwidth}{0.5pt}

\vspace{2cm}

{\Large Mémoire}

\vspace{1.5cm}

{\large
\textbf{Domaine~:} Gestion des risques financiers\\[0.5em]
\textbf{Méthode~:} Data science et optimisation sous contraintes
}

\vfill

{\large 2026}

\end{titlepage}

% ── Table des matières ────────────────────────────────────
\tableofcontents
\newpage

% ══════════════════════════════════════════════════════════
% INTRODUCTION (hors parties)
% ══════════════════════════════════════════════════════════

\chapter*{Introduction}
\addcontentsline{toc}{chapter}{Introduction}
\markboth{Introduction}{Introduction}
\label{chap:introduction}

% ── MOVE 1 : Établir le territoire ─────────────────────────

La crise financière de 2008 a révélé que les banques pouvaient détenir des portefeuilles individuellement acceptables mais collectivement fragiles.
En réponse, le législateur a profondément remanié le cadre normatif~: IFRS~9 \citep{ifrs9} a remplacé le provisionnement rétrospectif d'IAS~39 par un modèle prospectif en trois stades~; CRR3 \citep{crr3} a durci les pondérations de risque, en particulier sur les expositions en actions et en private equity~; Bâle~III \citep{bcbs_d424} a introduit des ratios de liquidité (LCR, NSFR) et un encadrement du risque de taux (IRRBB).
Ces réformes ont rendu le bilan bancaire plus résilient, mais aussi plus contraint~: chaque classe d'actifs obéit désormais à un corpus réglementaire propre, et l'allocation du capital entre elles est devenue un problème d'optimisation sous contraintes multiples.

La littérature sur l'allocation de capital bancaire est abondante mais fragmentée.
Les travaux fondateurs de \citet{merton1974} et \citet{vasicek2002} ont posé les bases de la modélisation du risque de crédit en portefeuille, prolongées par \citet{gordy2003} dans le cadre du modèle ASRF qui sous-tend les exigences de capital de Bâle~II et~III.
L'estimation des paramètres de risque (PD, LGD, EAD) a fait l'objet d'un cadre réglementaire détaillé \citep{eba_gl_2017_06, eba_gl_2019_03} et d'une littérature appliquée sur les modèles de scoring \citep{siddiqi2006, hosmer_lemeshow2013}.
Plus récemment, les méthodes d'apprentissage automatique --- TabNet \citep{arik_pfister2021}, XGBoost \citep{chen_guestrin2016} --- ont démontré leur capacité à capturer des non-linéarités que les modèles traditionnels ignorent \citep{grinsztajn2022}.

Du côté de l'optimisation de portefeuille, le modèle de \citet{black_litterman1992} permet de combiner un équilibre de marché avec des vues subjectives, tandis que la CVaR \citep{rockafellar_uryasev2000} offre une mesure de risque cohérente adaptée aux distributions à queue épaisse.
L'allocation de capital au sein d'une banque a été formalisée par \citet{tasche2008} via le principe d'Euler, et le cadre général du risque quantitatif est couvert par \citet{mcneil_frey_embrechts2015}.

Enfin, la valorisation du private equity en juste valeur (IFRS~13) fait l'objet de guidelines spécifiques \citep{ipev2022}, et les données de marché --- décotes de liquidité \citep{jefferies_lazard2023}, performance par vintage \citep{cambridge2023}, taux de défaut par notation \citep{moodys_default2023} --- permettent de calibrer les modèles.

% ── MOVE 2 : Établir la niche (gap) ────────────────────────

Toutefois, ces travaux abordent chaque dimension du problème de manière isolée.
Les modèles de risque de crédit ne traitent pas de la valorisation PE.
Les méthodes d'optimisation de portefeuille ne modélisent pas le staging IFRS~9.
Les guidelines de provisionnement n'intègrent pas les contraintes de liquidité.
En pratique, le provisionnement crédit, la valorisation PE, l'optimisation du capital et le suivi de la liquidité sont réalisés par des équipes distinctes, avec des outils distincts, et remontent au comité des risques comme des rapports séparés.

Il n'existe pas, à notre connaissance, de cadre intégré qui permette d'optimiser l'allocation du capital d'une banque sous l'ensemble des contraintes réglementaires --- IFRS~9, IFRS~13, CRR3, LCR, NSFR, IRRBB --- en modélisant conjointement toutes les classes d'actifs face aux mêmes chocs macroéconomiques.

% ── MOVE 3 : Occuper la niche ──────────────────────────────

Ce mémoire se propose de combler cette lacune en répondant à la question suivante~:

\medskip
\begin{quote}
\itshape
Sous les contraintes conjointes d'IFRS~9, d'IFRS~13 et de Bâle~III, quelle est l'allocation optimale du capital d'un établissement financier entre ses classes d'actifs~?
\end{quote}
\medskip

Pour y répondre, nous construisons un cadre d'analyse intégré et l'appliquons à deux échelles successives.
Ce travail apporte quatre contributions~:

\begin{enumerate}
  \item \textbf{Un cadre unifié multi-normes.}
  Nous proposons un pipeline qui modélise conjointement le provisionnement IFRS~9, la valorisation IFRS~13, la consommation de capital CRR3, et les ratios de liquidité (LCR, NSFR), alimenté par les mêmes variables macroéconomiques et testé sur douze scénarios historiquement calibrés.

  \item \textbf{Des résultats sur la nature de l'arbitrage.}
  En partant du cas crédit contre PE --- deux classes aux profils de risque opposés --- puis en étendant l'analyse à quatorze classes d'actifs via un optimiseur Black-Litterman CVaR, nous identifions où se situe réellement le levier d'allocation.

  \item \textbf{Une découverte sur les déterminants de l'allocation.}
  L'analyse systématique de 96 configurations de modélisation révèle qu'un paramètre rarement discuté dans la littérature a un impact de premier ordre sur l'allocation --- plus que l'algorithme d'optimisation lui-même.

  \item \textbf{Un outillage complet.}
  Un pipeline intégré --- du générateur de données synthétiques au tableau de bord CRO --- incluant trois familles de modèles PD, dix moteurs de gouvernance, et un analyste neuro-symbolique.
\end{enumerate}

L'approche repose sur des données synthétiques calibrées sur des sources institutionnelles (EBA, Moody's, Cambridge Associates), trois familles de modèles de probabilité de défaut (régression logistique, TabNet, XGBoost), un optimiseur multi-contraintes, et une batterie de dix moteurs de gouvernance pour la validation.
Le tout est intégré dans un tableau de bord interactif destiné au directeur des risques.

Le mémoire est organisé en deux parties.
La \textbf{première partie} part du cas crédit/PE sur cinq secteurs économiques et construit l'ensemble de l'outillage analytique~: génération de données, modélisation du risque de défaut, moteur de provisionnement IFRS~9, métriques de portefeuille, et analyse du private equity sous contraintes réglementaires.
La \textbf{seconde partie} passe à l'échelle d'un bilan complet --- quatorze classes d'actifs, optimisation multi-contraintes, et analyse systématique des déterminants de l'allocation.


% ══════════════════════════════════════════════════════════
% PARTIE I
% ══════════════════════════════════════════════════════════
\part{Fondation sectorielle}
\label{part:fondation}

% ──────────────────────────────────────────────────────────
% Chapitre 1 — Données synthétiques et processus générateur
% ──────────────────────────────────────────────────────────

\chapter{Données synthétiques et processus générateur}
\label{chap:donnees}

Ce chapitre présente l'architecture du générateur de données synthétiques qui sous-tend l'ensemble de l'analyse.

Le recours à des données synthétiques est un choix méthodologique délibéré, motivé par trois raisons.
Premièrement, les données bancaires réelles sont confidentielles~: aucun portefeuille de crédit d'une banque européenne n'est librement disponible à la granularité nécessaire (30\,000 contreparties, 15 variables financières, historique de retard de paiement).
Deuxièmement, un processus générateur de données (DGP) contrôlé permet de connaître \emph{exactement} les relations entre variables et défaut, ce qui autorise un diagnostic précis des forces et faiblesses de chaque modèle de scoring --- une propriété impossible avec des données réelles où le vrai DGP est inconnu.
Troisièmement, les scénarios de stress extrêmes (Stagflation, GFC, COVID) n'apparaissent qu'une fois dans les données historiques~; un générateur permet de les reproduire à volonté.

La contrepartie de ce choix est que les amplitudes des résultats dépendent de la calibration du générateur.
Pour limiter cet arbitraire, tous les paramètres sont calibrés sur des sources institutionnelles publiques (EBA, Moody's, Cambridge Associates, Eurostat, BCE), et nous documentons systématiquement l'origine de chaque choix.


% ════════════════════════════════════════════════════════
\section{Architecture générale}
\label{sec:donnees:architecture}

Le générateur produit un quadruplet de jeux de données~:
%
\begin{equation}
\label{eq:generate_dataset}
\texttt{generate\_dataset}(n, s) \;\to\; \bigl(\mathbf{D}_{\text{credit}},\; \mathbf{D}_{\text{PE}},\; \mathbf{D}_{\text{hist}},\; \mathbf{D}_{\text{bilan}}\bigr)
\end{equation}
%
où $n = 30\,000$ est le nombre de clients crédit et $s = 123$ la graine aléatoire.

Le choix de $n = 30\,000$ résulte d'un compromis~: c'est suffisant pour que les taux de défaut sectoriels (1{,}2 à 2{,}5\,\%) produisent un nombre de défauts statistiquement exploitable (360 à 750 défauts par exécution), tout en restant rapide à générer ($< 2$ secondes).

Chaque composante du quadruplet a un rôle distinct~:
\begin{itemize}[nosep]
  \item $\mathbf{D}_{\text{credit}}$~: caractéristiques financières, secteur, PD d'origination et flag de défaut --- alimente les modèles de scoring (chapitre~2)~;
  \item $\mathbf{D}_{\text{PE}}$~: fonds de private equity associés (vintage, NAV, levier, multiples) --- alimente la valorisation IFRS~13 (chapitre~5)~;
  \item $\mathbf{D}_{\text{hist}}$~: historique mensuel (soldes, retards de paiement, variables macro) --- alimente le staging IFRS~9 (chapitre~3)~;
  \item $\mathbf{D}_{\text{bilan}}$~: bilan agrégé par classe d'actifs (EAD, PD, LGD, RW) --- alimente les métriques de portefeuille (chapitre~4).
\end{itemize}

La graine aléatoire est déclinée par module~: $s + 50$ pour les variables causales, $s + 100$ pour les caractéristiques crédit, $s + 200$ pour le PE, $s + 300$ pour l'historique.
Ce découpage garantit que la modification d'un module n'altère pas les tirages des autres --- propriété essentielle pour isoler l'effet d'un changement lors de l'analyse de sensibilité.


% ════════════════════════════════════════════════════════
\section{Univers sectoriel}
\label{sec:donnees:secteurs}

Le portefeuille couvre cinq secteurs économiques.
Le choix de cinq secteurs (et non dix ou vingt) est dicté par le modèle ASRF mono-facteur de \citet{gordy2003} qui sous-tend Bâle~II/III~: dans ce cadre, le risque systématique est capturé par un facteur unique, et la granularité sectorielle sert principalement à différencier les sensibilités macroéconomiques.
Cinq secteurs suffisent pour couvrir les principaux profils de risque (cyclique, défensif, taux-sensible, immobilier, services) sans introduire de redondance.

\begin{table}[htbp]
\centering
\caption{Paramètres sectoriels~: pondérations, taux de défaut et taux zombie.}
\label{tab:secteurs}
\begin{tabular}{lcccp{5cm}}
\toprule
Secteur & Poids & $\text{PD}_{\text{base}}$ & Zombie & Justification \\
\midrule
Technologie   & 25\,\% & 2{,}5\,\% & 25\,\% & Taux NPL EBA 2--3\,\% en régime normal~; entreprises à forte croissance mais fragiles (burn rate élevé des start-ups) \\
Industrie     & 20\,\% & 2{,}0\,\% & 18\,\% & Milieu de spectre Moody's (manufacturing)~; levier opérationnel élevé, cyclicité modérée \\
Santé         & 12\,\% & 1{,}2\,\% &  5\,\% & Taux NPL Moody's 0{,}5--1{,}5\,\% (revenus régulés, demande inélastique)~; très peu d'entreprises zombies grâce au remboursement public \\
Immobilier    & 28\,\% & 2{,}2\,\% & 30\,\% & Taux NPL ECB/EBA hypothécaire 1{,}5--3{,}0\,\% post-crise~; forte proportion de promoteurs à marge négative (30\,\% en cycle baissier) \\
Services      & 15\,\% & 1{,}8\,\% & 22\,\% & Taux NPL EBA PME-services 1{,}5--2{,}5\,\%~; pression salariale en régime inflationniste \\
\bottomrule
\end{tabular}
\end{table}

Les pondérations reflètent la structure d'un portefeuille corporate européen diversifié.
L'Immobilier à 28\,\% est cohérent avec la part du crédit immobilier dans les bilans des banques françaises (25--30\,\% selon l'ACPR).
La Technologie à 25\,\% est volontairement surpondérée par rapport à la réalité ($\sim$15\,\%) pour tester la robustesse du modèle face à un secteur à forte dispersion.

Le \emph{taux zombie} --- proportion cible d'entreprises à marge EBITDA négative --- est un paramètre structurant du DGP.
Il détermine la bimodalité de la distribution des marges et, par conséquent, la difficulté du problème de classification.
Les valeurs retenues sont cohérentes avec les estimations de la BRI~: 12\,\% de firmes zombies en zone euro en 2019 \citep{bcbs_d424}, avec une forte hétérogénéité sectorielle (Santé 5\,\%, Immobilier 30\,\%).

\subsection{Sensibilités macroéconomiques}

Chaque secteur est paramétré par un vecteur de sensibilités $\mathbf{s}_k = (s_{\text{chômage}}, s_{\text{PIB}}, s_{\text{taux}}, s_{\text{HPI}}, s_{\text{inflation}})$.
Ces sensibilités déterminent comment un choc macroéconomique se transmet à la PD conditionnelle via le modèle de Vasicek (chapitre~3).

\begin{table}[htbp]
\centering
\caption{Sensibilités macroéconomiques par secteur (canal crédit) et justification.}
\label{tab:sensibilites_credit}
\begin{tabular}{lccccp{4.5cm}}
\toprule
Secteur & Chôm. & PIB & Taux & HPI & Justification dominante \\
\midrule
Technologie & 2{,}5 & 1{,}8 & 1{,}2 & 0{,}5 & Chômage $\to$ coupes budgets IT $\to$ churn SaaS $\to$ défaut \\
Industrie   & 1{,}5 & 2{,}0 & 1{,}0 & 0{,}8 & PIB dominant (demande cyclique d'équipements) \\
Santé       & 0{,}5 & 0{,}4 & 0{,}6 & 0{,}3 & Défensif~: demande inélastique, revenus publics \\
Immobilier  & 1{,}0 & 1{,}2 & 2{,}0 & 2{,}5 & HPI dominant (LTV, collatéral)~; taux via service de dette \\
Services    & 1{,}2 & 1{,}0 & 0{,}8 & 0{,}5 & Inflation dominante ($s_{\text{infl}} = 1{,}8$, pression salariale) \\
\bottomrule
\end{tabular}
\end{table}

La logique de calibration est la suivante~:
\begin{itemize}[nosep]
  \item \textbf{Immobilier, HPI = 2{,}5}~: c'est la sensibilité la plus élevée du portefeuille. Une baisse de 10\,\% du HPI dégrade mécaniquement le LTV (ratio prêt/valeur), ce qui déclenche des reclassements en Stage~2 IFRS~9 et augmente la LGD \textit{downturn}. Les données Eurostat confirment que les prix immobiliers en zone euro ont reculé de $-3{,}2\,\%$ en 2009 et de $-2{,}5\,\%$ en 2012, avec des effets amplifiés sur les promoteurs (LTV~$> 80\,\%$).
  \item \textbf{Technologie, chômage = 2{,}5}~: le mécanisme est indirect mais puissant~: le chômage réduit les budgets IT des entreprises clientes (poste compressible en premier), ce qui entraîne un \emph{churn} sur les abonnements SaaS et menace la viabilité des start-ups en phase de croissance.
  \item \textbf{Santé, toutes sensibilités $< 1$}~: ce secteur est structurellement protégé par le remboursement public (sécurité sociale, assurance maladie), ce qui découple sa performance des cycles macroéconomiques. Les données Moody's montrent un taux de défaut inférieur à 1{,}5\,\% même en 2009.
\end{itemize}


% ════════════════════════════════════════════════════════
\section{Structure de dépendance~: copule de Student}
\label{sec:donnees:copule}

Les variables financières d'une même entreprise ne sont pas indépendantes~: une firme fortement endettée a généralement une marge faible, un ratio de couverture des intérêts bas, et une volatilité élevée.
Pour reproduire cette structure de dépendance sans imposer la normalité jointe, nous utilisons une copule de Student sectorielle.

\subsection{Pourquoi une copule de Student ?}

La copule gaussienne, largement utilisée dans les modèles de risque de crédit \citep{vasicek2002}, sous-estime la dépendance de queue~: elle suppose que les co-mouvements extrêmes sont aussi rares que le prévoit la loi normale bivariée.
Or, les crises financières montrent que les corrélations augmentent en période de stress --- phénomène que \citet{mcneil_frey_embrechts2015} documentent abondamment.
La copule de Student, paramétrée par un nombre de degrés de liberté $\nu$, permet de contrôler explicitement l'épaisseur des queues jointes.

\subsection{Modèle factoriel}

La matrice de corrélation est construite via un modèle à quatre facteurs latents (profitabilité, levier, liquidité, maturité)~:
%
\begin{equation}
\label{eq:copula_corr}
\boldsymbol{\Sigma} = \mathbf{L}\mathbf{L}^\top + \boldsymbol{\Psi}, \quad \text{où} \quad \Psi_{ii} = 1 - h_i^2, \quad h_i^2 = \sum_{j=1}^{4} L_{ij}^2
\end{equation}
%
où $\mathbf{L} \in \mathbb{R}^{15 \times 4}$ est la matrice de charges factorielles.
Les communalités $h_i^2$ sont plafonnées à 0{,}95~: au-delà, la variance spécifique $\Psi_{ii}$ deviendrait quasi-nulle, rendant certaines variables presque déterministes.

Chaque secteur modifie cette structure via un vecteur d'échelle $\mathbf{s}_k \in \mathbb{R}^4$~:
%
\begin{equation}
\mathbf{L}_k = \mathbf{L} \odot \mathbf{s}_k^\top
\end{equation}
%
La Technologie, par exemple, amplifie le facteur profitabilité ($s_1 = 1{,}20$) et atténue le levier ($s_2 = 0{,}70$), reflétant le modèle \emph{asset-light} de ce secteur (marges élevées, peu de dette physique).
L'Énergie fait l'inverse ($s_2 = 1{,}30$), traduisant l'intensité capitalistique des projets extractifs.

\subsection{Échantillonnage et régime de stress}

L'échantillonnage suit le schéma standard de la copule de Student~:
%
\begin{equation}
\label{eq:copula_sampling}
\mathbf{Z} \sim \mathcal{N}_d(\mathbf{0}, \mathbf{C}), \quad
W \sim \chi^2(\nu)/\nu, \quad
\mathbf{X}_t = \mathbf{Z} / \sqrt{W}, \quad
\mathbf{U} = F_t(\mathbf{X}_t; \nu)
\end{equation}
%
Le choix de $\nu = 5$ en régime normal suit la recommandation de \citet{mcneil_frey_embrechts2015}, qui préconisent $\nu \in [4, 8]$ pour les portefeuilles de crédit corporate --- assez bas pour capturer les co-mouvements de queue, assez haut pour éviter une variance infinie ($\nu > 4$).

En période de stress (trimestres T9--T11 du cycle), deux modifications s'appliquent~:
\begin{itemize}[nosep]
  \item $\nu$ passe de 5 à 3 (queues plus lourdes, variance finie car $\nu > 2$)~;
  \item la corrélation est amplifiée par un facteur $\lambda_s = 1{,}3$~:
\end{itemize}
%
\begin{equation}
\mathbf{C}_{\text{stress}} = \mathbf{C}_{\text{normal}} + \lambda_s \bigl(\mathbf{C}_{\text{normal}} - \mathbf{I}\bigr)
\end{equation}
%
Ce facteur $\lambda_s = 1{,}3$ est calibré pour que la corrélation moyenne hors diagonale passe d'environ 0{,}25 à 0{,}33 en stress --- un ordre de grandeur cohérent avec l'augmentation observée des corrélations de crédit pendant la crise de 2008 \citep{mcneil_frey_embrechts2015}.


% ════════════════════════════════════════════════════════
\section{Distributions marginales}
\label{sec:donnees:marginales}

Les uniformes de la copule sont transformées en variables financières via la méthode PIT (Probability Integral Transform).
Le choix de chaque distribution marginale est guidé par les propriétés statistiques de la variable~:

\begin{table}[htbp]
\centering
\caption{Distributions marginales et justification du choix.}
\label{tab:marginales}
\begin{tabular}{llp{6.5cm}}
\toprule
Variable & Distribution & Justification \\
\midrule
Chiffre d'affaires      & Log-normale & Toujours positif, fortement asymétrique à droite (loi de puissance des tailles d'entreprises) \\
Marge EBITDA            & Mélange gaussien & Bimodale~: 80\,\% d'entreprises saines ($\mu = +15\,\%$) et 20\,\% de zombies ($\mu = -5\,\%$) \\
Ratio d'endettement     & Gamma ($\alpha\!=\!2$, $\beta\!=\!0{,}3$) & Positif, mode à $\sim$0{,}30 (zone de confort), queue droite vers les LBO \\
Volatilité des CF       & Bêta ($\alpha\!=\!1{,}5$, $\beta\!=\!4$) & Bornée $[0, 1]$, mode $\sim$0{,}10~; queue droite pour les firmes spéculatives \\
Couverture des intérêts  & Log-normale ($\sigma\!=\!0{,}7$) & Toujours positive, médiane ${\sim}\,3{,}0\times$ (seuil BBB), forte asymétrie \\
Liquidité courante      & Gamma ($\alpha\!=\!5$, $\beta\!=\!0{,}3$) & Mode $\sim$1{,}2 (au-dessus du seuil critique de 1{,}0), queue droite modérée \\
Loan-to-value           & Bêta ($\alpha\!=\!2{,}5$, $\beta\!=\!3$) & Bornée $[0, 1]$, mode $\sim$0{,}45 (portefeuille européen moyen, ECB MFI) \\
\bottomrule
\end{tabular}
\end{table}

Le mélange bimodal de la marge EBITDA mérite une explication.
Dans un portefeuille bancaire réel, la distribution des marges n'est pas unimodale~: elle présente un pic principal autour de 10--20\,\% (entreprises saines) et un second mode négatif (entreprises en difficulté qui brûlent du cash).
La proportion de 20\,\% de la composante zombie correspond à la valeur de base avant ajustement sectoriel~; le tableau~\ref{tab:secteurs} montre que ce taux varie de 5\,\% (Santé) à 30\,\% (Immobilier).

\subsection{Ajustements sectoriels}

Après transformation marginale, chaque secteur applique un décalage de localisation pour différencier les profils.
Ces décalages reflètent les caractéristiques économiques de chaque secteur~:
\begin{itemize}[nosep]
  \item \textbf{Technologie}~: $+0{,}05$ sur la marge EBITDA (modèle SaaS à forte marge) et $-0{,}15$ sur les actifs tangibles (modèle \emph{asset-light}, peu d'immobilisations corporelles).
  \item \textbf{Industrie}~: $+0{,}10$ sur les actifs tangibles (usines, équipements) et $+0{,}05$ sur le ratio d'endettement (financement d'actifs physiques par la dette).
  \item \textbf{Énergie}~: $+0{,}08$ sur le ratio d'endettement et $+0{,}06$ sur la volatilité des flux de trésorerie --- reflétant l'intensité capitalistique et la dépendance aux prix des matières premières.
\end{itemize}


% ════════════════════════════════════════════════════════
\section{Processus générateur de défaut}
\label{sec:donnees:dgp}

Le DGP est construit en trois blocs additifs dont la répartition du signal est volontairement structurée pour tester les capacités discriminantes de chaque famille de modèles.
Cette architecture s'inspire du cadre de \citet{grinsztajn2022}, qui montrent que l'avantage des modèles à base d'arbres sur les données tabulaires provient précisément de leur capacité à capturer les interactions et les effets non linéaires.

\subsection{Bloc 1~: effets monotones par paliers ($\sim$40\,\% du signal)}

Le premier bloc capture les relations univariées entre chaque variable et le risque de défaut.
Les seuils sont choisis pour correspondre aux conventions bancaires~:

\begin{itemize}[nosep]
  \item \textbf{Ratio d'endettement}~: le seuil bas de 0{,}30 correspond à la zone de confort des covenants bancaires~; le seuil haut de 0{,}60 marque la zone de détresse (au-delà, le service de la dette consomme l'essentiel du cash-flow)~:
\end{itemize}
%
\begin{equation}
\label{eq:bloc1_dette}
f_{\text{dette}}(x) = 0{,}15 \cdot \frac{\max(0,\, x - 0{,}30)}{0{,}30} + 0{,}80 \cdot \frac{\max(0,\, x - 0{,}60)}{0{,}40}
\end{equation}
%
Le coefficient 0{,}80 du second palier (5{,}3 fois supérieur au premier) traduit l'accélération non linéaire du risque au-delà du seuil critique --- un phénomène documenté par la littérature sur l'instabilité financière \citep{merton1974}.

\begin{itemize}[nosep]
  \item \textbf{Couverture des intérêts (ICR)}~: le seuil de 1{,}5 correspond au point où l'EBITDA ne couvre plus confortablement les charges financières~; en dessous de 1{,}0, l'entreprise ne génère pas assez de trésorerie pour servir sa dette. Les guidelines EBA fixent typiquement le covenant ICR à 2{,}0.
  \item \textbf{Marge EBITDA}~: le seuil à 0 sépare les entreprises génératrices de cash (marge positive) des entreprises zombies (marge négative). Le coefficient fort ($0{,}60$) sur la marge négative reflète le fait qu'une firme qui brûle du cash a un horizon de survie limité.
\end{itemize}

Les effets protecteurs (taille, ancienneté) agissent en sens inverse. La taille est modélisée avec une saturation logarithmique~:
%
\begin{equation}
f_{\text{taille}}(x) = -0{,}10 \cdot \min\!\bigl(3{,}0,\; \ln(1 + x/10^7)\bigr)
\end{equation}
%
La saturation à 3{,}0 empêche les très grandes entreprises de bénéficier d'une protection illimitée --- un choix motivé par le fait que les grandes faillites (Lehman, Wirecard) montrent que la taille ne protège pas indéfiniment.

\subsection{Bloc 2~: interactions multiplicatives ($\sim$35\,\% du signal)}

Le deuxième bloc introduit des termes d'interaction entre variables.
La motivation économique est que les facteurs de risque ne sont pas additifs~: un endettement élevé est supportable si les flux de trésorerie sont stables, mais devient fatal si la volatilité est forte.
Ce phénomène correspond à la notion de \emph{zones de stabilité} de Minsky~:
%
\begin{equation}
\label{eq:bloc2}
g(\mathbf{x}) = \sum_{(i,j) \in \mathcal{I}} \alpha_{ij} \cdot \max(0, x_i - \tau_i) \cdot \mathbb{1}(x_j \in R_j) \cdot m_{\text{stress}}
\end{equation}
%
Le multiplicateur de stress $m_{\text{stress}} = 1{,}5$ en période de récession (T9--T11 du cycle) et 1{,}0 sinon capture la pro-cyclicité~: les mêmes caractéristiques de risque ont un pouvoir prédictif plus fort en récession, car les marges de sécurité sont érodées.
Le facteur 1{,}5 est calibré pour que le taux de défaut au creux du cycle soit environ 4 à 6 fois supérieur au taux tendanciel --- un ordre de grandeur cohérent avec les données EBA 2008--2012.

Les trois interactions les plus fortes sont~:
\begin{itemize}[nosep]
  \item $\text{dette} \times \text{volatilité}$ (coefficient $0{,}50$)~: la combinaison endettement élevé + flux instables est le premier prédicteur de défaut dans les LBO industriels \citep{hosmer_lemeshow2013}~;
  \item $\text{dette} \times \text{ICR}$ ($0{,}50$)~: un ratio de couverture faible transforme l'endettement en piège de liquidité~;
  \item $\text{dette} \times \text{marge}$ ($0{,}40$)~: la combinaison endettement + marge faible indique une incapacité structurelle à rembourser.
\end{itemize}

Ce bloc est volontairement conçu pour que la régression logistique, qui modélise $\log(p/(1-p)) = \mathbf{x}^\top \boldsymbol{\beta}$ sans termes croisés natifs, ne capture que $\sim$30\,\% de ce signal.
Les modèles à base d'arbres (XGBoost, TabNet), qui partitionnent l'espace des variables, en capturent 60--80\,\%.

\subsection{Bloc 3~: conjonctions ($\sim$25\,\% du signal)}

Le troisième bloc modélise les \emph{poches de risque} --- situations où le défaut ne survient que lorsque plusieurs conditions sont simultanément remplies~:
%
\begin{equation}
\label{eq:bloc3}
h(\mathbf{x}) = \sum_{S \in \mathcal{S}} \beta_S \cdot \prod_{i \in S} \mathbb{1}(x_i \in R_i) \cdot m_{\text{stress}}
\end{equation}
%
La conjonction la plus forte ($\beta = 0{,}80$) est $\text{ICR} < 2{,}5 \wedge \text{marge} < 8\,\%$~: c'est la \og spirale mortelle \fg{} où l'entreprise ne couvre pas ses intérêts \emph{et} ne génère pas assez de marge pour se redresser --- une situation qui, dans le modèle de \citet{merton1974}, correspond à une distance au défaut quasi-nulle.

Un terme quadratique sur la dette nette / EBITDA capture l'explosion du risque au-delà de 4{,}0$\times$~:
%
\begin{equation}
h_{\text{quad}}(x) = 0{,}30 \cdot \max(0,\, x_{\text{NDE}} - 4{,}0)^2 \cdot m_{\text{stress}}
\end{equation}
%
Le seuil de 4{,}0 correspond à la limite haute des covenants bancaires standard (les LBO structurés au-delà de 4{,}0$\times$ dette nette / EBITDA sont considérés comme \emph{highly leveraged} par l'EBA).

\subsection{Calibration de la PD}

Le score brut total $r = f(\mathbf{x}) + g(\mathbf{x}) + h(\mathbf{x})$ est transformé en probabilité de défaut via la fonction logistique~:
%
\begin{equation}
\label{eq:pd_calibration}
\text{PD}_i = \sigma\!\bigl(\alpha_0 + r_i\bigr) = \frac{1}{1 + e^{-(\alpha_0 + r_i)}}
\end{equation}
%
L'intercept $\alpha_0$ est calibré par la méthode de Brent (recherche de racine) pour que le taux de défaut moyen du portefeuille égale exactement la cible sectorielle $\text{PD}_{\text{base}}$ du tableau~\ref{tab:secteurs}~:
%
\begin{equation}
\alpha_0^* = \arg\min_{\alpha_0} \left| \frac{1}{n} \sum_{i=1}^n \sigma(\alpha_0 + r_i) - \text{PD}_{\text{cible}} \right|
\end{equation}

Le défaut est réalisé par tirage de Bernoulli~: $D_i \sim \text{Bernoulli}(\text{PD}_i)$.
Ce tirage stochastique introduit un bruit irréductible~: même un modèle parfait ne peut prédire le défaut avec certitude, puisque celui-ci est intrinsèquement probabiliste.

\subsection{Bruit hétéroscédastique et données manquantes}

Le DGP ajoute un bruit à queues lourdes dont la variance dépend du profil de risque de l'emprunteur~:
%
\begin{equation}
\epsilon_i = \sigma_i \cdot \frac{t_5}{\sqrt{3/5}}, \quad \sigma_i = \sigma_0 \bigl(1 + 0{,}3 \cdot I_{\text{PME}} + 0{,}2 \cdot \max(0, d_i - 0{,}3) + 0{,}2 \cdot \max(0, v_i - 0{,}2)\bigr)
\end{equation}
%
Le choix d'une Student à $\nu = 5$ degrés de liberté (plutôt qu'une gaussienne) produit des valeurs extrêmes plus fréquentes, ce qui est réaliste pour les scores de crédit.
L'hétéroscédasticité --- le bruit est 30\,\% plus fort pour les PME et augmente avec l'endettement et la volatilité --- reflète le fait que les petites entreprises endettées sont structurellement plus difficiles à discriminer (moins de données publiques, comptes moins audités).

Cinq variables de bruit pur (gaussien, uniforme, corrélées, permutées) sont ajoutées pour tester la capacité des modèles à ignorer les variables non informatives.

Les données manquantes suivent deux mécanismes distincts, motivés par la pratique bancaire~:
\begin{itemize}[nosep]
  \item \textbf{MNAR} (environ 10\,\%)~: les PME surendettées ne reportent pas leur marge EBITDA, les petites entreprises ne collectent pas de données sur les actifs tangibles --- un biais de sélection réaliste documenté par \citet{siddiqi2006}.
  \item \textbf{MCAR} (2\,\%)~: erreurs administratives (saisie, pannes système), indépendantes des valeurs~; le taux de 2\,\% est le seuil sous lequel l'estimation reste peu biaisée \citep{hosmer_lemeshow2013}.
\end{itemize}


% ════════════════════════════════════════════════════════
\section{Variables d'interaction pour l'apprentissage}
\label{sec:donnees:interactions}

Douze variables d'interaction sont pré-calculées à partir des variables brutes.
Leur raison d'être est de donner à la régression logistique accès aux effets croisés du Bloc~2, qu'elle ne peut pas découvrir seule.
XGBoost n'en a pas besoin~: il reconstruit nativement ces interactions via ses splits d'arbres.

Les interactions reproduisent exactement les termes du DGP (tableau~\ref{tab:interactions}), ce qui place la régression logistique dans les meilleures conditions possibles.
L'écart d'AUC résiduel entre LR et XGBoost mesure donc la part du signal qui provient des conjonctions d'ordre~3 et des non-linéarités que les interactions pré-calculées ne capturent pas.

\begin{table}[htbp]
\centering
\caption{Variables d'interaction (miroir du Bloc 2 du DGP).}
\label{tab:interactions}
\begin{tabular}{lp{7cm}}
\toprule
Variable & Définition et motivation \\
\midrule
\texttt{debt\_x\_vol}         & $\max(0, d - 0{,}30) \times I(v > 0{,}22)$ --- endettement $\times$ instabilité des CF \\
\texttt{debt\_x\_icr}         & $\max(0, d - 0{,}30) \times I(\text{ICR} < 2{,}5)$ --- endettement $\times$ couverture faible \\
\texttt{triple\_stress}       & $I(d > 0{,}45) \times I(v > 0{,}22) \times I(\text{ICR} < 2{,}5)$ --- conjonction triple de détresse \\
\texttt{nde\_quad}            & $\max(0, \text{NDE} - 4{,}0)^2$ --- explosion quadratique du levier \\
\bottomrule
\end{tabular}
\end{table}

Ces variables sont exclues de l'analyse de multicolinéarité (VIF) car elles sont corrélées par construction avec leurs composantes.


% ════════════════════════════════════════════════════════
\section{Données de private equity}
\label{sec:donnees:pe}

Pour chaque entreprise du portefeuille crédit, un fonds PE est associé.
Les paramètres de génération sont calibrés sur les indices Preqin 2023 (IRR médian par stratégie) et les données Cambridge Associates (MOIC par vintage)~:

\begin{table}[htbp]
\centering
\caption{Paramètres PE par secteur et sources de calibration.}
\label{tab:pe_params}
\begin{tabular}{lcccl}
\toprule
Secteur & Entrée (EV/EBITDA) & Levier & Sortie & Source \\
\midrule
Technologie & $[4;\; 8]$ & $[10;\; 40]\,\%$ & 9{,}0 & Preqin tech buyout~: IRR 12\,\%, MOIC $1{,}6\times$ \\
Industrie   & $[4;\; 8]$ & $[30;\; 60]\,\%$ & 8{,}5 & Preqin industrial~: IRR 11\,\%, MOIC $1{,}5\times$ \\
Santé       & $[10;\; 15]$ & $[20;\; 50]\,\%$ & 17{,}0 & Preqin healthcare~: IRR 10\,\%, MOIC $1{,}4\times$ \\
Immobilier  & $[4;\; 7]$ & $[50;\; 70]\,\%$ & 8{,}0 & Norme marché~: LTV 60--70\,\%, actifs collatéralisés \\
Services    & $[6;\; 10]$ & $[20;\; 50]\,\%$ & 11{,}0 & Preqin services~: IRR 10\,\%, MOIC $1{,}4\times$ \\
\bottomrule
\end{tabular}
\end{table}

La Santé affiche des multiples d'entrée très supérieurs (10--15$\times$ vs 4--8$\times$ pour l'Industrie)~: cela reflète la prime défensive d'un secteur à revenus régulés, où les acquéreurs acceptent de payer plus cher pour une moindre volatilité.
Le levier de l'Immobilier (50--70\,\%) est le plus élevé car les actifs physiques servent de collatéral --- une pratique standard dans les LBO immobiliers.
Inversement, la Technologie a le levier le plus faible (10--40\,\%) car la volatilité élevée des revenus limite la capacité d'endettement \citep{ipev2022}.

Le vintage est tiré uniformément entre 2018 et 2025, ce qui détermine la durée de détention ($h = 2026 - v$) et le calendrier de tirage du capital engagé~: 30\,\% en année~1, 60\,\% en année~2, 80\,\% en année~3, 100\,\% au-delà --- un profil en J standard conforme aux données ILPA.


% ════════════════════════════════════════════════════════
\section{Trajectoire macroéconomique}
\label{sec:donnees:macro}

Le générateur produit un historique de 24 trimestres (6 ans) simulant un cycle économique complet.
Les niveaux sont calibrés sur les données Eurostat et BCE 2010--2024~:

\begin{enumerate}[nosep]
  \item \textbf{Expansion} (T1--T4)~: PIB 1{,}5--2{,}2\,\% (croissance zone euro 2015--2018), chômage en baisse, taux en hausse~;
  \item \textbf{Décélération} (T5--T8)~: ralentissement progressif, type 2018--2019 pré-COVID~;
  \item \textbf{Stress} (T9--T11)~: PIB 0{,}8--1{,}3\,\%, chômage remonte --- le multiplicateur $m = 1{,}5$ s'active~;
  \item \textbf{Reprise} (T12--T16)~: amorce de reprise, taux bas~;
  \item \textbf{Normalisation} (T17--T24)~: retour vers les moyennes de long terme.
\end{enumerate}

Un bruit intra-trimestre est ajouté pour créer de l'hétérogénéité entre emprunteurs exposés au même trimestre~:
$\sigma_{\text{PIB}} = 1{,}80\,\text{pp}$ (dispersion régionale, Eurostat LFS), $\sigma_{\text{chômage}} = 1{,}00\,\text{pp}$ (Loire vs Île-de-France), $\sigma_{\text{taux}} = 1{,}70\,\text{pp}$ (spread de crédit variable selon l'emprunteur, Euribor + 150--300\,bp).

Cette trajectoire sert \emph{uniquement} à générer les données d'entraînement.
Les 12 scénarios de stress du chapitre~\ref{chap:metriques} sont des chocs appliqués \emph{a posteriori} au pipeline de provisionnement --- ils ne modifient pas le DGP.


% ════════════════════════════════════════════════════════
\section{Historique mensuel et réseau fournisseurs}
\label{sec:donnees:historique}

L'historique $\mathbf{D}_{\text{hist}}$ contient 12 mois de données par client.
Le retard de paiement mensuel (DPD) suit un processus dépendant du score de crédit~:
%
\begin{equation}
p_{\text{DPD},i} = \text{clip}\!\left(\frac{650 - \text{score}_i}{1\,000},\; 0,\; 0{,}3\right)
\end{equation}
%
Le seuil de 650 correspond au score FICO \emph{prime} aux États-Unis et, par extension, au seuil de bonne qualité dans les systèmes de scoring européens.
Les emprunteurs à score élevé ($> 650$) n'ont quasiment jamais de retard~; ceux à score faible ($< 450$) ont une probabilité de 20\,\% de retard chaque mois.

\subsection{Réseau d'approvisionnement}

Un réseau dirigé relie les entreprises via des relations fournisseur-client.
Les \emph{hubs} (5\,\% du portefeuille, sélectionnés par chiffre d'affaires) structurent le réseau~: cela reproduit la structure en étoile des chaînes d'approvisionnement industrielles, où quelques grands donneurs d'ordre concentrent les relations commerciales.

L'affinité inter-sectorielle gouverne la probabilité de connexion~: Technologie $\to$ Industrie ($2{,}0$) traduit la relation amont-aval naturelle (équipementiers industriels achetant des composants technologiques)~; les liens atypiques ont une affinité de $1{,}0$.
Ce réseau alimente le moteur de contagion du chapitre~\ref{chap:metriques}.


% ════════════════════════════════════════════════════════
\section{Décomposition du signal et performances attendues}
\label{sec:donnees:signal}

La répartition 40/35/25 entre les trois blocs est un choix de conception central.
Elle est calibrée pour produire un écart d'AUC d'environ 1 point entre la régression logistique et XGBoost --- un écart suffisant pour être statistiquement significatif sur 30\,000 observations, mais assez modeste pour que les conclusions sur le provisionnement IFRS~9 ne dépendent pas dramatiquement du choix de modèle.

\begin{table}[htbp]
\centering
\caption{Décomposition du signal et AUC attendu par famille de modèles.}
\label{tab:signal_decomposition}
\begin{tabular}{lcccc}
\toprule
Composante & Part du signal & LR & TabNet & XGBoost \\
\midrule
Monotone (paliers) & $\sim$40\,\% & $\sim$80\,\% & $\sim$85\,\% & $\sim$90\,\% \\
Interactions         & $\sim$35\,\% & $\sim$30\,\% & $\sim$60\,\% & $\sim$80\,\% \\
Conjonctions         & $\sim$25\,\% & $< 5\,\%$    & $\sim$40\,\% & $\sim$70\,\% \\
\midrule
\textbf{AUC attendu} & --- & \textbf{0{,}826} & \textbf{0{,}831} & \textbf{0{,}837} \\
\bottomrule
\end{tabular}
\end{table}

L'écart limité ($\sim$1{,}1 point d'AUC) s'explique par le bruit hétéroscédastique et les données manquantes, qui imposent un plafond de performance à tous les modèles.
Cette propriété est volontaire~: dans la réalité bancaire, l'AUC des modèles de scoring se situe entre 0{,}70 et 0{,}85 \citep{siddiqi2006}, et l'écart entre modèles linéaires et modèles de machine learning est rarement supérieur à 2 points \citep{grinsztajn2022}.


% ════════════════════════════════════════════════════════
\section{Synthèse}
\label{sec:donnees:synthese}

Le générateur produit un portefeuille de 30\,000 clients répartis sur cinq secteurs, avec une structure de dépendance réaliste (copule de Student), un DGP à trois blocs garantissant une hiérarchie de performances entre modèles, et des données annexes (PE, historique, bilan) pour tester l'ensemble du pipeline réglementaire.

Trois propriétés méritent d'être soulignées~:
\begin{enumerate}[nosep]
  \item \textbf{Reproductibilité}~: la graine $s = 123$ et son décalage par module garantissent des résultats identiques d'une exécution à l'autre.
  \item \textbf{Calibration}~: chaque paramètre est tracé vers une source institutionnelle (EBA, Moody's, Cambridge Associates, Eurostat, BCE) ou justifié par un raisonnement économique explicite.
  \item \textbf{Testabilité}~: la décomposition en trois blocs permet de diagnostiquer précisément quels types d'effets chaque modèle capture ou ignore.
\end{enumerate}

Le chapitre suivant exploite ces données pour entraîner et comparer trois familles de modèles de probabilité de défaut.


% ══════════════════════════════════════════════════════════
% CHAPITRE 2 — MODÈLES DE PROBABILITÉ DE DÉFAUT
% ══════════════════════════════════════════════════════════
\chapter{Modèles de probabilité de défaut}
\label{chap:modeles-pd}

L'estimation de la probabilité de défaut (PD) constitue la première brique du pipeline IFRS~9~: sans PD fiable, ni le staging, ni le calcul des pertes attendues ne peuvent prétendre à la robustesse. Ce chapitre présente les trois familles de modèles entraînés sur les données du chapitre précédent, en justifiant chaque choix architectural~--- du prétraitement des variables jusqu'à la calibration des probabilités.

\section{Pourquoi trois familles de modèles}
\label{sec:pd:trois-familles}

La littérature sur la modélisation tabulaire oppose traditionnellement les modèles linéaires, interprétables mais limités aux frontières convexes, aux modèles d'ensemble qui capturent les interactions sans ingénierie manuelle \citep{grinsztajn2022}. Le survey de \citet{borisov2022} montre qu'aucune famille ne domine universellement~: les réseaux profonds rivalisent avec le gradient boosting sur les jeux volumineux, mais souffrent sur les petits échantillons.

Ce constat motive le choix de trois familles complémentaires~:

\begin{enumerate}[nosep]
  \item \textbf{Régression logistique avec WoE} (LR-WoE)~: frontière de décision convexe, interprétabilité totale via un scorecard. C'est le standard réglementaire pour l'approche IRB \citep{anderson2007, siddiqi2006}.
  \item \textbf{TabNet} \citep{arik_pfister2021}~: réseau profond à attention séquentielle et Sparsemax. Il apprend des non-linéarités sans ingénierie de variables, tout en offrant des importances de variables natives (masques d'attention, non approximées comme SHAP).
  \item \textbf{XGBoost} \citep{chen_guestrin2016}~: ensemble d'arbres de décision par boosting de gradient. Frontières de décision en escalier, découverte automatique des interactions par les splits successifs.
\end{enumerate}

Ce triptyque n'est pas arbitraire. Il garantit que l'on dispose~:
\begin{itemize}[nosep]
  \item d'un modèle validable par un régulateur (LR-WoE)~;
  \item d'un modèle à haute capacité pour les données volumineuses (TabNet)~;
  \item d'un benchmark de performance maximale (XGBoost).
\end{itemize}

En production, la plupart des établissements utilisent XGBoost en arrière-plan et LR-WoE pour l'audit et la validation réglementaire \citep{anderson2007}.

\section{Prétraitement~: du brut au signal}
\label{sec:pd:pretraitement}

\subsection{Écrêtage des valeurs aberrantes}

Avant tout apprentissage, les variables sont écrêtées pour prévenir l'influence d'observations extrêmes sur les points de coupure du WoE et les splits d'arbre~:

\begin{table}[htbp]
\centering
\caption{Bornes d'écrêtage --- portefeuille corporate}
\label{tab:clipping-corporate}
\begin{tabular}{lccl}
\toprule
Variable & Min & Max & Justification \\
\midrule
\texttt{debt\_ratio} & 0{,}0 & 1{,}5 & Ratio d'endettement > 150\,\% = restructuration, hors scope \\
\texttt{credit\_score} & 300 & 850 & Plage FICO/Banque de France standard \\
\texttt{utilization\_rate} & 0{,}0 & 1{,}2 & Dépassement de 120\,\% = anomalie opérationnelle \\
\bottomrule
\end{tabular}
\end{table}

L'écrêtage est appliqué \emph{avant} l'ajustement du WoE-binner sur l'ensemble d'entraînement, de sorte que les points de coupure des quantiles ne soient pas déformés par des outliers.

\subsection{Ingénierie de variables~: les douze interactions}

Le chapitre précédent a montré que le DGP décompose le signal en trois blocs~: effets monotones (40\,\%), interactions multiplicatives (35\,\%) et conjonctions (25\,\%). La régression logistique, par construction linéaire dans l'espace WoE, ne peut capturer les blocs~2 et~3 qu'avec des variables d'interaction explicites.

Douze variables d'interaction sont calculées, chacune encodant un canal de risque conditionnel~:

\begin{table}[htbp]
\centering
\caption{Variables d'interaction (DGP blocs 2 et 3)}
\label{tab:pd-interactions}
\small
\begin{tabular}{lll}
\toprule
Variable & Formule & Canal économique \\
\midrule
\texttt{debt\_x\_hi\_vol} & $d \times \mathbf{1}(\sigma_{CF} > 0{,}22)$ & Levier amplifié par la volatilité \\
\texttt{debt\_x\_lo\_icr} & $d \times \mathbf{1}(\text{ICR} < 2{,}5)$ & Levier sous couverture d'intérêts faible \\
\texttt{debt\_x\_lo\_margin} & $d \times \mathbf{1}(m < 0{,}08)$ & Levier + marge EBITDA insuffisante \\
\texttt{margin\_neg\_x\_hi\_debt} & $\max(0, -m) \times \mathbf{1}(d > 0{,}45)$ & Marge négative sous endettement \\
\texttt{vol\_x\_hi\_debt} & $\sigma_{CF} \times \mathbf{1}(d > 0{,}45)$ & Volatilité sous levier \\
\texttt{icr\_low\_x\_hi\_debt} & $\max(0, 2 - \text{ICR}) \times \mathbf{1}(d > 0{,}45)$ & Couverture fragile sous levier \\
\texttt{incidents\_x\_hi\_debt} & $n_{\text{inc}} \times \mathbf{1}(d > 0{,}45)$ & Incidents de paiement sous levier \\
\texttt{util\_x\_hi\_vol} & $\max(0, u - 0{,}50) \times \mathbf{1}(\sigma_{CF} > 0{,}22)$ & Saturation du tirage + volatilité \\
\texttt{icr\_low\_AND\_margin\_low} & $\mathbf{1}(\text{ICR} < 2{,}5) \times \mathbf{1}(m < 0{,}08)$ & Conjonction ICR + marge (Minsky) \\
\texttt{debt\_hi\_AND\_vol\_hi} & $\mathbf{1}(d > 0{,}45) \times \mathbf{1}(\sigma_{CF} > 0{,}22)$ & Zone de fragilité financière \\
\texttt{util\_hi\_AND\_incidents} & $\mathbf{1}(u > 0{,}50) \times \mathbf{1}(n_{\text{inc}} \geq 2)$ & Utilisation + historique d'incidents \\
\texttt{nde\_squared\_excess} & $[\max(0, \text{NDE} - 4)]^2$ & Accélération du risque NDE > 4$\times$ \\
\bottomrule
\end{tabular}
\end{table}

Les seuils ne sont pas choisis arbitrairement. Ils correspondent aux conventions bancaires standard~:
\begin{itemize}[nosep]
  \item $d = 0{,}45$~: médiane du ratio d'endettement dans le portefeuille, seuil de vigilance des comités de crédit~;
  \item $\sigma_{CF} = 0{,}22$~: volatilité trimestrielle au-delà de laquelle le cash-flow ne couvre plus le service de la dette avec 95\,\% de confiance~;
  \item $\text{ICR} = 2{,}5$~: seuil covenant standard des prêts syndiqués (EBA GL 2017/06)~;
  \item $m = 0{,}08$~: marge EBITDA minimale pour absorber un choc de 200 pb sur le coût de la dette~;
  \item $u = 0{,}50$~: au-delà de 50\,\% d'utilisation, la flexibilité de financement se réduit significativement~;
  \item $\text{NDE} = 4$~: net debt/EBITDA au-delà duquel la capacité de remboursement devient non-linéairement fragile.
\end{itemize}

\paragraph{Pourquoi XGBoost n'utilise pas ces interactions.} Les arbres de décision apprennent les interactions nativement par splits successifs~: si $d > 0{,}45$ \emph{et} $\sigma_{CF} > 0{,}22$, l'arbre peut conditionner sur les deux variables à des nœuds imbriqués. Inclure les interactions manufacturées dans XGBoost diluerait l'importance des variables de base sans gain prédictif. Elles sont donc exclues de l'entraînement XGBoost et réservées aux modèles LR-WoE et TabNet.

\section{Weight of Evidence~: transformer pour interpréter}
\label{sec:pd:woe}

\subsection{Principe et formule}

La transformation Weight of Evidence (WoE) est un prétraitement standard en scoring crédit \citep{siddiqi2006, anderson2007}. Elle remplace chaque variable par son pouvoir discriminant bin par bin~:

\begin{equation}
\text{WoE}_b = \ln\!\left(\frac{n^{-}_b / N^{-}}{n^{+}_b / N^{+}}\right)
\label{eq:woe}
\end{equation}

\noindent où $n^{+}_b$ et $n^{-}_b$ désignent respectivement le nombre de défauts et de non-défauts dans le bin $b$, et $N^{+}$, $N^{-}$ les totaux du portefeuille. Un WoE positif indique un bin de meilleure qualité que la moyenne~; un WoE négatif, un bin plus risqué.

\paragraph{Lissage de Laplace.} Pour éviter la divergence $\ln(0)$ lorsqu'un bin ne contient aucun défaut (ou aucun non-défaut), on applique un lissage additif avec $\varepsilon = 0{,}5$~:
\begin{equation}
\text{WoE}_b = \ln\!\left(\frac{(n^{-}_b + \varepsilon) / (N^{-} + 2\varepsilon)}{(n^{+}_b + \varepsilon) / (N^{+} + 2\varepsilon)}\right)
\label{eq:woe-laplace}
\end{equation}
La valeur $\varepsilon = 0{,}5$ est le standard Bayésien (prior de Jeffreys), qui minimise le biais asymptotique.

\subsection{Information Value~: sélectionner les variables}

La valeur informationnelle (IV) mesure le pouvoir prédictif global d'une variable~:
\begin{equation}
\text{IV} = \sum_{b=1}^{B} \left(\frac{n^{-}_b}{N^{-}} - \frac{n^{+}_b}{N^{+}}\right) \times \text{WoE}_b
\label{eq:iv}
\end{equation}

\begin{table}[htbp]
\centering
\caption{Grille d'interprétation de l'Information Value}
\label{tab:iv-grille}
\begin{tabular}{cl}
\toprule
IV & Interprétation \\
\midrule
$< 0{,}02$ & Non prédictif~: variable exclue \\
$0{,}02 - 0{,}10$ & Pouvoir faible \\
$0{,}10 - 0{,}30$ & Pouvoir moyen \\
$0{,}30 - 0{,}50$ & Pouvoir fort \\
$> 0{,}50$ & Suspect~: risque de sur-apprentissage \\
\bottomrule
\end{tabular}
\end{table}

Les variables avec $\text{IV} < 0{,}02$ sont éliminées avant l'ajustement logistique. Ce seuil, issu de \citet{siddiqi2006}, évite d'introduire du bruit dans le modèle.

\subsection{Pipeline de discrétisation en six étapes}

La discrétisation WoE suit un pipeline en six étapes, conçu pour garantir à la fois la robustesse statistique et la cohérence économique~:

\begin{enumerate}
  \item \textbf{Quantiles}~: la variable est découpée en $B = 10$ bins de fréquence égale. Dix bins offrent un bon compromis entre résolution et stabilité~: trop peu lissent le signal, trop de bins fragmentent les observations.

  \item \textbf{Traitement des NaN}~: les valeurs manquantes forment un bin séparé avec leur propre WoE, plutôt que d'être imputées. Si le bin NaN représente moins de 5\,\% de la population, il est fusionné avec le bin numérique de WoE le plus proche.

  \item \textbf{Fusion des bins sous-peuplés}~: tout bin contenant moins de 20 défauts est fusionné avec son voisin le plus petit. Ce seuil assure qu'aucun bin ne repose sur une poignée d'observations (erreur standard du taux de défaut $\leq 1/\sqrt{20} \approx 22\,\%$ en relatif).

  \item \textbf{Détection de la direction}~: la corrélation de Spearman entre médiane du bin et taux de défaut détermine si le WoE doit être croissant ou décroissant. Cette détection automatique évite de spécifier manuellement le signe attendu de chaque variable.

  \item \textbf{Pool Adjacent Violators} (PAV)~: l'algorithme de régression isotonique fusionne les bins adjacents qui violent la monotonie. On obtient un WoE strictement monotone, cohérent avec le sens économique de la variable (par exemple, un endettement croissant ne peut pas devenir \emph{moins} risqué puis \emph{plus} risqué).

  \item \textbf{Taille minimale}~: les bins représentant moins de 5\,\% de la population sont fusionnés avec le voisin de WoE le plus proche, pour éviter les micro-segments instables.
\end{enumerate}

Le résultat est une grille de scoring stable, généralisable et économiquement interprétable~: chaque bin a un sens (zone de confort, zone grise, zone de danger), et la monotonie est garantie.

\section{Régression logistique~: le socle réglementaire}
\label{sec:pd:lr}

\subsection{Architecture et contraintes}

La régression logistique estime la log-cote de défaut comme combinaison linéaire des WoE~:
\begin{equation}
\ln\!\left(\frac{P(D=1|\mathbf{x})}{1 - P(D=1|\mathbf{x})}\right) = \beta_0 + \sum_{j=1}^{p} \beta_j \cdot \text{WoE}_j(\mathbf{x})
\label{eq:lr}
\end{equation}

\noindent soit, en inversant~:
\begin{equation}
\text{PD}(\mathbf{x}) = \frac{1}{1 + \exp\!\bigl(-\beta_0 - \sum_j \beta_j \cdot \text{WoE}_j(\mathbf{x})\bigr)}
\label{eq:pd-lr}
\end{equation}

L'ajustement utilise les paramètres suivants~:
\begin{itemize}[nosep]
  \item \textbf{Solveur}~: L-BFGS (méthode de Newton tronquée, convergence quadratique)~;
  \item \textbf{Régularisation}~: $C = 1{,}0$ (régularisation $L_2$ modérée, $\lambda = 1/C = 1$)~;
  \item \textbf{Pondération}~: \texttt{class\_weight="balanced"}, qui pondère chaque classe inversement à sa fréquence, ramenant la prévalence effective d'entraînement à $\tilde{\pi} = 0{,}5$.
\end{itemize}

\subsection{Filtrage VIF~: éliminer la multicolinéarité}

Avant l'ajustement, les variables sont filtrées par le \emph{Variance Inflation Factor} (VIF). Pour chaque variable $j$~:
\begin{equation}
\text{VIF}_j = \frac{1}{1 - R^2_j}
\label{eq:vif}
\end{equation}
où $R^2_j$ est le coefficient de détermination de la régression de la variable $j$ sur toutes les autres. Un VIF $> 5$ indique une multicolinéarité sévère \citep{marquardt1970}~: la variable est retirée itérativement (on supprime celle avec le VIF le plus élevé, on recalcule, et on répète).

\paragraph{Exemption des interactions.} Les douze variables d'interaction sont exclues du filtrage VIF. Elles sont corrélées \emph{par construction} avec les variables de base ($\texttt{debt\_x\_hi\_vol}$ est mécaniquement corrélée à $\texttt{debt\_ratio}$), mais encodent un signal conditionnel distinct (le risque du levier \emph{quand} la volatilité est élevée). Les soumettre au VIF les éliminerait systématiquement, détruisant la capacité du modèle linéaire à capturer les blocs~2 et~3 du DGP.

\subsection{Contrainte de signe~: cohérence économique}

Après le filtrage VIF, le modèle est ajusté et les signes des coefficients vérifiés. Pour les variables de base (non-interactions), on impose $\beta_j < 0$. La logique est la suivante~:
\begin{itemize}[nosep]
  \item Le WoE est positif pour les bins de bonne qualité ($\%$ non-défauts $>$ $\%$ défauts)~;
  \item Un coefficient $\beta_j < 0$ signifie qu'un WoE plus élevé \emph{diminue} la probabilité de défaut~;
  \item C'est économiquement cohérent~: un dossier de meilleure qualité a une PD plus basse.
\end{itemize}

Toute variable de base avec $\beta_j > 0$ est retirée itérativement~: on ajuste, on identifie la violation, on supprime, on réajuste. L'itération s'arrête quand tous les coefficients de base respectent la contrainte.

\paragraph{Exception pour les interactions.} Les variables d'interaction \emph{peuvent} avoir $\beta_j > 0$ sans violer la cohérence économique. Par exemple, pour \texttt{icr\_low\_AND\_margin\_low}, un WoE élevé signifie que les deux conditions sont réunies (ICR faible \emph{et} marge faible), ce qui est un signal de risque. Un $\beta > 0$ est alors correct~: la conjonction des deux faiblesses \emph{augmente} la probabilité de défaut.

\subsection{Correction du prior (King et Zeng)}

La pondération \texttt{balanced} crée un biais~: le modèle est ajusté comme si $\pi_{\text{train}} = 0{,}5$, alors que la prévalence réelle du portefeuille est $\pi_{\text{pop}} \approx 3$--5\,\%. \citet{king_zeng2001} proposent une correction de l'intercept~:

\begin{equation}
\hat{\beta}_0^{\text{corrigé}} = \hat{\beta}_0 + \ln\!\left(\frac{\pi_{\text{pop}}}{1 - \pi_{\text{pop}}}\right) - \ln\!\left(\frac{\pi_{\text{train}}}{1 - \pi_{\text{train}}}\right)
\label{eq:prior-correction}
\end{equation}

\noindent Avec $\pi_{\text{pop}} = 0{,}05$ et $\pi_{\text{train}} = 0{,}5$~:
$$\Delta\beta_0 = \ln\!\left(\frac{0{,}05}{0{,}95}\right) - \ln\!\left(\frac{0{,}5}{0{,}5}\right) = -2{,}944 - 0 = -2{,}944$$

L'intercept est décalé de $-2{,}944$, ramenant les probabilités prédites au niveau de la prévalence observée.

\subsection{Scorecard~: du modèle aux points}

La transformation en scorecard traduit la PD calibrée en un score entier, plus intuitif pour les analystes~:

\begin{equation}
\text{Score} = A + B \cdot \text{logit}(\text{PD}_{\text{cal}})
\label{eq:scorecard}
\end{equation}

\noindent avec~:
\begin{equation}
B = \frac{-\text{PDO}}{\ln 2}, \qquad A = S_0 - B \cdot \ln(O_0)
\label{eq:scorecard-params}
\end{equation}

\noindent où PDO (\emph{Points to Double the Odds}) $= 20$, $S_0 = 600$ (score à l'ancre), $O_0 = 50$ (cote à l'ancre, soit PD $\approx 1{,}96\,\%$). Ainsi~:
\begin{itemize}[nosep]
  \item PD $= 1{,}96\,\%$ $\Rightarrow$ Score $= 600$~;
  \item PD $= 0{,}98\,\%$ (cote doublée) $\Rightarrow$ Score $= 620$~;
  \item PD $= 3{,}85\,\%$ (cote divisée par deux) $\Rightarrow$ Score $= 580$.
\end{itemize}

Le PDO de 20 points est le standard industriel pour les scorecards bancaires \citep{siddiqi2006}. L'ancre à 600 / cote 50:1 est une convention qui place les PD typiques d'un portefeuille investment grade entre 500 et 700.

\section{TabNet~: l'attention au service du crédit}
\label{sec:pd:tabnet}

\subsection{Pourquoi un réseau profond sur données tabulaires}

Les données tabulaires sont le domaine naturel du gradient boosting. Pourtant, \citet{arik_pfister2021} montrent que TabNet rivalise avec XGBoost sur les grands jeux de données, avec deux avantages~:
\begin{enumerate}[nosep]
  \item Des \textbf{importances de variables natives}~: à chaque étape de décision, un masque d'attention Sparsemax sélectionne les variables pertinentes. L'importance agrégée est un sous-produit de l'inférence, non une approximation post-hoc (contrairement à SHAP).
  \item L'\textbf{apprentissage de représentations}~: les couches partagées entre étapes permettent de capturer des motifs non-linéaires sans ingénierie manuelle.
\end{enumerate}

\subsection{Architecture et hyperparamètres}

TabNet opère en $N_{\text{steps}}$ étapes séquentielles. À chaque étape $i$~:
\begin{enumerate}[nosep]
  \item Un masque d'attention $\mathbf{M}^{(i)} \in [0,1]^p$ est calculé par Sparsemax (variante parcimonieuse du softmax)~;
  \item Les variables sélectionnées $\mathbf{M}^{(i)} \odot \mathbf{h}$ alimentent un réseau de décision de dimension $N_d$~;
  \item Le coefficient de relaxation $\gamma$ contrôle la réutilisation des variables entre étapes ($\gamma = 1$ = pas de réutilisation, $\gamma > 1$ = réutilisation partielle).
\end{enumerate}

Deux variantes sont entraînées, adaptées à la taille des données~:

\begin{table}[htbp]
\centering
\caption{Hyperparamètres TabNet --- variantes Full et Light}
\label{tab:tabnet-config}
\begin{tabular}{lccp{5.5cm}}
\toprule
Paramètre & Full & Light & Justification \\
\midrule
$N_d$ (décision) & 64 & 16 & Full~: capacité pour 1{,}5M lignes. Light~: réduit le sur-apprentissage sur 50k. \\
$N_a$ (attention) & 64 & 16 & Symétrique avec $N_d$ (recommandation des auteurs) \\
$N_{\text{steps}}$ & 7 & 3 & 7 étapes $\times$ 64 dim $\approx$ 120k paramètres~; 3 étapes $\times$ 16 dim $\approx$ 8k \\
$\gamma$ (relaxation) & 1{,}3 & 1{,}5 & Light autorise davantage de réutilisation pour compenser la profondeur réduite \\
$\lambda_{\text{sparse}}$ & $10^{-3}$ & $10^{-2}$ & Light~: pénalité de parcimonie plus forte (moins de données) \\
Learning rate & 0{,}02 & 0{,}025 & Light~: convergence plus rapide pour compenser moins d'époques \\
Batch size & 8\,192 & 2\,048 & Proportionnel à la taille des données \\
Virtual batch (GBN) & 512 & 256 & Ghost Batch Normalization~: stabilise les gradients \\
Époques max & 200 & 100 & Light~: convergence plus rapide avec $N_d = 16$ \\
Patience & 20 & 15 & Arrêt précoce pour éviter le sur-apprentissage \\
\bottomrule
\end{tabular}
\end{table}

\paragraph{Pourquoi Full vs Light ?} La variante Full (120k paramètres) est dimensionnée pour le portefeuille corporate de 1{,}5M lignes~: le ratio données/paramètres est de $\sim$12, suffisant pour une généralisation robuste. La variante Light (8k paramètres, 15$\times$ plus petite) est utilisée pour les portefeuilles consommateur et hypothécaire (50k lignes chacun)~: avec 120k paramètres pour 50k observations, le réseau mémoriserait le jeu d'entraînement plutôt que d'apprendre des motifs généralisables.

\paragraph{Ghost Batch Normalization.} La normalisation par lots (\emph{batch normalization}) standard est instable avec les grands batchs ($B = 8\,192$) car les statistiques de chaque lot sont trop lisses. Le GBN découpe chaque lot en sous-lots virtuels de taille $B_v$ (512 ou 256), calcule les moyennes et variances sur ces sous-lots, puis les agrège. Cela introduit une régularisation stochastique sans réduire le débit de l'entraînement.

\subsection{Détection GPU}

L'entraînement détecte automatiquement la disponibilité CUDA~: sur GPU (RTX 5070 Ti dans notre configuration), un passage complet prend $\sim$30--60 secondes~; sur CPU, $\sim$5--10 minutes. XGBoost bénéficie également de l'accélération GPU via le backend \texttt{device="cuda"}.

\section{XGBoost~: la frontière de performance}
\label{sec:pd:xgb}

\subsection{Paramétrage}

XGBoost \citep{chen_guestrin2016} construit séquentiellement des arbres de décision, chacun corrigeant les erreurs résiduelles du précédent~:

\begin{table}[htbp]
\centering
\caption{Hyperparamètres XGBoost}
\label{tab:xgb-config}
\begin{tabular}{lcp{6cm}}
\toprule
Paramètre & Valeur & Justification \\
\midrule
\texttt{n\_estimators} & 400 & Suffisant pour convergence avec $\eta = 0{,}05$ (vérifié par courbe d'apprentissage) \\
\texttt{max\_depth} & 6 & Profondeur modérée~: capture les interactions 2-way et 3-way sans sur-apprentissage \\
\texttt{learning\_rate} $(\eta)$ & 0{,}05 & Shrinkage modéré~: permet un ajustement fin sans instabilité \\
\texttt{subsample} & 0{,}8 & Sous-échantillonnage ligne~: 80\,\% des observations par arbre (bagging) \\
\texttt{colsample\_bytree} & 0{,}8 & Sous-échantillonnage colonne~: 80\,\% des variables par arbre \\
\texttt{min\_child\_weight} & 3 & Poids minimum par feuille (empêche les feuilles avec 1-2 obs.) \\
\texttt{reg\_lambda} $(L_2)$ & 1{,}5 & Pénalité sur les poids des feuilles~: régularisation supplémentaire \\
\bottomrule
\end{tabular}
\end{table}

\paragraph{Pourquoi $\texttt{max\_depth} = 6$ ?} Chaque nœud d'un arbre binaire partitionne l'espace sur une variable. Avec une profondeur de 6, un arbre peut conditionner sur 6 variables successivement, capturant des interactions jusqu'à l'ordre~6. En pratique, les interactions d'ordre 2--3 dominent le DGP (blocs~2 et~3)~; une profondeur de 6 les capture confortablement tout en laissant une marge pour les effets de seuil.

\paragraph{Sous-échantillonnage 80/80.} Le double sous-échantillonnage (lignes et colonnes) réduit la corrélation entre arbres successifs, améliorant la généralisation. C'est l'analogue pour le boosting du bagging de Random Forest, mais appliqué à chaque arbre plutôt qu'à l'ensemble.

\paragraph{Gestion du déséquilibre.} Le paramètre \texttt{scale\_pos\_weight} est calculé automatiquement comme $n^{-}/n^{+}$ (ratio non-défauts/défauts), pondérant la fonction de perte pour traiter les deux classes équitablement.

\section{Calibration~: des scores aux probabilités}
\label{sec:pd:calibration}

\subsection{Pourquoi calibrer}

Un modèle discriminant (AUC élevée) n'est pas nécessairement bien calibré~: il peut classer correctement les défauts avant les non-défauts tout en produisant des probabilités systématiquement biaisées. Or, IFRS~9 exige des PD \emph{point-in-time} utilisées dans le calcul des ECL~:
$$\text{ECL} = \text{PD} \times \text{LGD} \times \text{EAD}$$
Si la PD est systématiquement surestimée de 50\,\%, l'ECL l'est aussi. La calibration corrige ce biais.

\subsection{Régression isotonique}

Les trois modèles sont calibrés par régression isotonique, une méthode non-paramétrique qui ajuste une fonction monotone croissante $f$ minimisant~:
\begin{equation}
\min_f \sum_{i=1}^{n} (y_i - f(s_i))^2 \quad \text{sous la contrainte} \quad s_i < s_j \Rightarrow f(s_i) \leq f(s_j)
\label{eq:isotonic}
\end{equation}
où $s_i$ est le score brut et $y_i \in \{0, 1\}$ le défaut observé.

\paragraph{Pourquoi isotonique plutôt que sigmoïde (Platt) ?} La calibration sigmoïde suppose que les log-cotes sont linéaires en le score brut~: $\ln(P/(1-P)) = a \cdot s + b$. Cette hypothèse est raisonnable pour la régression logistique (qui produit déjà des log-cotes), mais pas pour les arbres et les réseaux, dont les scores n'ont pas de forme fonctionnelle prédéfinie. La régression isotonique, libre de toute hypothèse paramétrique, s'adapte à n'importe quelle transformation monotone.

\paragraph{Validation croisée $k = 3$.} La calibration est ajustée par validation croisée à 3 plis~: à chaque pli, le modèle de base est entraîné sur 2/3 des données et l'isotonique est ajusté sur le tiers restant. Les trois courbes isotoniques sont moyennées. Le choix de $k = 3$ (plutôt que 5 ou 10) est motivé par la taille des données~: avec 1{,}5M observations, chaque pli contient 500k observations, largement suffisant pour ajuster une courbe isotonique stable.

\section{Métriques d'évaluation}
\label{sec:pd:metriques}

\subsection{Discrimination~: AUC et Gini}

L'aire sous la courbe ROC (AUC-ROC) mesure la probabilité qu'un défaut tiré au hasard ait un score plus élevé qu'un non-défaut~:
\begin{equation}
\text{AUC} = \frac{1}{n^+ \cdot n^-} \sum_{i \in \mathcal{D}^+} \sum_{j \in \mathcal{D}^-} \mathbf{1}(s_i > s_j)
\label{eq:auc}
\end{equation}

Le coefficient de Gini normalise l'AUC~:
\begin{equation}
\text{Gini} = 2 \times \text{AUC} - 1
\label{eq:gini}
\end{equation}

Un modèle aléatoire a Gini $= 0$~; un modèle parfait, Gini $= 1$.

\subsection{Séparation~: statistique de Kolmogorov-Smirnov}

La statistique KS mesure la distance maximale entre les fonctions de répartition empiriques des défauts et des non-défauts~:
\begin{equation}
\text{KS} = \max_s \left| F^+(s) - F^-(s) \right|
\label{eq:ks}
\end{equation}

Un KS $> 0{,}40$ indique une séparation excellente~; entre $0{,}20$ et $0{,}40$, acceptable. Nos trois modèles dépassent 0{,}40.

\subsection{Calibration~: score de Brier et log-loss}

Le score de Brier mesure l'erreur quadratique moyenne entre la PD prédite et le défaut observé~:
\begin{equation}
\text{Brier} = \frac{1}{N} \sum_{i=1}^{N} \bigl(\hat{p}_i - y_i\bigr)^2
\label{eq:brier}
\end{equation}

La log-loss (entropie croisée binaire) pénalise plus sévèrement les prédictions confiantes mais erronées~:
\begin{equation}
\text{LogLoss} = -\frac{1}{N} \sum_{i=1}^{N} \bigl[y_i \ln \hat{p}_i + (1 - y_i) \ln(1 - \hat{p}_i)\bigr]
\label{eq:logloss}
\end{equation}

\subsection{Stabilité~: Population Stability Index}

Le PSI mesure la dérive entre la distribution des scores d'entraînement et celle d'un nouveau jeu~:
\begin{equation}
\text{PSI} = \sum_{b=1}^{B} \left(\%_{\text{actuel},b} - \%_{\text{référence},b}\right) \times \ln\!\left(\frac{\%_{\text{actuel},b}}{\%_{\text{référence},b}}\right)
\label{eq:psi}
\end{equation}

Un PSI $< 0{,}10$ indique une distribution stable~; entre $0{,}10$ et $0{,}25$, une dérive modérée~; au-delà, un recalibrage est nécessaire.

\section{Résultats comparés}
\label{sec:pd:resultats}

\subsection{Portefeuille corporate (1{,}5M observations)}

\begin{table}[htbp]
\centering
\caption{Performance des trois modèles --- portefeuille corporate}
\label{tab:resultats-corporate}
\begin{tabular}{lccc}
\toprule
Métrique & LR-WoE & TabNet & XGBoost \\
\midrule
AUC (test) & 0{,}826 & 0{,}831 & 0{,}837 \\
Gini & 0{,}652 & 0{,}662 & 0{,}674 \\
KS & $> 0{,}42$ & $> 0{,}43$ & $> 0{,}44$ \\
Écart train--test & 0{,}000 & 0{,}000 & 0{,}000 \\
\midrule
Interactions utilisées & Oui (12) & Oui (12) & Non (natives) \\
Nombre de paramètres & $\sim$50 & $\sim$120k & $\sim$400 arbres \\
Temps d'entraînement & $\sim$5\,s & $\sim$45\,s & $\sim$15\,s \\
\bottomrule
\end{tabular}
\end{table}

\paragraph{Lecture des résultats.} L'écart AUC entre LR-WoE (0{,}826) et XGBoost (0{,}837) est de 1{,}1 point de pourcentage. Cet écart est \emph{exact}ement celui visé par le DGP~: les blocs~2 et~3 (interactions + conjonctions) représentent $\sim$60\,\% du signal, mais les 12 variables d'interaction manufacturées en capturent environ les deux tiers. Le tiers restant (conjonctions 3-way, non-stationnarité stress $\times 1{,}5$) ne peut être capté que par des modèles non-linéaires.

\paragraph{Zéro écart train--test.} L'absence de sur-apprentissage ($\Delta\text{AUC} = 0{,}000$) s'explique par la combinaison de trois facteurs~: (i) la calibration isotonique avec $k=3$ régularise naturellement, (ii) le jeu d'entraînement de 1{,}5M observations est suffisamment grand pour que les modèles convergent vers la distribution réelle, et (iii) les hyperparamètres sont conservateurs (profondeur 6 pour XGBoost, patience 20 pour TabNet).

\paragraph{Hiérarchie structurelle.} La hiérarchie LR $<$ TabNet $<$ XGBoost est \emph{garantie par construction} du DGP~: les effets monotones (bloc~1) sont parfaitement capturés par les trois modèles~; les interactions (bloc~2) sont partiellement capturées par LR via les 12 variables manufacturées et nativement par TabNet/XGBoost~; les conjonctions 3-way et la non-stationnarité (bloc~3) ne sont accessibles qu'aux modèles non-linéaires. L'écart $\sim$1\,pp est calibré par l'algorithme de Brent (cf. chapitre~1) pour être réaliste d'un portefeuille corporate.

\subsection{Portefeuilles consommateur et hypothécaire}

\begin{table}[htbp]
\centering
\caption{Performance --- portefeuilles consommateur et hypothécaire}
\label{tab:resultats-consumer-mortgage}
\begin{tabular}{lccl}
\toprule
Portefeuille & AUC XGBoost & AUC LR-WoE & Source des données \\
\midrule
Consommateur & 0{,}798 & 0{,}791 & Lending Club (50k réels) \\
Hypothécaire & 0{,}743 & 0{,}758 & Synthétique (50k) \\
\bottomrule
\end{tabular}
\end{table}

\paragraph{Consommateur.} Les données proviennent de Lending Club (50k prêts réels). L'AUC est plus basse que le corporate (0{,}798 vs 0{,}837) pour deux raisons~: (i) l'échantillon est 30$\times$ plus petit, limitant la capacité des modèles à apprendre des motifs fins, et (ii) les prêts à la consommation sont moins structurés que le crédit corporate (moins de ratios financiers, plus de bruit comportemental).

\paragraph{Hypothécaire.} L'inversion LR $>$ XGBoost (0{,}758 vs 0{,}743) sur le portefeuille hypothécaire illustre un phénomène classique~: sur les petits jeux synthétiques avec peu de variables, la régression logistique avec WoE est \emph{optimale par construction} (le DGP hypothécaire est plus simple, avec des effets principalement monotones). XGBoost sur-découpe les splits sans gain discriminant.

\paragraph{TabNet Light.} Les deux portefeuilles utilisent la variante Light (8k paramètres) de TabNet. Le ratio données/paramètres de $\sim$6 reste acceptable~: la patience de 15 époques et la pénalité de parcimonie $\lambda = 10^{-2}$ préviennent le sur-apprentissage.

\subsection{Prévention des fuites de données}

La rigueur de l'évaluation repose sur une séparation stricte entre entraînement et test~:

\begin{enumerate}[nosep]
  \item \textbf{Variables cibles masquées}~: \texttt{default\_flag} (cible), \texttt{pd\_latent} et \texttt{pd\_origination} (vérités terrain du DGP) sont retirées des prédicteurs. Leur inclusion donnerait des AUC $> 0{,}99$, sans signification prédictive.
  \item \textbf{WoE ajusté sur l'entraînement uniquement}~: les points de coupure, les valeurs de WoE et les IV sont calculés sur le jeu d'entraînement. Le jeu de test utilise les bins pré-calculés, sans recalcul.
  \item \textbf{Graines séparées}~: l'entraînement utilise $s = 42$ (indépendant de la génération du portefeuille, $s = 123$). Même si les données sont synthétiques, les deux processus aléatoires sont décorrélés.
  \item \textbf{Holdout consommateur}~: les données Lending Club sont séparées en 50k (entraînement) et 15k (test holdout). Le pipeline de production utilise exclusivement le holdout, garantissant qu'aucune observation d'entraînement ne fuit dans les résultats du tableau de bord.
\end{enumerate}

\section{Fonctions de transformation~: logit et expit}
\label{sec:pd:logit-expit}

Deux fonctions sont omniprésentes dans le pipeline~: le logit (log-cote) et son inverse, l'expit (sigmoïde). Elles transforment les probabilités en scores non bornés, et inversement~:

\begin{equation}
\text{logit}(p) = \ln\!\left(\frac{p}{1-p}\right), \qquad \text{expit}(x) = \frac{1}{1 + e^{-x}}
\label{eq:logit-expit}
\end{equation}

Le stress testing opère dans l'espace logit~: un choc macroéconomique est ajouté au logit de la PD de base, puis retransformé en probabilité par l'expit. Cette approche garantit que la PD stressée reste dans $[0, 1]$, même pour des chocs importants.

\paragraph{Stabilité numérique.} L'implémentation de l'expit utilise deux branches pour éviter les débordements~:
$$
\text{expit}(x) = \begin{cases}
\displaystyle\frac{1}{1 + e^{-x}} & \text{si } x \geq 0 \\[8pt]
\displaystyle\frac{e^x}{1 + e^x} & \text{si } x < 0
\end{cases}
$$
Pour $x \gg 0$, $e^{-x} \to 0$ et la première branche est stable~; pour $x \ll 0$, $e^x \to 0$ et la seconde branche évite le calcul de $e^{-x}$ qui déborderait.

\section{Synthèse et choix du modèle actif}
\label{sec:pd:synthese}

Trois constats émergent de la comparaison~:

\begin{enumerate}
  \item \textbf{La hiérarchie est structurelle}~: XGBoost $>$ TabNet $>$ LR-WoE sur le corporate, avec un écart de $\sim$1\,pp entre LR et XGBoost. Cette hiérarchie reflète la structure du DGP (effets monotones vs interactions vs conjonctions), et non un ajustement des hyperparamètres.

  \item \textbf{L'écart est modeste}~: 1{,}1\,pp d'AUC entre LR et XGBoost. Sur un portefeuille de 30\,000 clients, cela représente environ 330 clients mieux classés par XGBoost que par LR ($1{,}1\,\% \times 30\,000$). En termes d'ECL, l'impact dépend de la concentration du risque dans ces 330 clients.

  \item \textbf{Le choix est contextuel}~: en production, le modèle actif est sélectionnable. Pour un reporting réglementaire COREP, LR-WoE offre la transparence maximale. Pour un stress test interne, XGBoost exploite mieux les non-linéarités. Le pipeline permet de basculer sans re-entraîner.
\end{enumerate}

Le chapitre suivant utilise ces PD calibrées pour construire le moteur de pertes attendues (ECL) conforme à IFRS~9, avec le staging SICR et les ajustements LGD \emph{downturn}.


% ══════════════════════════════════════════════════════════
% CHAPITRE 3 — MOTEUR ECL IFRS 9
% ══════════════════════════════════════════════════════════
\chapter{Moteur de pertes attendues (ECL) sous IFRS~9}
\label{chap:ecl}

La probabilité de défaut, estimée au chapitre précédent, n'est que la première composante de la perte attendue. IFRS~9 impose un calcul de perte qui intègre trois dimensions supplémentaires~: le \emph{staging} (horizon de la perte), la perte en cas de défaut (LGD) et l'exposition au moment du défaut (EAD). Ce chapitre présente le moteur complet, du classement en stages jusqu'à l'agrégation multi-scénarios.

\section{Le staging IFRS~9~: classifier pour provisionner}
\label{sec:ecl:staging}

\subsection{Trois stages, trois horizons}

IFRS~9 classe chaque exposition dans l'un de trois stages, qui déterminent l'horizon de la perte attendue~:

\begin{table}[htbp]
\centering
\caption{Classification IFRS~9 en trois stages}
\label{tab:stages}
\begin{tabular}{clcl}
\toprule
Stage & Condition & Horizon ECL & Interprétation \\
\midrule
1 & Qualité de crédit stable & 12 mois & Provision sur la perte à un an \\
2 & Dégradation significative (SICR) & Durée de vie & Provision sur toute la durée \\
3 & Défaut observé & Durée de vie & Défaut avéré, PD $= 100\,\%$ \\
\bottomrule
\end{tabular}
\end{table}

L'assignation suit un ordre de priorité strict~:
\begin{enumerate}[nosep]
  \item \textbf{Stage~3}~: si \texttt{default\_flag} $= 1$ \emph{ou} DPD $\geq 90$ jours.
  \item \textbf{Stage~2}~: si le score SICR dépasse le seuil $\tau = 1{,}65$.
  \item \textbf{Stage~1}~: par défaut (pas de dégradation significative).
\end{enumerate}

\paragraph{Pourquoi le seuil DPD de 90 jours.} Ce seuil est la définition réglementaire du défaut selon Bâle~III et le CRR (art.~178)~: une contrepartie est en défaut lorsqu'elle est en retard de paiement de plus de 90 jours sur une obligation significative. C'est un critère objectif, observable, et harmonisé au niveau européen.

\paragraph{Suppression du seuil PD pour le Stage~3.} Une version antérieure du moteur classait en Stage~3 les expositions dont la PD courante dépassait 30\,\%. Ce critère a été retiré conformément à IFRS~9~B5.5.21, qui exige des \emph{événements observés} (défaut avéré, retard de paiement), et non des prédictions de modèle. Un modèle prédisant une PD de 35\,\% n'est pas équivalent à un défaut constaté~: le maintenir conduirait à des reclassements procycliques en Stage~3 lors des stress tests, sans fait générateur.

\subsection{Le score SICR~: quatre composantes}

Le \emph{Significant Increase in Credit Risk} (SICR) est le mécanisme central du staging IFRS~9. La norme (§B5.5.17) exige une évaluation multi-factorielle, combinant des indicateurs quantitatifs et qualitatifs. Notre implémentation agrège quatre composantes en un score unique~:

\begin{equation}
S_{\text{SICR}} = w_1 \cdot \underbrace{\left(\frac{\text{PD}_t}{\text{PD}_0} - 1\right)}_{\text{ratio PD}} + w_2 \cdot \underbrace{\max(0,\, \text{PD}_t - \text{PD}_0)}_{\text{delta PD}} + w_3 \cdot \underbrace{\frac{\text{DPD}}{30}}_{\text{retard}} + w_4 \cdot \underbrace{Z_{\text{macro}}}_{\text{forward-looking}}
\label{eq:sicr}
\end{equation}

\noindent avec les poids~:

\begin{table}[htbp]
\centering
\caption{Pondérations du score SICR}
\label{tab:sicr-poids}
\begin{tabular}{clcp{5.5cm}}
\toprule
Composante & Poids & Valeur & Justification \\
\midrule
Ratio PD ($w_1$) & 0{,}25 & $\text{PD}_t / \text{PD}_0 - 1$ & Détérioration \emph{relative}~: une PD passant de 0{,}1\,\% à 0{,}5\,\% ($\times 5$) est plus significative qu'une PD passant de 5\,\% à 5{,}4\,\% ($\times 1{,}08$). \\
Delta PD ($w_2$) & 0{,}25 & $\max(0, \text{PD}_t - \text{PD}_0)$ & Détérioration \emph{absolue}~: capte les cas où le ratio est modéré mais le niveau absolu est préoccupant. Unilatéral (seule la hausse compte). \\
Retard ($w_3$) & 0{,}20 & DPD / 30 & Signal comportemental observable, non modélisé. 30 jours $= 1{,}0$~: un mois de retard contribue significativement au score. \\
Macro ($w_4$) & 0{,}30 & $Z_{\text{macro}}$ & Information \emph{forward-looking} exigée par IFRS~9~B5.5.17(a). Poids le plus élevé~: les conditions macro affectent l'ensemble du portefeuille simultanément. \\
\bottomrule
\end{tabular}
\end{table}

\paragraph{Pourquoi 50\,\% pour la PD.} Les deux composantes PD (ratio et delta) totalisent 50\,\% du score. C'est le signal le plus prédictif de la dégradation de crédit~: il intègre à la fois les fondamentaux de la contrepartie (ratios financiers, comportement de paiement) et les conditions de marché (via le stress de la PD). Le ratio capte les dégradations relatives (importantes pour les contreparties de haute qualité)~; le delta capte les dégradations absolues (importantes pour les contreparties déjà fragiles).

\paragraph{Pourquoi 30\,\% pour le macro.} Le poids macro est le plus élevé individuellement. IFRS~9 insiste sur l'intégration d'informations \emph{forward-looking}~: un scénario macroéconomique adverse doit déclencher des migrations en Stage~2 \emph{avant} que les retards de paiement ne se matérialisent. Ce poids assure que les provisions augmentent dès la dégradation du contexte macro, sans attendre les premiers défauts.

\paragraph{Seuil $\tau = 1{,}65$.} Le seuil est calibré empiriquement pour produire un taux de migration en Stage~2 de $\sim$10\,\% en scénario de base. Ce taux est cohérent avec les publications EBA sur les banques européennes (8--15\,\% de l'encours en Stage~2).

\paragraph{Écrêtage du ratio PD.} Le ratio PD est écrêté à 5{,}0 pour éviter une sensibilité excessive lorsque la PD d'origination est très basse. Sans écrêtage, une PD passant de $0{,}01\,\%$ à $0{,}10\,\%$ donnerait un ratio de $9{,}0$, déclenchant le SICR pour une exposition de qualité investment grade avec un risque absolu négligeable.

\subsection{Le Z-score macroéconomique}

La composante macro du SICR traduit les conditions macroéconomiques en un score standardisé. Cinq variables sont normalisées par leur volatilité historique (2010--2024) et combinées en tenant compte de la direction économique~:

\begin{equation}
Z_{\text{macro}} = \frac{1}{5} \sum_{i=1}^{5} \frac{\epsilon_i \cdot (x_i - \mu_i)}{\sigma_i}
\label{eq:z-macro}
\end{equation}

\noindent où $\epsilon_i = +1$ pour les variables dont la hausse est adverse (chômage, taux d'intérêt), et $\epsilon_i = -1$ pour celles dont la hausse est favorable (PIB, HPI). Les constantes de normalisation~:

\begin{table}[htbp]
\centering
\caption{Normalisation du Z-score macroéconomique}
\label{tab:z-norm}
\begin{tabular}{lcccl}
\toprule
Variable & $\mu$ (base) & $\sigma$ & $\epsilon$ & Justification \\
\midrule
Chômage & 7{,}5\,\% & 3{,}0\,pp & $+1$ & Hausse $\Rightarrow$ adverse (Okun) \\
PIB & 1{,}2\,\% & 2{,}7\,pp & $-1$ & Baisse $\Rightarrow$ adverse \\
Taux d'intérêt & 3{,}5\,\% & 1{,}5\,pp & $+1$ & Hausse $\Rightarrow$ adverse (charge de la dette) \\
HPI & 2{,}0\,\% & 10{,}0\,pp & $-1$ & Baisse $\Rightarrow$ adverse (collatéral) \\
Inflation & 2{,}5\,\% & 3{,}0\,pp & $+1$ & Hausse $\Rightarrow$ adverse (érosion des marges) \\
\bottomrule
\end{tabular}
\end{table}

La normalisation par $\sigma$ empêche les variables à large amplitude (HPI~: $\sim$30\,pp de plage) de dominer celles à faible amplitude (taux~: $\sim$2\,pp). Chaque variable contribue en nombre de déviations standard par rapport au scénario de base.

\subsection{Exemption souveraine (IFRS~9 B5.5.25)}

Les expositions souveraines sont exemptées du staging par la norme IFRS~9 (paragraphe B5.5.25). Elles restent en Stage~1 en permanence et ne supportent qu'une ECL à 12 mois. La justification normative est que le risque souverain est de nature systémique~: une dégradation du risque souverain affecte l'ensemble de l'économie, et les provisions individuelles ne constituent pas l'instrument approprié pour gérer ce risque.

Dans le moteur, cette exemption est implémentée par le drapeau \texttt{exempt\_from\_staging = True} dans le profil de la classe d'actifs souveraine.

\section{La perte en cas de défaut (LGD)}
\label{sec:ecl:lgd}

\subsection{LGD de base~: trois ajustements}

La LGD de base est calculée en trois étapes successives~:

\begin{equation}
\text{LGD}_{\text{base}} = \text{LGD}_{\text{secteur}} \times m_{\text{type}} + \Delta_{\text{score}}
\label{eq:lgd-base}
\end{equation}

\paragraph{1. LGD sectorielle.} Chaque secteur a une LGD moyenne calibrée sur les taux de recouvrement observés. L'immobilier a la LGD la plus élevée ($\sim$38\,\%) en raison des coûts de liquidation des actifs immobiliers~; la santé la plus basse ($\sim$33\,\%) grâce à des flux de trésorerie plus stables.

\paragraph{2. Multiplicateur par type de prêt.} Les crédits revolving reçoivent un multiplicateur $m = 1{,}15$ (LGD majorée de 15\,\%) car ils sont structurellement subordonnés aux prêts à terme dans les conventions de priorité. Les prêts à terme bénéficient d'un multiplicateur $m = 0{,}90$ (LGD réduite de 10\,\%).

\paragraph{3. Ajustement par score de crédit.} Un ajustement continu, borné à $\pm 2\,\%$, reflète la qualité intrinsèque de la contrepartie~:
\begin{equation}
\Delta_{\text{score}} = \text{clip}\!\left(\frac{650 - \text{score}}{200}, -0{,}10, 0{,}10\right) \times 0{,}20
\label{eq:lgd-score}
\end{equation}
Un score de 650 (médiane) produit un ajustement nul~; un score de 750 réduit la LGD de ${\sim}\,2\,\%$~; un score de 550 l'augmente de ${\sim}\,2\,\%$. L'amplitude est volontairement modeste~: la LGD est principalement déterminée par les caractéristiques de l'exposition (séniorité, collatéral), pas par la qualité de l'emprunteur.

\subsection{Dispersion par distribution Bêta}

Pour capturer l'hétérogénéité des taux de recouvrement, la LGD de base est dispersée par une distribution Bêta, dont les paramètres sont calculés par la méthode des moments~:

\begin{equation}
\alpha = \mu \cdot \kappa, \qquad \beta = (1 - \mu) \cdot \kappa, \qquad \kappa = \frac{\mu(1-\mu)}{\sigma^2} - 1
\label{eq:beta-params}
\end{equation}

\noindent avec $\mu = \text{LGD}_{\text{base}}$ et $\sigma = 0{,}15$ (écart-type TTC). La distribution Bêta est le choix naturel pour les taux de recouvrement~: elle est définie sur $[0, 1]$, peut être asymétrique (recouvrement concentré autour de 40\,\% mais avec une queue vers les pertes totales), et ne nécessite que deux paramètres.

La LGD est bornée entre 5\,\% (plancher de recouvrement~: même un défaut total laisse une valeur résiduelle) et 95\,\% (plafond~: une perte totale est théoriquement possible mais extrêmement rare pour des expositions collatéralisées).

\subsection{LGD \emph{downturn}~: le cycle amplifie la perte}

IFRS~9 et l'EBA GL/2019/03 exigent que la LGD reflète les conditions de \emph{downturn}~: en récession, les taux de recouvrement chutent car les actifs sont vendus dans un marché déprimé. L'ajustement \emph{downturn} est proportionnel à l'intensité du stress macroéconomique~:

\begin{equation}
\text{LGD}_{\text{DT}} = \text{LGD}_{\text{TTC}} \times \bigl(1 + \rho_{\text{LGD}} \cdot |Z_{\text{stress}}|\bigr)
\label{eq:lgd-downturn}
\end{equation}

\noindent où $\rho_{\text{LGD}}$ est la corrélation sectorielle entre la LGD et le cycle économique, et $Z_{\text{stress}}$ mesure l'intensité du stress en nombre de déviations standard ($Z = 0$ pour le scénario de base, $Z = 2$ pour un stress adverse).

\begin{table}[htbp]
\centering
\caption{Corrélation LGD-cycle par secteur ($\rho_{\text{LGD}}$)}
\label{tab:rho-lgd}
\begin{tabular}{lcp{6cm}}
\toprule
Secteur & $\rho_{\text{LGD}}$ & Justification \\
\midrule
Technologie & 0{,}25 & Actifs incorporels (PI, logiciels)~: valeur de liquidation volatile mais modérée \\
Industrie & 0{,}30 & Actifs tangibles mais spécifiques~: machines industrielles difficiles à revendre en récession \\
Santé & 0{,}10 & Demande inélastique, flux défensifs~: la récession affecte peu les recouvrements \\
Immobilier & 0{,}35 & Actifs immobiliers fortement procycliques~: corrélation maximale avec le marché \\
Services & 0{,}20 & Actifs légers (capital humain)~: la LGD dépend du goodwill client, modérément cyclique \\
\bottomrule
\end{tabular}
\end{table}

\paragraph{Exemple.} Pour l'immobilier ($\rho = 0{,}35$) en scénario adverse ($Z = 2$)~:
$$\text{LGD}_{\text{DT}} = 0{,}38 \times (1 + 0{,}35 \times 2) = 0{,}38 \times 1{,}70 = 0{,}646$$
La LGD passe de 38\,\% à 64{,}6\,\%~: en récession immobilière, les recouvrements chutent car les actifs sont vendus à des prix déprimés, les procédures s'allongent, et les frais juridiques augmentent.

\paragraph{Invariant.} La formule garantit $\text{LGD}_{\text{DT}} \geq \text{LGD}_{\text{TTC}}$ pour $Z \geq 0$~: la LGD \emph{downturn} ne peut jamais être inférieure à la LGD \emph{through-the-cycle}. Cet invariant est vérifié par les tests unitaires.

\subsection{Canal HPI~: le collatéral immobilier}

Pour les secteurs sensibles au prix de l'immobilier, un canal supplémentaire ajuste la LGD en fonction de la variation du HPI~:

\begin{equation}
\Delta_{\text{HPI}} = \frac{2{,}0 - \text{HPI}_{\text{scénario}}}{100} \times h_s
\label{eq:lgd-hpi}
\end{equation}

\noindent où $h_s$ est la sensibilité HPI du secteur $s$ (2{,}5 pour l'immobilier, 0{,}3 pour la santé), et 2{,}0\,\% est la croissance HPI de base. Le mécanisme est bidirectionnel~: une baisse du HPI augmente la LGD (collatéral déprécié), une hausse la réduit.

\paragraph{Exemple (crise immobilière).} HPI $= -3\,\%$, secteur Immobilier ($h = 2{,}5$)~:
$$\Delta_{\text{HPI}} = \frac{2{,}0 - (-3{,}0)}{100} \times 2{,}5 = 0{,}05 \times 2{,}5 = 0{,}125$$
La LGD augmente de 12{,}5 points~: un effondrement des prix immobiliers se traduit directement en pertes accrues sur les expositions collatéralisées.

\section{L'exposition au défaut (EAD)}
\label{sec:ecl:ead}

L'EAD mesure l'exposition anticipée au moment du défaut. Elle diffère de l'encours actuel pour les facilités renouvelables, où l'emprunteur peut tirer des montants supplémentaires avant de faire défaut.

\subsection{Formule}

\begin{equation}
\text{EAD} = \begin{cases}
\text{montant du prêt} & \text{si prêt à terme (CCF} = 1{,}0\text{)} \\
\text{tiré} + \text{CCF} \times \text{non tiré} & \text{si revolving}
\end{cases}
\label{eq:ead}
\end{equation}

\noindent avec~:
\begin{itemize}[nosep]
  \item $\text{tiré} = \text{montant} \times \text{taux d'utilisation}$~;
  \item $\text{non tiré} = \text{montant} \times (1 - \text{taux d'utilisation})$~;
  \item $\text{CCF} = 0{,}75$ pour les revolving (Bâle~III, CRR3 art.~166--167).
\end{itemize}

\paragraph{Pourquoi CCF $= 0{,}75$.} Le facteur de conversion de crédit reflète le fait qu'un emprunteur en difficulté tend à tirer sur ses lignes disponibles avant le défaut~: empiriquement, environ 75\,\% de la partie non tirée est consommée. Ce taux est le standard réglementaire Bâle~III pour les facilités engagées.

\subsection{Stress de l'EAD}

En scénario adverse, deux ajustements supplémentaires s'appliquent~:
\begin{enumerate}[nosep]
  \item Le CCF est majoré de 20\,pp (de 0{,}75 à 0{,}95), reflétant l'accélération des tirages en période de stress~;
  \item Les secteurs fragiles (sensibilité au chômage $> 1{,}5$) subissent un multiplicateur de 1{,}05 ($+5\,\%$), capturant l'intensification des besoins de financement dans les secteurs les plus touchés.
\end{enumerate}

\section{Le stress de la PD~: du macro au micro}
\label{sec:ecl:pd-stress}

\subsection{Transformation logit~: garantir les bornes}

Le stress macroéconomique est appliqué à la PD dans l'espace logit, puis retransformé en probabilité par l'expit~:

\begin{equation}
\text{PD}_{\text{stressée}} = \text{expit}\!\left(\text{logit}(\text{PD}_{\text{base}}) + A \cdot \sum_i \epsilon_i \cdot \Delta_i \cdot s_{i,k}\right)
\label{eq:pd-stress}
\end{equation}

\noindent où~:
\begin{itemize}[nosep]
  \item $A \approx 4{,}94$ est l'amplitude logit (calibrée ci-dessous)~;
  \item $\Delta_i = (x_i - \mu_i)/100$ est le choc normalisé de la variable macro $i$~;
  \item $s_{i,k}$ est la sensibilité du secteur $k$ à la variable $i$~;
  \item $\epsilon_i$ encode la direction ($+1$ pour chômage/taux, $-1$ pour PIB/HPI).
\end{itemize}

\paragraph{Pourquoi l'espace logit.} Le logit transforme $[0, 1] \to \mathbb{R}$~: les chocs sont additifs dans cet espace, et la retransformation par l'expit garantit que la PD stressée reste dans $[0, 1]$, même pour des chocs extrêmes. L'alternative (multiplication directe, $\text{PD} \times (1 + \delta)$) risque de produire des PD supérieures à 1 pour les expositions déjà fragiles.

\subsection{Calibration de l'amplitude}

L'amplitude $A$ est calibrée pour reproduire le stress EBA~: une PD de base de 6\,\% doit atteindre 18\,\% ($\times 3$) sous un choc composite de 25\,\%~:

\begin{equation}
A = \frac{\text{logit}(0{,}18) - \text{logit}(0{,}06)}{0{,}25} = \frac{-1{,}506 - (-2{,}752)}{0{,}25} \approx 4{,}94
\label{eq:amplitude}
\end{equation}

Cette calibration n'est pas arbitraire~: elle ancre le modèle sur les résultats de l'EBA \emph{stress test} 2023, où les PD corporate ont typiquement triplé entre le scénario de base et le scénario adverse.

\subsection{Sensibilités sectorielles}

Chaque secteur réagit différemment aux chocs macro, reflétant sa structure économique~:

\begin{table}[htbp]
\centering
\caption{Sensibilités sectorielles aux variables macroéconomiques}
\label{tab:sensibilites}
\begin{tabular}{lccccc}
\toprule
Secteur & Chômage & PIB & Taux & HPI & Inflation \\
\midrule
Technologie & 2{,}5 & 1{,}8 & 1{,}2 & 0{,}5 & 1{,}5 \\
Industrie & 1{,}5 & 2{,}0 & 1{,}0 & 0{,}8 & 1{,}2 \\
Santé & 0{,}5 & 0{,}4 & 0{,}6 & 0{,}3 & 0{,}8 \\
Immobilier & 1{,}0 & 1{,}2 & 2{,}0 & 2{,}5 & 1{,}0 \\
Services & 1{,}2 & 1{,}0 & 0{,}8 & 0{,}5 & 1{,}8 \\
\bottomrule
\end{tabular}
\end{table}

\paragraph{Lecture.} L'immobilier est le secteur le plus sensible aux taux d'intérêt (2{,}0) et au HPI (2{,}5)~: le financement immobilier est structurellement exposé au coût de l'emprunt et à la valeur du collatéral. La technologie est la plus sensible au chômage (2{,}5)~: les startups sont intensives en capital humain, et une hausse du chômage reflète une contraction de la demande pour leurs services. La santé est le secteur le plus défensif~: toutes ses sensibilités sont inférieures à 1, reflétant la demande inélastique de soins.

\subsection{Modèle de Vasicek conditionnel (ASRF)}

Pour l'extension multi-actifs (Partie~II), le stress de la PD utilise le modèle mono-facteur de \citet{vasicek2002}, standardisé par le CRR3 (art.~153)~:

\begin{equation}
\text{PD}_{\text{cond}}(Z) = \Phi\!\left(\frac{\Phi^{-1}(\text{PD}) + \sqrt{\rho} \cdot Z}{\sqrt{1 - \rho}}\right)
\label{eq:vasicek}
\end{equation}

\noindent où $\rho$ est la corrélation d'actifs (spécifique à chaque classe, cf. table~\ref{tab:rho-actifs}), $Z \sim \mathcal{N}(0,1)$ est le facteur systématique, et $\Phi$ la fonction de répartition normale standard.

\begin{table}[htbp]
\centering
\caption{Corrélation d'actifs $\rho$ par classe (CRR3)}
\label{tab:rho-actifs}
\small
\begin{tabular}{lcp{5.5cm}}
\toprule
Classe & $\rho$ & Base réglementaire \\
\midrule
Crédit corporate & 0{,}18 & CRR3 art.~153, PD~$\approx 4{,}5\,\%$ \\
Private equity & 0{,}24 & Traitement equity (moins liquide) \\
Souverain & 0{,}06 & Longstaff (2011), empirique \\
Hypothécaire résidentiel & 0{,}15 & CRR3 art.~154(2)(a), fixe \\
Consommation & 0{,}04 & CRR3 art.~154(2)(b), revolving \\
Interbancaire & 0{,}24 & CRR3 art.~153, risque systémique \\
Titrisation & 0{,}30 & Cliff de corrélation implicite \\
\bottomrule
\end{tabular}
\end{table}

\paragraph{Intuition.} Lorsque $Z > 0$ (stress), le numérateur de l'équation~\eqref{eq:vasicek} augmente, donc $\text{PD}_{\text{cond}} > \text{PD}$~: le facteur systématique amplifie les défauts. L'amplification est d'autant plus forte que $\rho$ est élevé~: un secteur très corrélé au cycle (titrisation, $\rho = 0{,}30$) subit des pertes beaucoup plus concentrées qu'un secteur idiosyncratique (consommation, $\rho = 0{,}04$).

\section{Calcul de l'ECL~: assembler les composantes}
\label{sec:ecl:calcul}

\subsection{Formule générale}

L'ECL par exposition est le produit des trois composantes, actualisé~:

\begin{equation}
\text{ECL} = \text{PD}_{\text{eff}} \times \text{LGD} \times \text{EAD} \times \text{DF}
\label{eq:ecl}
\end{equation}

\noindent avec~:
\begin{itemize}[nosep]
  \item $\text{PD}_{\text{eff}}$~: PD 12 mois pour le Stage~1, PD \emph{lifetime} pour le Stage~2, 100\,\% pour le Stage~3~;
  \item LGD~: TTC pour le scénario de base, \emph{downturn} pour le scénario adverse~;
  \item EAD~: exposition ajustée du CCF~;
  \item DF~: facteur d'actualisation.
\end{itemize}

\subsection{PD \emph{lifetime}~: la méthode du taux de hasard}

La PD \emph{lifetime} extrapole la PD 12 mois sur l'horizon résiduel $T$ (typiquement 3 ans) via le taux de hasard constant~:

\begin{equation}
h = -\ln(1 - \text{PD}_{12\text{m}}), \qquad \text{PD}_{\text{life}} = 1 - e^{-h \cdot T}
\label{eq:pd-lifetime}
\end{equation}

\paragraph{Pourquoi la méthode du taux de hasard.} L'alternative $\text{PD}_{\text{life}} = 1 - (1-\text{PD})^T$ est une approximation acceptable pour les faibles PD, mais sous-estime la PD \emph{lifetime} pour les expositions à haute PD (car elle ignore le taux de hasard instantané). La transformation exponentielle est exacte sous l'hypothèse d'un taux de hasard constant, et garantit $\text{PD}_{\text{life}} \geq \text{PD}_{12\text{m}}$ (monotonie).

\paragraph{Horizon résiduel.} $T = 1$ an pour le Stage~1 (la PD \emph{lifetime} est alors égale à la PD 12 mois)~; $T = 3$ ans pour le Stage~2, conformément à la durée de vie résiduelle moyenne d'un portefeuille corporate non revolving.

\subsection{Facteur d'actualisation}

\paragraph{Stage~1.} Actualisation simple sur un an~:
\begin{equation}
\text{DF}_1 = \frac{1}{1 + r} \approx 0{,}957 \quad \text{avec } r = 4{,}5\,\%
\label{eq:df1}
\end{equation}

Le taux $r = 4{,}5\,\%$ est un proxy du taux d'intérêt effectif (TIE) moyen du portefeuille. IFRS~9 (§5.5.17) exige l'actualisation au TIE de l'instrument~; en pratique, un taux unique est utilisé pour l'ensemble du portefeuille.

\paragraph{Stage~2 et 3.} Le facteur d'actualisation est pondéré par les PD marginales, concentrant l'actualisation sur les années où le défaut est le plus probable~:

\begin{equation}
\text{DF}_{\text{pondéré}} = \frac{\sum_{t=1}^{T} \text{PD}_{\text{marg}}(t) \cdot (1+r)^{-t}}{\sum_{t=1}^{T} \text{PD}_{\text{marg}}(t)}
\label{eq:df-weighted}
\end{equation}

\noindent où $\text{PD}_{\text{marg}}(t) = S(t-1) - S(t) = e^{-h(t-1)} - e^{-ht}$ est la probabilité de défaut dans l'année $t$ conditionnellement à la survie jusqu'en $t-1$.

\paragraph{Intuition.} Pour une exposition à haute PD, la plupart du défaut est concentré dans la première année~: le DF pondéré est proche de $1/(1+r)$. Pour une exposition à faible PD, le défaut est réparti plus uniformément sur $T$ années~: le DF pondéré est plus bas (actualisation sur un horizon plus long). Ce raffinement, conforme à IFRS~9 IE24, évite de surestimer l'ECL des expositions à faible PD et long horizon.

\section{Pondération multi-scénarios}
\label{sec:ecl:scenarios}

IFRS~9 (§B5.5.42) exige que l'ECL soit calculée comme une espérance pondérée par la probabilité de chaque scénario, et non comme le maximum ou le scénario le plus probable~:

\begin{equation}
\text{ECL}_{\text{pondérée}} = \sum_{s \in \mathcal{S}} \pi_s \cdot \text{ECL}_s
\label{eq:ecl-weighted}
\end{equation}

Trois scénarios sont utilisés~:

\begin{table}[htbp]
\centering
\caption{Pondération des scénarios}
\label{tab:scenarios-poids}
\begin{tabular}{lcp{6cm}}
\toprule
Scénario & $\pi$ & Justification \\
\midrule
Base & 50\,\% & Scénario central (tendance zone euro) \\
Adverse & 25\,\% & Stress sévère mais plausible (GFC-like) \\
Favorable & 25\,\% & Conditions favorables (reprise) \\
\bottomrule
\end{tabular}
\end{table}

\paragraph{Pourquoi pas un maximum.} Un calcul basé sur le pire scénario (\emph{worst-case}) violerait IFRS~9, qui exige une \emph{espérance}. De plus, le scénario favorable contribue à l'ECL pondérée~: il n'est pas ignoré. La symétrie des poids (25/25) reflète l'incertitude sur la direction de l'économie.

\paragraph{Interaction LGD--scénario.} La LGD varie selon le scénario~: TTC pour le scénario de base, \emph{downturn} pour le scénario adverse. Le scénario favorable utilise la LGD TTC (pas de réduction~: la norme ne permet pas de provisionner moins que la perte attendue en conditions normales).

\section{Le spread de crédit~: lien avec le marché}
\label{sec:ecl:spread}

Le spread de crédit est calculé comme la prime de risque implicite de l'exposition, selon le modèle de Merton~:

\begin{equation}
\text{spread} = -\frac{\ln(1 - \text{PD} \times \text{LGD})}{T} + \lambda_{\text{liq}}
\label{eq:spread}
\end{equation}

\noindent avec $\lambda_{\text{liq}} = 50$\,pb (prime de liquidité) et un plancher/plafond de 50/2000\,pb. La prime de liquidité reflète le coût de portage des expositions illiquides~: même à PD nulle, une obligation non cotée commande une prime par rapport au taux sans risque.

\paragraph{Exemple.} PD $= 2\,\%$, LGD $= 40\,\%$~:
$$\text{spread} = -\ln(1 - 0{,}008) + 0{,}0050 = 0{,}00803 + 0{,}0050 \approx 130\,\text{pb}$$

\section{Agrégation et métriques de portefeuille}
\label{sec:ecl:agregation}

\subsection{ECL par stage et par secteur}

L'ECL est agrégée en croisant stage et secteur, produisant une matrice $3 \times 5$~:
\begin{itemize}[nosep]
  \item \textbf{Nombre d'expositions}~: granularité du portefeuille~;
  \item \textbf{ECL totale}~: provision nécessaire~;
  \item \textbf{Taux de couverture}~: $\text{ECL} / \text{EAD}$, indicateur de provisionnement~;
  \item \textbf{PD moyenne}~: qualité de crédit du segment.
\end{itemize}

\subsection{Matrice de transition}

Le moteur produit une matrice de transition $3 \times 3$ entre stages~:
$$\mathbf{T} = \begin{pmatrix}
p_{1 \to 1} & p_{1 \to 2} & p_{1 \to 3} \\
p_{2 \to 1} & p_{2 \to 2} & p_{2 \to 3} \\
p_{3 \to 1} & p_{3 \to 2} & p_{3 \to 3}
\end{pmatrix}$$

\noindent Les transitions les plus informatives sont~:
\begin{itemize}[nosep]
  \item $1 \to 2$~: migration SICR (dégradation, provision augmentée)~;
  \item $2 \to 1$~: cure SICR (amélioration, provision réduite)~;
  \item $1 \to 3$~: défaut sans SICR préalable (signal d'alerte~: le modèle SICR a raté la dégradation).
\end{itemize}

\subsection{Coût du risque annualisé}

Le coût du risque rapporte l'ECL totale à l'exposition annualisée~:
\begin{equation}
\text{CoR}_{\text{pb}} = \frac{\text{ECL}_{\text{total}}}{\text{EAD} \times \bar{T}} \times 10\,000
\label{eq:cor}
\end{equation}
où $\bar{T}$ est l'horizon moyen pondéré ($\sim$1{,}3 an pour un portefeuille mixte Stage~1 / Stage~2).

\subsection{Ratio Texas synthétique}

Le ratio Texas compare les expositions en défaut aux provisions et au capital alloué~:
\begin{equation}
\text{Texas} = \frac{\text{EAD}_{\text{Stage 3}}}{\text{ECL}_{\text{Stage 3}} + \text{Capital}_{\text{alloué}}}
\label{eq:texas}
\end{equation}
Un ratio $< 1$ indique que les provisions et le capital couvrent les expositions en défaut~; au-delà de $1{,}5$, la situation est critique.

\section{Synthèse du moteur ECL}
\label{sec:ecl:synthese}

Le moteur ECL assemble six composantes en un pipeline cohérent~:

\begin{enumerate}
  \item \textbf{PD stressée}~: transformation logit avec amplitude calibrée EBA ($A \approx 4{,}94$), sensibilités sectorielles (tableau~\ref{tab:sensibilites}), et modèle de Vasicek pour les classes multi-actifs.

  \item \textbf{Staging SICR}~: score multi-factoriel à quatre composantes (équation~\ref{eq:sicr}), seuil $\tau = 1{,}65$, pondération macro 30\,\% pour assurer la \emph{forward-looking}.

  \item \textbf{LGD}~: base sectorielle $\times$ type de prêt $+$ ajustement score, dispersion Bêta, \emph{downturn} via $\rho_{\text{LGD}}$, canal HPI.

  \item \textbf{EAD}~: CCF 75\,\% (revolving Bâle~III), stress 20\,pp en scénario adverse.

  \item \textbf{Actualisation}~: PD-pondérée pour le Stage~2 (IE24), taux $r = 4{,}5\,\%$.

  \item \textbf{Pondération multi-scénarios}~: espérance 50/25/25 (base/adverse/favorable), conforme B5.5.42.
\end{enumerate}

Chaque paramètre est tracé vers une source normative (IFRS~9, CRR3, EBA GL), un calibrage empirique (stress test EBA 2023, taux de recouvrement observés), ou un raisonnement économique explicite. Le chapitre suivant utilise ce moteur pour construire les métriques de portefeuille --- RAROC, NII, ALM --- qui alimentent les décisions d'allocation.


% ══════════════════════════════════════════════════════════
% CHAPITRE 4 — MÉTRIQUES DE PORTEFEUILLE
% ══════════════════════════════════════════════════════════
\chapter{Métriques de portefeuille~: du risque à la rentabilité}
\label{chap:metriques}

Le moteur ECL du chapitre précédent quantifie la perte attendue par exposition. Mais le CRO ne gère pas des expositions individuelles~: il alloue du capital entre classes d'actifs, secteurs et canaux. Ce chapitre construit les métriques qui transforment les pertes individuelles en indicateurs de décision~: RAROC, NII, EVA, et mesures de concentration.

\section{Revenu net d'intérêts (NII)}
\label{sec:metriques:nii}

\subsection{Formule de base}

Le NII d'une exposition est le produit de l'EAD et du spread de crédit~:

\begin{equation}
\text{NII} = \text{EAD} \times \text{spread}_{\text{crédit}}
\label{eq:nii}
\end{equation}

\noindent avec le spread défini au chapitre précédent (équation~\ref{eq:spread})~:
\begin{equation}
\text{spread} = -\ln(1 - \text{PD} \times \text{LGD}) + \lambda_{\text{liq}} + m_{\text{com}}
\label{eq:spread-nii}
\end{equation}

\noindent où $m_{\text{com}} = 200$\,pb est la marge commerciale. Ce paramètre n'est pas arbitraire~: il correspond au benchmark EBA FINREP 2023 pour les prêts corporate en zone euro (fourchette 150--250\,pb, médiane $\approx$200\,pb).

\paragraph{Décomposition du spread.} Le spread a trois composantes~:
\begin{itemize}[nosep]
  \item $-\ln(1 - \text{PD} \times \text{LGD})$~: la prime de risque de crédit, qui rémunère la perte attendue~;
  \item $\lambda_{\text{liq}} = 50$\,pb~: la prime de liquidité, qui rémunère l'illiquidité de l'actif par rapport au marché monétaire~;
  \item $m_{\text{com}} = 200$\,pb~: la marge commerciale, qui couvre les frais opérationnels, la rémunération du capital, et le profit.
\end{itemize}

\subsection{Impact des taux d'intérêt~: le canal ALM}

En scénario de hausse des taux, le NII est affecté par le \emph{repricing gap}~: les actifs à taux fixe ne reprennent pas immédiatement le nouveau taux de marché, tandis que le coût de refinancement augmente. L'ajustement ALM capture cet effet~:

\begin{equation}
\text{NII}_{\text{ajusté}} = \text{NII} - \text{EAD} \times f_b \times \Delta_{\text{IR}}
\label{eq:nii-alm}
\end{equation}

\noindent avec~:
\begin{equation}
f_b = \min\!\left(\frac{D}{D_{\text{ref}}}, 1\right) \times \beta_{\text{net}}
\label{eq:fb}
\end{equation}

\noindent où $D$ est la duration de l'exposition, $D_{\text{ref}} = 10$ ans (duration de référence d'un portefeuille bancaire long), $\beta_{\text{net}} = 0{,}30$ (taux de transmission net), et $\Delta_{\text{IR}} = (\text{IR}_{\text{scénario}} - \text{IR}_{\text{base}}) / 100$.

\paragraph{Pourquoi $\beta_{\text{net}} = 0{,}30$.} Le taux de transmission net combine deux effets~:
\begin{itemize}[nosep]
  \item Le \textbf{ratio de couverture ALM}~: les grandes banques européennes couvrent 30--50\,\% de leur risque de duration par des swaps de taux (EBA 2024)~;
  \item Le \textbf{bêta des dépôts}~: $\beta_{\text{dépôts}} \approx 0{,}45$ (repricing BCE IMIR 2023), c'est-à-dire que seulement 45\,\% d'une hausse des taux se transmet au coût des dépôts.
\end{itemize}
Le produit $(1 - 0{,}30) \times 0{,}45 \approx 0{,}30$ donne le taux de transmission net~: pour chaque point de taux de hausse, le NII est réduit de 30\,\% de la duration normalisée.

\paragraph{Trois chemins de calcul.} Le code implémente trois variantes~:
\begin{enumerate}[nosep]
  \item \textbf{Crédit corporate}~: formule~\eqref{eq:nii-alm} avec la duration du prêt~;
  \item \textbf{Classes ECL} (multi-actifs)~: même formule, avec la duration spécifique à chaque classe d'actifs (tableau~\ref{tab:durations})~;
  \item \textbf{Private equity}~: pas d'ajustement NII (le PE est valorisé en NAV, pas en flux d'intérêts).
\end{enumerate}

\begin{table}[htbp]
\centering
\caption{Durations par classe d'actifs}
\label{tab:durations}
\small
\begin{tabular}{lcp{5.5cm}}
\toprule
Classe & Duration & Justification \\
\midrule
Crédit corporate & 3{,}0 ans & Ténor 4 ans, prépaiement $\Rightarrow$ 3 ans effectif \\
Private equity & 5{,}0 ans & Illiquide, sensibilité au taux d'actualisation \\
Souverain & 6{,}0 ans & Benchmark 7--10 ans zone euro \\
Hypothécaire & 4{,}0 ans & Ténor 20 ans, CPR $\Rightarrow$ 4 ans effectif \\
Project finance & 7{,}0 ans & Ténor long, pas de prépaiement \\
Covered bonds & 5{,}0 ans & Benchmark Pfandbrief 5 ans \\
Consommation & 2{,}0 ans & Ténor court ($\sim$3 ans) \\
Trade finance & 0{,}5 an & Auto-liquidant, $< 6$ mois \\
Interbancaire & 0{,}25 an & O/N à 3 mois \\
Repos/SFT & 0{,}1 an & Moyenne $\sim$1 mois \\
\bottomrule
\end{tabular}
\end{table}

\subsection{Ajustement du CIR à l'inflation}

Le ratio coût/revenu (CIR) est ajusté en fonction du scénario d'inflation~: 60\,\% des charges opérationnelles sont indexées (50\,\% personnel, 10\,\% locaux)~:

\begin{equation}
\text{CIR}_{\text{eff}} = \text{CIR}_{\text{base}} \times (1 + 0{,}60 \times \Delta_{\text{infl}})
\label{eq:cir-infl}
\end{equation}

\noindent avec $\text{CIR}_{\text{base}} = 45\,\%$ (moyenne EBA FINREP 2023 pour les grandes banques européennes) et $\Delta_{\text{infl}} = (\text{infl}_{\text{scénario}} - \text{infl}_{\text{base}}) / 100$. En stagflation ($\Delta_{\text{infl}} = +3\,\text{pp}$), le CIR passe à $45\,\% \times 1{,}018 = 45{,}8\,\%$, comprimant le revenu net.

\section{RAROC~: le rendement ajusté du risque}
\label{sec:metriques:raroc}

\subsection{Formule}

Le RAROC (\emph{Risk-Adjusted Return on Capital}) est la métrique centrale de l'allocation~:

\begin{equation}
\text{RAROC} = \frac{(\text{NII} \times (1 - \text{CIR}) - \text{EL}) \times (1 - \tau)}{K_{\text{alloué}}}
\label{eq:raroc}
\end{equation}

\noindent avec~:
\begin{itemize}[nosep]
  \item NII~: revenu net d'intérêts (section~\ref{sec:metriques:nii})~;
  \item CIR~: ratio coût/revenu (45\,\%, ajusté de l'inflation)~;
  \item EL~: perte attendue annuelle ($\sum \text{PD}_{\text{stressée}} \times \text{LGD} \times \text{EAD}$)~;
  \item $\tau = 25\,\%$~: taux d'imposition effectif (moyenne zone euro)~;
  \item $K_{\text{alloué}} = \text{RWA} \times \text{CET1}_{\text{cible}}$~: capital alloué.
\end{itemize}

\paragraph{Pourquoi CET1 $= 13\,\%$.} Le ratio cible de 13\,\% comprend~:
\begin{itemize}[nosep]
  \item Minimum CET1 Pilier~1~: 4{,}5\,\% (CRR3)~;
  \item Coussin de conservation~: 2{,}5\,\% (Bâle~III)~;
  \item Coussin contracyclique~: $\sim$1\,\% (BCE, variable)~;
  \item Coussin systémique (G-SIB)~: $\sim$2\,\% (FSB)~;
  \item Buffer SREP~: $\sim$2{,}5\,\% (supervision BCE)~;
  \item Buffer de management~: $\sim$0{,}5\,\%.
\end{itemize}
Le total de $\sim$13\,\% est le niveau effectif des grandes banques européennes (EBA Risk Dashboard 2023).

\subsection{RAROC du private equity}

Le PE est traité différemment du crédit~: le revenu est la performance de la NAV, pas un flux d'intérêts~:

\begin{equation}
\text{RAROC}_{\text{PE}} = \frac{(\text{NAV} \times \text{IRR}_{\text{moyen}} \times (1 - \text{CIR}) - \text{EL}_{\text{PE}}) \times (1 - \tau)}{K_{\text{PE}}}
\label{eq:raroc-pe}
\end{equation}

\noindent avec $\text{IRR}_{\text{moyen}}$ l'IRR moyen du portefeuille PE (stressé par le scénario), et $K_{\text{PE}} = \text{RWA}_{\text{PE}} \times \text{CET1}_{\text{cible}}$.

\paragraph{Contrainte de levier.} Pour les classes à RW nul ou très faible (souverains AAA, repos collatéralisés), le capital alloué est planché par le ratio de levier Bâle~III~:
\begin{equation}
K = \max(\text{RWA} \times \text{CET1}_{\text{cible}}, \text{EAD} \times 3{,}3\,\%)
\label{eq:leverage-floor}
\end{equation}
Le plancher de 3{,}3\,\% empêche d'allouer un capital quasi nul à des expositions volumineuses mais faiblement pondérées.

\section{Pondérations en risque (RWA)}
\label{sec:metriques:rwa}

\subsection{Budget RWA et capital}

Le bilan de référence a un budget RWA de 3{,}69 billions d'euros, soit un capital CET1 de~:
$$K_{\text{total}} = 3{,}69\text{T} \times 13\,\% = 479{,}7\text{Md EUR}$$

Ce budget contraint l'allocation~: les classes à pondération élevée consomment plus de capital par euro d'exposition, réduisant la capacité d'investissement dans les autres classes.

\subsection{Pondérations CRR3 par classe}

\paragraph{Private equity (CRR3 art.~133).} Les pondérations PE sont les plus élevées du bilan, variant de 190\,\% à 400\,\% selon un score de qualité~:

\begin{equation}
\text{Score} = 0{,}30 \times f_{\text{perf}} + 0{,}25 \times f_{\text{qualité}} + 0{,}25 \times f_{\text{levier}} + 0{,}20 \times f_{\text{macro}}
\label{eq:pe-score}
\end{equation}

\begin{itemize}[nosep]
  \item Score $< 0{,}25$~: RW $= 190\,\%$ (PE diversifié, performance IRB)~;
  \item Score $\in [0{,}25, 0{,}50)$~: RW $= 250\,\%$ (equity générale, défaut CRR3)~;
  \item Score $\geq 0{,}50$~: RW $= 400\,\%$ (spéculatif, en difficulté).
\end{itemize}

Ces pondérations pénalisent structurellement le PE~: à RAROC égal, un euro investi en PE consomme 2 à 4 fois plus de capital qu'un euro en crédit corporate ($\sim$100\,\%).

\paragraph{Crédit corporate (CRR3 art.~122 + 501).} La pondération dépend du score de crédit, avec un facteur PME~:

\begin{table}[htbp]
\centering
\caption{Pondérations corporate (CRR3 approche standard)}
\label{tab:rw-corporate}
\begin{tabular}{lcl}
\toprule
Score de crédit & RW & Équivalent rating \\
\midrule
$\geq 750$ & 20\,\% & AAA/AA \\
$680 - 750$ & 50\,\% & A \\
$600 - 680$ & 75\,\% & BBB \\
$500 - 600$ & 100\,\% & BB \\
$< 500$ & 150\,\% & Below BB \\
\midrule
\multicolumn{3}{l}{\emph{Facteur PME (art.~501)}~: $\times 0{,}7619$ si CA $< 50$ M EUR} \\
\multicolumn{3}{l}{\emph{Défaut (art.~122)}~: $\text{RW} = 150\,\%$ pour Stage~3} \\
\bottomrule
\end{tabular}
\end{table}

\paragraph{Hypothécaire résidentiel (CRR3 art.~124--125).} La pondération est fonction du LTV (\emph{Loan-to-Value}), avec un barème progressif~:

\begin{table}[htbp]
\centering
\caption{Barème LTV--RW pour les hypothécaires résidentiels}
\label{tab:rw-ltv}
\begin{tabular}{cc}
\toprule
LTV & RW \\
\midrule
$\leq 50\,\%$ & 20\,\% \\
$50 - 60\,\%$ & 25\,\% \\
$60 - 70\,\%$ & 30\,\% \\
$70 - 80\,\%$ & 35\,\% \\
$80 - 90\,\%$ & 50\,\% \\
$90 - 100\,\%$ & 70\,\% \\
\bottomrule
\end{tabular}
\end{table}

La pondération blendée est calculée sur la distribution LTV du portefeuille (distribution française/européenne typique~: $\sim$33\,\%).

\paragraph{Titrisation (CRR3 art.~242--270, SEC-SA).} Les pondérations dépendent de la séniorité et du label STS~:

\begin{table}[htbp]
\centering
\caption{Pondérations SEC-SA par tranche}
\label{tab:rw-sec}
\begin{tabular}{lc}
\toprule
Tranche & RW \\
\midrule
Senior STS & 20\,\% \\
Senior non-STS & 60\,\% \\
Mezzanine & 100\,\% \\
Junior & 250\,\% \\
\bottomrule
\end{tabular}
\end{table}

\subsection{Dispatch dynamique des pondérations}

Le moteur sélectionne automatiquement le barème approprié~:
\begin{enumerate}[nosep]
  \item Si le profil contient une distribution LTV $\Rightarrow$ barème hypothécaire~;
  \item Si le profil contient un mix de titrisation $\Rightarrow$ barème SEC-SA~;
  \item Sinon $\Rightarrow$ pondération plate (RW fixe du profil).
\end{enumerate}

Ce dispatch est crucial pour l'optimiseur BL-CVaR (chapitre~8)~: il permet d'évaluer la consommation de capital de chaque allocation sans codage spécifique par classe.

\section{EVA~: la valeur économique ajoutée}
\label{sec:metriques:eva}

L'EVA mesure la valeur créée (ou détruite) au-delà du coût du capital~:

\begin{equation}
\text{EVA} = (\text{RAROC} - \text{CET1}_{\text{cible}}) \times K_{\text{alloué}}
\label{eq:eva}
\end{equation}

\begin{itemize}[nosep]
  \item $\text{EVA} > 0$~: l'activité crée de la valeur (le rendement dépasse le coût du capital)~;
  \item $\text{EVA} < 0$~: l'activité détruit de la valeur (mieux vaut réallouer le capital)~;
  \item $\text{EVA} = 0$~: l'activité rémunère exactement le capital investi.
\end{itemize}

\paragraph{Utilisation.} L'EVA est la métrique d'arbitrage entre classes d'actifs~: entre deux classes de même RAROC, celle qui mobilise le plus de capital produit une EVA plus élevée (en valeur absolue). Inversement, une classe à RAROC élevé mais à capital minuscule (par exemple, le trade finance) a un impact limité sur le P\&L consolidé.

\section{Concentration~: l'indice HHI}
\label{sec:metriques:hhi}

\subsection{HHI inter-cellules}

L'indice de Herfindahl-Hirschman \citep{herfindahl1950} mesure la concentration du portefeuille sur 10 cellules (5 secteurs $\times$ 2 canaux~: crédit et PE)~:

\begin{equation}
\text{HHI} = \sum_{i=1}^{10} \left(\frac{E_i}{E_{\text{total}}}\right)^2 \times 10\,000
\label{eq:hhi}
\end{equation}

\noindent où $E_i$ est l'exposition de la cellule $i$ (EAD pour le crédit, NAV pour le PE). Le facteur 10\,000 normalise l'indice entre 100 (répartition uniforme sur 100 cellules) et 10\,000 (concentration totale sur une cellule).

\begin{table}[htbp]
\centering
\caption{Seuils d'appétit pour le risque --- concentration HHI}
\label{tab:hhi-seuils}
\begin{tabular}{ccl}
\toprule
HHI & Signal & Interprétation \\
\midrule
$< 1\,500$ & Vert & Portefeuille bien diversifié \\
$1\,500 - 2\,500$ & Ambre & Concentration modérée \\
$\geq 2\,500$ & Rouge & Concentration excessive \\
\bottomrule
\end{tabular}
\end{table}

\subsection{HHI de concentration de noms}

Un second HHI est calculé au niveau des contreparties individuelles~:
\begin{equation}
\text{HHI}_{\text{noms}} = \sum_{j=1}^{N} \left(\frac{\text{EAD}_j}{\text{EAD}_{\text{total}}}\right)^2 \times 10\,000
\label{eq:hhi-noms}
\end{equation}

Les seuils sont plus stricts~: vert $< 30$, ambre $< 50$, rouge $\geq 50$. Un HHI de noms élevé signale un risque idiosyncratique de concentration (ICAAP, Pilier~2).

\subsection{Pourquoi le HHI plutôt qu'une autre mesure}

Le HHI est le standard réglementaire de mesure de concentration~:
\begin{itemize}[nosep]
  \item Il est \emph{convexe}~: la concentration est pénalisée de manière quadratique (une exposition représentant 50\,\% du portefeuille contribue 25$\times$ plus qu'une exposition de 10\,\%)~;
  \item Il est \emph{additif}~: le HHI du portefeuille est la somme des carrés des parts~;
  \item Il est \emph{borné}~: entre $10\,000/N$ (uniforme) et $10\,000$ (concentré)~;
  \item Il est utilisé par les autorités de concurrence (Commission européenne, FTC) et par les régulateurs bancaires (EBA, FSB).
\end{itemize}

\section{CET1 et contrainte de capital}
\label{sec:metriques:cet1}

\subsection{La contrainte qui lie}

Le ratio CET1 est la contrainte qui domine l'allocation~:
\begin{equation}
\text{CET1}_{\text{ratio}} = \frac{K_{\text{CET1}}}{\sum_i w_i \times \text{RW}_i \times E_i} \geq 13\,\%
\label{eq:cet1}
\end{equation}

\noindent où $w_i$ est le poids de la classe $i$ dans le portefeuille et $E_i$ son exposition. La contrainte lie dans 10 scénarios sur 11~: seule la «~Trappe à liquidité~» (taux négatifs $\Rightarrow$ RWA réduits) passe sans ajustement.

\paragraph{Pas de double déduction ECL.} Le capital CET1 est calculé comme $\text{RWA}_{\text{budget}} \times \text{CET1}_{\text{cible}}$, \emph{sans} déduire l'ECL. En effet, l'ECL est déjà incorporée dans les RWA via les pondérations de risque (CRR3 SA)~: la déduire du numérateur reviendrait à la compter deux fois.

\subsection{Interaction PE--CET1}

Le PE est la classe la plus consommatrice de CET1~: avec des RW de 190--400\,\%, chaque euro de NAV PE consomme 2 à 5 fois plus de capital qu'un euro de crédit corporate. Toutefois, le split crédit/PE est \textbf{endogène}~: l'optimiseur BL-CVaR ajuste la part PE entre $\sim$18\,\% (scénarios neutres et favorables) et $\sim$62\,\% (crise souveraine), selon le différentiel de RAROC entre canaux~:
\begin{itemize}[nosep]
  \item Le PE est utilisé comme \emph{diversificateur} quand le RAROC crédit est négatif (GFC~$\to$ PE~38\,\%, Souveraine~$\to$ PE~62\,\%)~;
  \item En stress symétrique (Stagflation, les deux RAROC négatifs), l'allocation reste majoritairement crédit (78\,\%)~;
  \item Le CET1 reste le plafond naturel~: même à 62\,\% de PE, le ratio CET1 est de 23{,}3\,\% (bien au-dessus des 13\,\% réglementaires).
\end{itemize}

Ce résultat, obtenu après suppression des caps de marché rigides qui écrasaient la sortie de l'optimiseur, est documenté quantitativement au chapitre~6 (Partie~I) et validé à l'échelle des 14 classes d'actifs en Partie~II.

\section{Synthèse~: le tableau de bord du CRO}
\label{sec:metriques:synthese}

Les métriques s'organisent en quatre familles~:

\begin{table}[htbp]
\centering
\caption{Synthèse des métriques de portefeuille}
\label{tab:metriques-synthese}
\small
\begin{tabular}{llll}
\toprule
Famille & Métrique & Formule & Seuils \\
\midrule
Rentabilité & RAROC & $(R_{\text{net}} - \text{EL}) \times (1-\tau) / K$ & $> 13\,\%$ vert \\
 & EVA & $(\text{RAROC} - 13\,\%) \times K$ & $> 0$ vert \\
\midrule
Risque & ECL totale & $\sum \text{PD} \times \text{LGD} \times \text{EAD}$ & $-$ \\
 & Coût du risque & $\text{ECL} / (\text{EAD} \times \bar{T}) \times 10^4$ & $-$ \\
 & Texas ratio & $\text{EAD}_{S3} / (\text{ECL}_{S3} + K)$ & $< 1$ vert \\
\midrule
Capital & CET1 ratio & $K / \text{RWA}$ & $\geq 13\,\%$ \\
 & Ratio de levier & $K / \text{EAD}_{\text{total}}$ & $\geq 3{,}3\,\%$ \\
\midrule
Concentration & HHI inter-cellules & $\sum s_i^2 \times 10^4$ & $< 1\,500$ vert \\
 & HHI de noms & $\sum s_j^2 \times 10^4$ & $< 30$ vert \\
\bottomrule
\end{tabular}
\end{table}

Ces métriques sont calculées pour chaque scénario et chaque allocation, formant une surface de décision multidimensionnelle. Le chapitre suivant applique ce cadre au private equity, où les contraintes CRR3 et l'illiquidité IFRS~13 créent des tensions spécifiques avec le crédit.


% ══════════════════════════════════════════════════════════
% CHAPITRE 5 — PRIVATE EQUITY SOUS CONTRAINTES
% ══════════════════════════════════════════════════════════
\chapter{Private equity sous contraintes réglementaires}
\label{chap:pe}

Le private equity occupe une place singulière dans le bilan bancaire~: c'est la classe d'actifs à la fois la plus rentable (IRR cible de 12--20\,\%) et la plus coûteuse en capital (RW de 190 à 400\,\%). Ce chapitre analyse les trois contraintes qui façonnent l'arbitrage crédit--PE~: la valorisation IFRS~13, les pondérations CRR3, et l'illiquidité structurelle.

\section{Valorisation~: trois canaux de stress}
\label{sec:pe:valorisation}

\subsection{Architecture du modèle}

La NAV d'une position PE est le produit de trois composantes, chacune soumise à un canal de stress distinct~:

\begin{equation}
\text{NAV}_i = \underbrace{M_i^{\text{stressé}}}_{\text{Canal 1~: métrique}} \times \underbrace{m_{\text{exit}}^{\text{comprimé}}}_{\text{Canal 2~: multiple}} \times \underbrace{(1 - L_i^{\text{stressé}})}_{\text{Canal 3~: levier}} \times \varepsilon_i
\label{eq:nav}
\end{equation}

\noindent Les trois canaux capturent des risques économiquement distincts~:

\paragraph{Canal 1~: stress de la métrique opérationnelle.} L'EBITDA (ou le chiffre d'affaires pour la technologie) est stressé par les conditions macro~:
\begin{equation}
M^{\text{stressé}} = M \times \exp(-S \times 3{,}0)
\label{eq:metric-stress}
\end{equation}
avec $S = \delta_{\text{PIB}} \cdot s^{\text{PE}}_{\text{PIB}} + \delta_{\text{chômage}} \cdot s^{\text{PE}}_{\text{chômage}} + \delta_{\text{infl}} \cdot s^{\text{PE}}_{\text{infl}}$. Le facteur d'échelle 3{,}0 est calibré pour qu'une baisse de 3\,pp du PIB produise une réduction de l'EBITDA de 5--8\,\%, conformément aux données Preqin 2023 sur les fonds buyout européens.

\paragraph{Canal 2~: compression du multiple de sortie.} Le multiple d'exit est comprimé par les taux d'intérêt et la conjoncture~:
\begin{equation}
m_{\text{exit}}^{\text{comprimé}} = m_{\text{base}} \times \exp(-C \times 5{,}0)
\label{eq:multiple-compression}
\end{equation}
avec $C = \delta_{\text{IR}} \cdot s^{\text{PE}}_{\text{IR}} + \delta_{\text{PIB}} \cdot s^{\text{PE}}_{\text{PIB}} + \delta_{\text{HPI}} \cdot s^{\text{PE}}_{\text{HPI}}$. Le facteur 5{,}0 est calibré pour qu'une hausse de 100\,pb des taux comprime le multiple de 3--5\,\%, cohérent avec les stress tests EBA 2023.

\paragraph{Canal 3~: stress du levier.} En hausse des taux, le levier effectif augmente (charge de la dette accrue)~:
\begin{equation}
L^{\text{stressé}} = \text{clip}(L + \delta_{\text{IR}} \cdot s^{\text{PE}}_{\text{IR}} \times 0{,}5, \; 0, \; 0{,}95)
\label{eq:leverage-stress}
\end{equation}
Le facteur 0{,}5 implique qu'une hausse de 100\,pb augmente le levier de 0{,}5--1\,pp. Le plafond à 95\,\% empêche les NAV négatives.

\paragraph{Dispersion intra-sectorielle.} Un terme de bruit $\varepsilon_i$ capture l'hétérogénéité des fonds au sein d'un secteur, avec une structure factorielle~:
$$\varepsilon_i = 1 + \sigma(\sqrt{\rho} \cdot Z_s + \sqrt{1-\rho} \cdot Z_i)$$
avec $\sigma = 3\,\%$, $\rho = 0{,}40$ (corrélation intra-secteur). Ce modèle est l'analogue PE de l'ASRF de \citet{vasicek2002}~: 40\,\% du risque est systématique (secteur), 60\,\% est idiosyncratique (fonds).

\subsection{Deux corrections asymétriques}

Le modèle intègre deux corrections qui brisent la symétrie favorable/adverse~:

\paragraph{Amortissement de récession (85\,\%).} Lorsque le PIB baisse mais que les taux baissent aussi (réponse monétaire accommodante), le bénéfice des taux bas sur les multiples est annulé à 85\,\%~:
$$C_{\text{ajusté}} = C - 0{,}85 \times \underbrace{(\delta_{\text{IR}} \cdot s^{\text{PE}}_{\text{IR}})}_{\text{bénéfice taux}} \times \min\!\left(\frac{\delta_{\text{PIB}}}{0{,}02}, 1\right)$$

\paragraph{Pourquoi 85\,\%.} En récession, la baisse des taux ne se transmet pas aux multiples PE~: les investisseurs exigent une prime de risque accrue, les transactions se tarissent, et les comités d'investissement appliquent des décotes. L'EBA Stress Test 2023 montre que les multiples PE chutent même lorsque les taux baissent en scénario adverse. Les 15\,\% restants reflètent le cas marginal où un fonds bien positionné (dette à taux variable, dry powder) profite effectivement des taux bas.

\paragraph{Compensation de croissance (50\,\%).} Inversement, lorsque le PIB croît mais que les taux montent (resserrement monétaire), la compression du multiple est partiellement compensée par la croissance des métriques~:
$$C_{\text{ajusté}} = C - 0{,}50 \times \underbrace{(\delta_{\text{IR}} \cdot s^{\text{PE}}_{\text{IR}})}_{\text{compression taux}} \times \min\!\left(\frac{|\delta_{\text{PIB}}|}{0{,}03}, 1\right)$$

Les 50\,\% reflètent l'observation empirique que la croissance de l'EBITDA compense en partie la hausse du coût du capital, mais pas intégralement (le re-rating du risque prend du temps).

\subsection{IRR~: du MOIC au rendement annualisé}

L'IRR est dérivé du multiple MOIC (\emph{Multiple on Invested Capital})~:
\begin{equation}
\text{IRR} = \text{MOIC}^{1/T_h} - 1, \qquad \text{MOIC} = \frac{\text{NAV}}{\text{Capital investi}}
\label{eq:irr}
\end{equation}
avec $T_h$ la durée de détention. L'IRR est borné à $[-50\,\%, +35\,\%]$~: le plafond de 35\,\% correspond au top décile Cambridge Associates 2023 pour les fonds buyout européens~; le plancher de $-50\,\%$ borne les pertes extrêmes.

\paragraph{Simplification.} Ce calcul est une simplification~: il ignore la courbe en J (appels de capitaux progressifs), les distributions intermédiaires, et le \emph{timing} des flux. Pour un outil d'allocation stratégique (horizon 1 an), cette approximation est suffisante.

\section{Probabilité de détresse~: le signal PE}
\label{sec:pe:distress}

\subsection{Modèle logit}

La probabilité de détresse d'une position PE est modélisée dans l'espace logit, avec trois composantes~:

\begin{equation}
P_{\text{distress}} = \text{expit}\!\left(\underbrace{\text{logit}(0{,}05)}_{\text{base}} + \underbrace{A_{\text{MOIC}}}_{\text{performance}} + \underbrace{A_{\text{macro}} \times 3{,}0}_{\text{cycle}}\right)
\label{eq:p-distress}
\end{equation}

\paragraph{Base $= 5\,\%$.} Le taux de détresse de base est calibré sur Preqin (2022)~: environ 8\,\% des fonds PE sont classés «~en difficulté~» à un instant donné, mais en incluant les fonds en restructuration réussie, le taux effectif est de $\sim$5\,\%. Cambridge Associates rapporte 3--5\,\% de \emph{write-offs} sur la durée de vie d'un fonds.

\paragraph{Ajustement MOIC (logarithmique).}
$$A_{\text{MOIC}} = \begin{cases} -2 \ln(\max(\text{MOIC}, 0{,}01)) & \text{si MOIC} < 1 \\ 0 & \text{si MOIC} \geq 1 \end{cases}$$
Un MOIC de 0{,}5 ($\Rightarrow A = 1{,}4$) double approximativement la probabilité de détresse~; un MOIC de 0{,}1 ($\Rightarrow A = 4{,}6$) la porte à $\sim$50\,\%. Le seuil à MOIC $= 1$ (capital retourné) annule la pénalité~: un fonds qui a remboursé son capital n'est pas en détresse.

\paragraph{Ajustement macro.}
$$A_{\text{macro}} = \sum_i \delta_i \cdot s^{\text{PE}}_i$$
Le facteur d'échelle 3{,}0 (inférieur au 4{,}94 du crédit) reflète le fait que le PE est une tranche equity~: les sensibilités sont déjà plus élevées que le crédit, et une échelle supplémentaire trop forte produirait des probabilités de détresse irréalistes.

\subsection{Classification et perte attendue}

\begin{table}[htbp]
\centering
\caption{Classification PE et perte attendue}
\label{tab:pe-classif}
\begin{tabular}{lccl}
\toprule
Catégorie & Seuil $P_{\text{distress}}$ & $\text{EL}_{\text{PE}}$ & Source \\
\midrule
Performing & $< 10\,\%$ & $P \times 0{,}60 \times \text{NAV}$ & Preqin (2022)~: $\sim$8\,\% distressed \\
Watchlist & $10 - 30\,\%$ & $P \times 0{,}60 \times \text{NAV}$ & Preqin~: $\sim$25\,\% watchlist \\
Distressed & $\geq 30\,\%$ & $P \times 0{,}60 \times \text{NAV}$ & Write-off probable \\
\bottomrule
\end{tabular}
\end{table}

La LGD equity de 60\,\% reflète le taux de recouvrement moyen des investissements en equity~: Moody's (2023) rapporte un taux de recouvrement de ${\sim}\,40\,\%$ pour les tranches equity en restructuration.

\section{Pondérations CRR3~: le coût du capital PE}
\label{sec:pe:rw}

\subsection{Score composite}

Le CRR3 (art.~133) impose des pondérations de 190\,\%, 250\,\% ou 400\,\% pour les expositions en actions non cotées. L'attribution de la pondération utilise un score composite à quatre facteurs~:

\begin{equation}
\text{Score} = 0{,}30 \times F + 0{,}25 \times Q + 0{,}25 \times L + 0{,}20 \times E
\label{eq:pe-rw-score}
\end{equation}

\noindent avec~:
\begin{itemize}[nosep]
  \item $F = \text{clip}(1 - \text{MOIC}, 0, 1)$ si MOIC $< 1{,}5$, sinon $0$ --- performance financière~;
  \item $Q = P_{\text{distress}}$ --- qualité du portefeuille~;
  \item $L = \text{clip}(\text{levier}/0{,}95, 0, 1)$ --- risque de levier~;
  \item $E = (L + F)/2$ --- environnement macro (moyenne des deux facteurs les plus cycliques).
\end{itemize}

\begin{table}[htbp]
\centering
\caption{Barème RW PE (CRR3 art.~133)}
\label{tab:rw-pe}
\begin{tabular}{cll}
\toprule
Score & RW & Catégorie CRR3 \\
\midrule
$< 0{,}25$ & 190\,\% & Equity diversifié (IRB) \\
$[0{,}25, 0{,}50)$ & 250\,\% & Equity générale (défaut SA) \\
$\geq 0{,}50$ & 400\,\% & Equity spéculatif \\
\bottomrule
\end{tabular}
\end{table}

\paragraph{Pourquoi $E = (L + F)/2$.} Utiliser directement $E = Q$ (probabilité de détresse) créerait un double comptage~: le score $Q$ est déjà fonction du levier et de la performance. La moyenne $(L + F)/2$ capture l'environnement sans redondance.

\subsection{Impact sur l'allocation}

Avec des RW de 190--400\,\%, le PE consomme 2 à 4 fois plus de capital qu'un crédit corporate standard ($\sim$100\,\%)~:

\begin{equation}
\frac{K_{\text{PE}}}{K_{\text{crédit}}} = \frac{\text{RW}_{\text{PE}}}{\text{RW}_{\text{crédit}}} = \frac{250\,\%}{100\,\%} = 2{,}5
\label{eq:pe-capital-ratio}
\end{equation}

Pour maintenir le ratio CET1 à 13\,\%, chaque euro de NAV PE à 250\,\% de RW «~bloque~» 32{,}5 centimes de capital (vs 13 centimes pour le crédit). Ce surcoût en capital limite structurellement la part du PE dans le bilan.

\section{Illiquidité~: le DLOM}
\label{sec:pe:dlom}

\subsection{Décote pour absence de marché}

Le DLOM (\emph{Discount for Lack of Marketability}) capture le coût de sortie anticipée d'un fonds PE. Contrairement aux obligations cotées, les parts de fonds PE ne sont négociables que sur un marché secondaire illiquide, avec des décotes significatives.

Le DLOM effectif dépend de la maturité du fonds~:

\begin{equation}
\text{DLOM}_{\text{eff}} = d_{\text{sec}} \times \left(1 + f_v \times \frac{\max(0, T_v - t)}{T_v}\right)
\label{eq:dlom}
\end{equation}

\noindent avec~:
\begin{itemize}[nosep]
  \item $d_{\text{sec}} = 10\,\%$~: décote de base sur le marché secondaire (Jefferies/Lazard 2023, médiane buyout 8--12\,\%)~;
  \item $f_v = 0{,}50$~: surcharge pour les fonds récents (jeunes vintages plus difficiles à vendre)~;
  \item $T_v = 5$ ans~: seuil de maturité (au-delà, le fonds est considéré «~mature~»)~;
  \item $t$~: ancienneté du fonds.
\end{itemize}

\paragraph{Exemples.}
\begin{itemize}[nosep]
  \item Fonds de 5 ans (mature)~: $\text{DLOM} = 10\,\%$~;
  \item Fonds de 2{,}5 ans~: $\text{DLOM} = 10\,\% \times 1{,}25 = 12{,}5\,\%$~;
  \item Fonds récent (0 an)~: $\text{DLOM} = 10\,\% \times 1{,}50 = 15\,\%$.
\end{itemize}

Le coût de sortie est alors~:
$$\text{Exit Cost} = \text{NAV} \times (1 - \text{DLOM}_{\text{eff}})$$

\paragraph{Sources.} L'AICPA (2013) et \citet{pratt_grabowski2014} documentent des DLOM de 15--30\,\% pour les PE illiquides. Notre fourchette de 10--15\,\% est conservatrice~: elle reflète le développement du marché secondaire PE depuis 2015 (Jefferies/Lazard 2023 reportent des volumes record de transactions secondaires).

\section{L'arbitrage crédit--PE~: un levier étroit}
\label{sec:pe:arbitrage}

\subsection{Trois contraintes qui se modulent}

L'arbitrage entre crédit et PE est encadré par trois mécanismes~:

\begin{enumerate}
  \item \textbf{CRR3}~: les RW de 190--400\,\% pénalisent fortement le PE en capital. Néanmoins, le CET1 reste satisfait même à des niveaux élevés de PE (23{,}3\,\% pour 62\,\% de PE en scénario Souveraine), grâce à l'excédent de capital du portefeuille.

  \item \textbf{IFRS~13}~: l'illiquidité (DLOM 10--15\,\%) réduit la NAV réalisable, comprimant le RAROC effectif par rapport au RAROC comptable.

  \item \textbf{Asymétrie des corrections}~: l'amortissement de récession (85\,\%) et la compensation de croissance (50\,\%) rendent le profil rendement/risque du PE \emph{asymétrique}~: les pertes en scénario adverse sont plus sévères que les gains en scénario favorable.
\end{enumerate}

\subsection{Deux leviers complémentaires}

Le constat central de cette Partie~I est que l'optimisation repose sur \textbf{deux leviers complémentaires}~: l'\emph{arbitrage crédit--PE} (split inter-canal) et l'\emph{allocation sectorielle} (intra-canal). Le split crédit/PE est endogène via l'optimiseur BL-CVaR et varie de 18\,\% à 62\,\% de PE selon le scénario (6 régimes distincts). Quand le RAROC crédit est négatif mais le RAROC PE reste positif (GFC, Souveraine), l'optimiseur concentre le PE~; en stress symétrique (Stagflation), il reste majoritairement crédit. Simultanément, l'allocation sectorielle au sein de chaque canal reste un déterminant majeur~: les cinq secteurs ont des sensibilités macro très différentes (tableau~\ref{tab:sensibilites}), et réallouer entre technologie et santé produit des variations de RAROC significatives sans pénalité de capital supplémentaire.

Ce résultat est validé par les 12 scénarios historiquement calibrés~: les allocations sectorielles sont scénario-réactives (écart-type moyen de 3{,}6\,\% en crédit et 7{,}4\,\% en PE), avec un alignement RAROC--poids de 71\,\% (17/24 cellules).

\section{Sensibilités sectorielles PE}
\label{sec:pe:sensibilites}

\begin{table}[htbp]
\centering
\caption{Paramètres PE par secteur}
\label{tab:pe-secteurs}
\small
\begin{tabular}{lccccccl}
\toprule
Secteur & Part & $s_{\text{PIB}}$ & $s_{\text{chôm}}$ & $s_{\text{IR}}$ & $s_{\text{HPI}}$ & $s_{\text{infl}}$ & Méthode \\
\midrule
Technologie & 25\,\% & 1{,}0 & $-$1{,}5 & 2{,}5 & 0{,}3 & 0{,}8 & EV/CA \\
Industrie & 25\,\% & 2{,}5 & 1{,}0 & 1{,}5 & 0{,}5 & 1{,}0 & EV/EBITDA \\
Santé & 15\,\% & 0{,}5 & 0{,}3 & 0{,}5 & 0{,}2 & 0{,}5 & EV/EBITDA \\
Immobilier & 20\,\% & 1{,}0 & 0{,}8 & 2{,}0 & 3{,}0 & 0{,}5 & Cap rate/NOI \\
Services & 15\,\% & 1{,}2 & 1{,}0 & 1{,}0 & 0{,}3 & 2{,}0 & EV/EBITDA \\
\bottomrule
\end{tabular}
\end{table}

\paragraph{Lecture.} La technologie a une sensibilité au chômage \emph{négative} ($-1{,}5$)~: une hausse du chômage peut bénéficier aux entreprises tech (automatisation, réduction des coûts salariaux). En revanche, elle est très sensible aux taux (2{,}5)~: les multiples tech sont les plus exposés aux variations du taux d'actualisation. L'immobilier cumule sensibilité aux taux (2{,}0) et au HPI (3{,}0)~: c'est le secteur PE le plus vulnérable en scénario de crise immobilière.

\paragraph{Méthode d'évaluation.} Chaque secteur utilise la méthode IPEV \citep{ipev2022} la plus adaptée~: EV/CA pour la technologie (entreprises pré-profit), EV/EBITDA pour l'industrie et les services, cap rate/NOI pour l'immobilier (après déduction du ratio de charges CBRE/JLL 2023 de 15\,\%).

\section{Synthèse}
\label{sec:pe:synthese}

Le PE offre un rendement attractif mais sous trois contraintes~:
\begin{enumerate}[nosep]
  \item \textbf{Capital}~: RW 190--400\,\% (CRR3 art.~133), fortement consommateurs de CET1~;
  \item \textbf{Valorisation}~: NAV stressée par trois canaux (métrique, multiple, levier) avec des corrections asymétriques (récession 85\,\%, croissance 50\,\%)~;
  \item \textbf{Liquidité}~: DLOM de 10--15\,\% (Jefferies/Lazard 2023), dépendant du vintage.
\end{enumerate}

Le résultat central est que l'\emph{arbitrage inter-classes} (crédit vs PE) est un levier significatif quand le capital est suffisant~: le split PE varie de 18\,\% à 62\,\% selon le scénario, l'optimiseur concentrant le PE dans les scénarios où le RAROC PE domine le RAROC crédit. L'\emph{arbitrage sectoriel} reste complémentaire et constitue le levier principal dans les scénarios de stress symétrique. Les deux leviers sont validés au chapitre suivant par l'analyste CRO et les moteurs de gouvernance.


% ══════════════════════════════════════════════════════════
% CHAPITRE 6 — ANALYSTE CRO ET GOUVERNANCE
% ══════════════════════════════════════════════════════════
\chapter{Analyste CRO et moteurs de gouvernance}
\label{chap:gouvernance}

Les chapitres précédents ont construit un pipeline complet~: données $\to$ PD $\to$ ECL $\to$ métriques $\to$ allocation. Mais un pipeline sans garde-fous est un risque en soi. Ce chapitre présente les dix moteurs de gouvernance qui encadrent les résultats du pipeline, et l'analyste CRO local qui synthétise l'ensemble en un rapport actionnable.

\section{L'analyste CRO local}
\label{sec:gouv:cro}

\subsection{Architecture}

L'analyste CRO est un moteur d'analyse déterministe (sans modèle de langage) qui produit un rapport de gouvernance en huit sections. Il prend en entrée l'état du portefeuille sous un scénario macroéconomique donné et génère~:

\begin{itemize}[nosep]
  \item Un \textbf{score de risque} sur 10 (agrégation pondérée de l'ECL, du staging, de la dérive, et des stress macro)~;
  \item Un \textbf{diagnostic} en prose, calibré sur le scénario (pas de template figé)~;
  \item Des \textbf{alertes} classées par sévérité (critique $>$ alerte $>$ avertissement)~;
  \item Des \textbf{recommandations} priorisées (haute $>$ moyenne $>$ standard).
\end{itemize}

\subsection{Huit sections du rapport}

\begin{enumerate}
  \item \textbf{Synthèse exécutive}~: score de risque, niveau (élevé/modéré/bon), couleur (rouge/orange/vert). Un score $\geq 5$ est «~élevé~»~: il se déclenche si l'ECL varie de plus de 15\,\%, si le Stage~2 dépasse 20\,\%, ou si le PSI dépasse 0{,}25.

  \item \textbf{Attribution de l'ECL}~: décomposition en cascade (\emph{waterfall}) des contributions par segment et par facteur. Identifie le segment dominant et le facteur macro le plus explicatif.

  \item \textbf{Sensibilité macro}~: pour chaque variable adverse, identifie le secteur le plus sensible, le mécanisme de transmission, et l'impact estimé. Si deux variables ou plus sont adverses simultanément, signale les boucles de rétroaction.

  \item \textbf{Analyse du staging}~: distribution Stage~1/2/3, matrice de transition. Une migration $1 \to 2$ supérieure à 5\,\% est qualifiée de «~significative~»~; une migration $2 \to 3$ supérieure à 5\,\% signale un passage en défaut massif.

  \item \textbf{Gouvernance des modèles}~: comparaison AUC/Gini/KS entre modèles, alerte PSI (dérive de distribution), backtesting \emph{walk-forward}, cohérence SHAP.

  \item \textbf{Concentration}~: HHI par segment et par nom, alertes si le segment dominant dépasse 40\,\% de l'ECL.

  \item \textbf{Perspectives \emph{forward-looking}}~: risques émergents (chômage $> 8\,\%$, taux $> 4\,\%$, HPI $< 0$), estimation du pipeline Stage~2 $\to$ Stage~3.

  \item \textbf{Recommandations}~: plan d'action priorisé. Une escalade vers la direction des risques est recommandée sous 48h si l'ECL varie de plus de 20\,\%.
\end{enumerate}

\subsection{Modèle d'élasticité}

L'analyste estime l'impact de chaque variable macro sur chaque segment via un multiplicateur PD~:
\begin{equation}
\hat{m}_{s,v} = 1 + \frac{\text{choc}_v}{100} \times s_{v,k} \times 10
\label{eq:elasticity}
\end{equation}
où $\text{choc}_v$ est la déviation de la variable $v$ par rapport au scénario de base, et $s_{v,k}$ la sensibilité du secteur $k$. Ce modèle linéaire est une approximation locale de la transformation logit~: il surestime légèrement l'impact pour les grands chocs, mais offre une interprétabilité immédiate au CRO.

\section{Prédiction conforme~: intervalles sans hypothèse}
\label{sec:gouv:conformal}

La prédiction conforme produit des intervalles de prédiction avec une garantie de couverture $1 - \alpha$, sans hypothèse distributionnelle~:

\begin{equation}
P\!\left(\text{PD}_{\text{réelle}} \in [\hat{\text{PD}} - \hat{q}, \hat{\text{PD}} + \hat{q}]\right) \geq 1 - \alpha
\label{eq:conformal}
\end{equation}

\noindent où $\hat{q}$ est le quantile $\lceil(n+1)(1-\alpha)\rceil / n$ des scores de non-conformité $|\hat{\text{PD}} - \text{PD}_{\text{réelle}}|$ calculés sur un ensemble de calibration.

\paragraph{Pourquoi la prédiction conforme.} Contrairement aux intervalles de confiance paramétriques (qui supposent une distribution gaussienne des erreurs), la prédiction conforme est \emph{distribution-free}~: elle garantit la couverture marginale même si le modèle est mal spécifié. Pour IFRS~9, cela permet un staging conservateur~: en utilisant la borne supérieure $\hat{\text{PD}} + \hat{q}$ pour le classement en Stage~2, on garantit que les expositions à risque ne sont pas sous-classées.

\section{Analyse de sensibilité de Sobol}
\label{sec:gouv:sobol}

L'analyse de Sobol décompose la variance de l'ECL en contributions de chaque variable macro et de leurs interactions~:

\begin{equation}
\text{Var}(Y) = \sum_i V_i + \sum_{i < j} V_{ij} + \cdots
\label{eq:sobol-anova}
\end{equation}

Les indices de premier ordre $S^1_i = V_i / \text{Var}(Y)$ mesurent l'effet direct~; les indices totaux $S^T_i = 1 - V_{\sim i}/\text{Var}(Y)$ incluent toutes les interactions de $X_i$ avec les autres variables.

\paragraph{Échantillonnage Saltelli.} La méthode de \citet{saltelli2002} utilise $N(2D+2)$ évaluations quasi-aléatoires, avec $N = 512$ et $D = 5$ variables, soit 3\,584 passages du pipeline ECL. Chaque passage prend $\sim$1{,}1 seconde~; l'analyse complète dure $\sim$1 heure.

\paragraph{Résultats typiques.} Le chômage explique $\sim$35\,\% de la variance de l'ECL, le PIB $\sim$25\,\%, les taux et le HPI $\sim$15\,\% chacun, l'inflation $\sim$10\,\%. Les interactions (indices $S^2$) représentent $\sim$10\,\%, principalement entre chômage et taux d'intérêt (loi d'Okun).

\section{Détection de régime (HMM)}
\label{sec:gouv:hmm}

Un modèle de Markov caché à trois états classifie le régime macroéconomique courant~:

\begin{table}[htbp]
\centering
\caption{Régimes HMM et paramètres conditionnels}
\label{tab:hmm-regimes}
\begin{tabular}{lccc}
\toprule
Paramètre & Contraction & Recovery & Expansion \\
\midrule
Poids Base/Adverse/Favorable & 35/45/20 & 50/25/25 & 45/15/40 \\
Multiplicateur $\tau$ (BL) & 0{,}60 & 1{,}00 & 1{,}20 \\
$\alpha$ CVaR & 0{,}80 & 0{,}95 & 0{,}95 \\
Scaling $\Omega$ (incertitude) & 2{,}5 & 1{,}0 & 0{,}8 \\
\bottomrule
\end{tabular}
\end{table}

\paragraph{Pourquoi trois régimes.} Hamilton (1989) montre que trois états (contraction, récupération, expansion) capturent les principales dynamiques macroéconomiques sans sur-paramétrer. La matrice de transition est fortement persistante (95\,\% d'auto-transition)~: les régimes changent lentement, conformément aux données du NBER.

\paragraph{Impact sur le pipeline.} En contraction, le poids du scénario adverse passe de 25\,\% à 45\,\%, le paramètre $\tau$ de Black-Litterman est réduit à 0{,}60 (plus conservateur), et le CVaR utilise un quantile plus profond ($\alpha = 0{,}80$, soit les 20\,\% pires cas). Ces ajustements rendent l'optimiseur plus défensif en période de crise.

\section{GFlowNet~: exploration de scénarios}
\label{sec:gouv:gflownet}

Le GFlowNet (\emph{Generative Flow Network}) explore l'espace des scénarios macroéconomiques pour découvrir des configurations critiques que l'échantillonnage uniforme manquerait.

\subsection{Deux modes}

\begin{itemize}
  \item \textbf{Mode défensif} (contraction)~: recherche les scénarios qui font franchir un seuil d'ECL, pondérés par leur plausibilité~:
  $$R_{\text{déf}}(\mathbf{z}) = \max(0, \text{ECL}(\mathbf{z}) - \text{seuil}) \times \exp\!\left(-\tfrac{1}{2}\max(0, \|\mathbf{z}\| - \sigma_b)^2\right)$$
  Le terme de plausibilité pénalise les scénarios à plus de $\sigma_b = 4$ déviations standard du centre de la distribution.

  \item \textbf{Mode offensif} (expansion)~: recherche les scénarios maximisant le RORAC sous contrainte d'ECL~:
  $$R_{\text{off}}(\mathbf{z}) = \text{RORAC}(\mathbf{z}) \times \max(0, 1 - \text{ECL}/\text{plafond}) \times \text{plausibilité}(\mathbf{z})$$
\end{itemize}

\subsection{Transformation de Cholesky}

Les scénarios sont générés dans l'espace $Z$ (variables standard indépendantes) puis transformés dans l'espace macro par la décomposition de Cholesky de la matrice de covariance~:
$$\mathbf{X} = \boldsymbol{\mu} + \mathbf{L} \cdot \mathbf{Z}, \qquad \mathbf{L} = \text{Cholesky}(\boldsymbol{\Sigma}_{\text{macro}})$$
Cette transformation garantit que les scénarios générés respectent les corrélations macro (Okun, Taylor) par construction.

\subsection{Apprentissage par Trajectory Balance}

Le GFlowNet est entraîné par la perte \emph{Trajectory Balance} de Malkin et al. (2022)~:
$$\mathcal{L} = \left(\log Z + \log P_F - \log R - \log P_B\right)^2$$
où $P_F$ et $P_B$ sont les probabilités forward et backward le long de la trajectoire, $R$ la récompense, et $Z$ la fonction de partition (apprise). L'architecture est un MLP à deux couches (5 $\to$ 32 $\to$ 16 $\to$ 1), entraîné sur 200 itérations avec un buffer de replay de 500 trajectoires.

\section{Moteurs complémentaires}
\label{sec:gouv:complementaires}

\subsection{Prime de risque de variance (VRP)}

La VRP mesure l'écart entre la volatilité implicite (options VSTOXX) et la volatilité réalisée~:
$$\text{VRP} = \text{IV}^2 - \text{RV}^2$$
Quatre régimes sont définis~: complaisance (VRP $< 0$), normal ($0$--$4\,\%$), stress ($4$--$8\,\%$), panique ($> 8\,\%$). Le régime VRP ajuste le paramètre $\tau$ de Black-Litterman~: en panique ($\tau \times 0{,}40$), l'optimiseur suit les poids de marché~; en complaisance ($\tau \times 1{,}30$), il surpondère les vues fondamentales.

\subsection{Théorie des matrices aléatoires (RMT)}

La matrice de covariance empirique des rendements est bruitée par l'estimation~: avec 14 classes d'actifs et 252 observations quotidiennes, $\sim$25--35\,\% de la variance est du bruit d'estimation. Le nettoyage RMT~:
\begin{enumerate}[nosep]
  \item décompose la matrice en valeurs propres~;
  \item identifie le seuil de Marchenko-Pastur~: $\lambda_{\max} = \sigma^2(1 + \sqrt{N/T})^2$~;
  \item remplace les valeurs propres sous le seuil par leur moyenne (préserve la trace).
\end{enumerate}
Le résultat est une matrice de covariance «~propre~» qui réduit le turnover de l'optimiseur de $\sim$40\,\% \citep{ledoit_wolf2004}.

\subsection{Signatures de chemins}

Les signatures de chemins (Lyons, 1998) compriment l'historique macro de 60 mois en un vecteur de dimension fixe. Pour 5 variables et un ordre de troncature $K = 2$~:
$$\text{dim}(\text{signature}) = \sum_{k=1}^{2} 5^k = 5 + 25 = 30$$
Les 5 composantes d'ordre~1 capturent les tendances~; les 25 composantes d'ordre~2 capturent les interactions temporelles (par exemple, «~le chômage a augmenté \emph{après} que le PIB a baissé~»). Ce vecteur alimente les classifieurs de régime et les modèles de survie.

\subsection{Analyse topologique (TDA)}

L'homologie persistante détecte la fragmentation de la diversification du portefeuille. La matrice de corrélation sectorielle est convertie en matrice de distances ($d_{ij} = 1 - |\rho_{ij}|$), puis analysée par filtration de Vietoris-Rips~:

\begin{equation}
\text{Fragilité} = 0{,}6 \times \bar{|\rho|} + 0{,}4 \times h_1
\label{eq:fragility}
\end{equation}

\noindent où $\bar{|\rho|}$ est la corrélation moyenne absolue et $h_1$ mesure l'absence de structure topologique riche (cycles). Une fragilité $> 0{,}6$ signale une perte de diversification~: tous les secteurs se corrèlent en période de crise.

\subsection{Portes de conformité}

Dix portes de conformité vérifient déterministement chaque sortie du système~:
\begin{itemize}[nosep]
  \item \textbf{CRR3}~: CET1 $\geq 4{,}5\,\%$ (Pilier~1), CET1 $\geq 10{,}5\,\%$ (cible), levier $\geq 3\,\%$, LCR $\geq 100\,\%$~;
  \item \textbf{IFRS~9}~: ratio PMA $\leq 25\,\%$ (overlay de gestion borné, EBA GL/2021/02), couverture $\geq 0{,}1\,\%$, Stage~3 $\leq 10\,\%$~;
  \item \textbf{Appétit pour le risque}~: HHI $\leq 4\,000$~;
  \item \textbf{AI Act}~: explicabilité SHAP + Sobol disponible (art.~13), supervision humaine activée (art.~14).
\end{itemize}

Toute porte en échec \textbf{bloque} les recommandations d'allocation jusqu'à résolution. Ce mécanisme empêche le système de proposer des allocations non conformes, même si elles sont optimales au sens du RAROC.

\section{Synthèse~: dix couches de gouvernance}
\label{sec:gouv:synthese}

Les dix moteurs forment un système de défense en profondeur~:

\begin{table}[htbp]
\centering
\caption{Synthèse des dix moteurs de gouvernance}
\label{tab:gouvernance-synthese}
\small
\begin{tabular}{llll}
\toprule
Moteur & Entrée & Sortie & Rôle \\
\midrule
Conforme & PD prédites & Intervalles $[\underline{\text{PD}}, \overline{\text{PD}}]$ & Staging conservateur \\
Sobol & ECL(macro) & Indices $S^1$, $S^T$ & Décomposition de variance \\
Morris & Paramètres & Ranking OAT & Screening rapide \\
VRP & IV, RV & Régime, $\tau_{\text{BL}}$ & Signal de marché \\
RMT & $\hat{\Sigma}$ & $\Sigma_{\text{propre}}$ & Débruitage covariance \\
Signatures & Historique 60m & Vecteur (30) & Encodage temporel \\
TDA & $R_{\text{corr}}$ & Fragilité $\in [0,1]$ & Alerte diversification \\
Compliance & KPIs & Pass/Warn/Fail & Portes réglementaires \\
HMM & Macro courant & Régime, poids & Classification régime \\
GFlowNet & $R(\mathbf{z})$ & Scénarios critiques & Exploration adversariale \\
\bottomrule
\end{tabular}
\end{table}

Ce chapitre clôt la Partie~I (fondation sectorielle sur 5 secteurs $\times$ 2 canaux). La Partie~II étend ces résultats à 14 classes d'actifs, avec l'optimiseur BL-CVaR et la découverte sur les conventions de signe.


% ══════════════════════════════════════════════════════════
% PARTIE II
% ══════════════════════════════════════════════════════════
\part{Validation multi-actifs}
\label{part:validation}

% ══════════════════════════════════════════════════════════
% CHAPITRE 8 — Extension multi-actifs
% ══════════════════════════════════════════════════════════
\chapter{Extension multi-actifs~: des secteurs aux classes d'actifs}
\label{chap:multi-actifs}

La Partie~I a démontré que l'arbitrage sectoriel (5 secteurs $\times$ 2~canaux) constitue le levier principal d'allocation sous contraintes IFRS~9 et CRR3. Ce résultat, aussi robuste soit-il, repose sur une simplification~: un portefeuille bancaire réel ne se réduit pas au crédit corporate et au private equity. Ce chapitre étend le cadre à \textbf{14 classes d'actifs} couvrant l'intégralité du bilan d'une banque universelle européenne, afin de vérifier si les conclusions sectorielles survivent à l'échelle.

\section{Pourquoi 14 classes~?}
\label{sec:multi-pourquoi}

Le bilan d'une banque universelle de la zone euro comprend typiquement~: du crédit corporate, du private equity, des obligations souveraines, du crédit immobilier, du financement de projet, des \textit{covered bonds}, du crédit à la consommation, du \textit{trade finance}, de l'interbancaire, des produits structurés (titrisation), des actions cotées, des obligations corporate, des repos/SFT et des dérivés. Chacune de ces classes obéit à des règles comptables, prudentielles et économiques distinctes~:

\begin{itemize}
\item \textbf{Trois traitements comptables}~: coût amorti (IFRS~9 §4.1.2), juste valeur par OCI (FVOCI, §4.1.2A) et juste valeur par résultat (FVTPL, §4.1.4). Le traitement détermine si l'ECL s'applique ou non.
\item \textbf{Cinq méthodes de pondération CRR3}~: SA standard (Art.~122), LTV résidentiel (Art.~124--125), SEC-SA titrisation (Art.~242--270), \textit{slotting} financement de projet (Art.~153(5)), ECRA \textit{rating} (Art.~120--122).
\item \textbf{Quatre profils de liquidité HQLA}~: L1 (souverain, 0\% décote), L2A (\textit{covered bonds}, obligations corporate IG, 15\%), L2B (interbancaire, 50\%), non éligible (crédit, PE, dérivés).
\end{itemize}

L'enjeu n'est pas la complétude pour elle-même, mais la \textbf{validité des conclusions}~: si l'arbitrage sectoriel domine dans un univers restreint, domine-t-il encore quand 14 classes se disputent le capital?

\subsection{Taxonomie en trois niveaux}

Les 14 classes sont organisées en trois niveaux de modélisation, selon la granularité des données disponibles~:

\begin{table}[htbp]
\centering
\caption{Taxonomie des 14 classes d'actifs par niveau de modélisation}
\label{tab:taxonomie-14}
\begin{tabular}{llll}
\toprule
\textbf{Niveau} & \textbf{Classe} & \textbf{Poids typique} & \textbf{Spécificité} \\
\midrule
\multirow{4}{*}{L1 — Détaillé} & Prêts corporate & 22\% & DGP v4.5, 3 suites PD \\
 & Private equity & 4\% & NAV 3 canaux \\
 & Actions cotées & 3\% & Merton, FVTPL \\
 & Obligations corporate & 6\% & \textit{Rating}-PD, FVOCI \\
\midrule
\multirow{6}{*}{L2 — Paramétrique} & Souverain & 14\% & HQLA L1, exempt staging \\
 & Crédit immobilier & 15\% & LTV, CRR3 Art.~124 \\
 & Financement de projet & 4\% & SPV, \textit{slotting} \\
 & \textit{Covered bonds} & 7\% & Double recours \\
 & Repos / SFT & 4\% & Collatéralisé, O/N \\
 & Dérivés / CVA & 4\% & SA-CCR, netting \\
\midrule
\multirow{4}{*}{L3 — Simple} & Crédit consommation & 5\% & Lending Club réel \\
 & \textit{Trade finance} & 4\% & Auto-liquidant \\
 & Interbancaire & 3\% & BRRD, ultra-court \\
 & Titrisation & 2\% & SEC-SA, \textit{cliff} \\
\bottomrule
\end{tabular}
\end{table}

Les poids typiques totalisent $\sum w_i = 97\%$ (les 3\% restants couvrent les arrondis et les actifs hors périmètre). Le total des encours du portefeuille est calibré à~:

\begin{equation}
\text{EAD}_{\text{total}} = \frac{\text{RWA}_{\text{budget}}}{\overline{\text{RW}}_{\text{base}}} \approx \frac{3{,}69 \text{T}\texteuro{}}{0{,}54} \approx 6{,}8 \text{T}\texteuro{}
\label{eq:total-ead}
\end{equation}

\noindent où $\text{RWA}_{\text{budget}} = 3{,}69\text{T}\texteuro{}$ est le budget RWA CRR3 et $\overline{\text{RW}}_{\text{base}} = \sum w_i \cdot \text{RW}_i \approx 0{,}54$ la pondération moyenne au portefeuille de base.


\section{Paramétrisation des 14 classes}
\label{sec:multi-params}

Chaque classe d'actifs est décrite par un profil \texttt{AssetClassProfile} comprenant 22 champs, organisés en 6 blocs~: crédit (PD, LGD, ténor, corrélation), réglementaire (RW CRR3, planchers), climat/ESG, sensibilités macro, liquidité/financement et risque de marché. Le tableau~\ref{tab:params-14} synthétise les paramètres clés et leurs justifications.

\begin{table}[htbp]
\centering
\small
\caption{Paramètres clés des 14 classes d'actifs}
\label{tab:params-14}
\begin{tabular}{lcccccc}
\toprule
\textbf{Classe} & \textbf{PD} & \textbf{LGD} & \textbf{RW} & $\rho$ & \textbf{RSF} & \textbf{HQLA} \\
 & (\%) & (\%) & (\%) & & & \\
\midrule
Prêts corporate & 4,50 & 35 & 100 & 0,18 & 0,50 & --- \\
Private equity & 6,00 & 45 & 250 & 0,24 & 1,00 & --- \\
Souverain & 0,05 & 45 & 0 & 0,06 & 0,00 & L1 \\
Crédit immobilier & 1,20 & 15 & 35 & 0,15 & 0,35 & --- \\
Fin.\ de projet & 1,80 & 35 & 130 & 0,20 & 0,85 & --- \\
\textit{Covered bonds} & 0,01 & 10 & 10 & 0,15 & 0,15 & L2A \\
Crédit conso & 3,50 & 65 & 75 & 0,04 & 0,50 & --- \\
\textit{Trade finance} & 0,80 & 30 & 20 & 0,20 & 0,10 & --- \\
Interbancaire & 0,05 & 45 & 20 & 0,24 & 0,00 & L2B \\
Titrisation & 1,00 & 40 & 42\textsuperscript{a} & 0,30 & 0,85 & --- \\
Actions cotées & 2,00 & 85 & 100 & 0,24 & 0,85 & --- \\
Oblig.\ corporate & 1,20 & 45 & 65 & 0,06 & 0,50 & L2A \\
Repos / SFT & 0,05 & 10 & 10 & 0,24 & 0,00 & --- \\
Dérivés / CVA & 0,30 & 35 & 60 & 0,20 & 0,50 & --- \\
\bottomrule
\multicolumn{7}{l}{\footnotesize \textsuperscript{a} RW blended SEC-SA~: 20\%$\times$0,60 + 60\%$\times$0,25 + 100\%$\times$0,15 = 42\%}
\end{tabular}
\end{table}

\subsection{Justification des corrélations d'actif}

La corrélation d'actif $\rho$ détermine l'amplification du risque systématique via la formule de Vasicek (équation~\ref{eq:vasicek-multi}). Les valeurs retenues proviennent de trois sources~:

\begin{enumerate}
\item \textbf{CRR3 réglementaire}~: $\rho = 0{,}24$ pour les expositions corporate non-PME (Art.~153), $\rho = 0{,}15$ pour le résidentiel (Art.~154(2)(a)), $\rho = 0{,}04$ pour le \textit{revolving} retail (Art.~154(2)(b)).
\item \textbf{Empirique (\citealt{longstaff2011})~:} $\rho = 0{,}06$ pour le souverain zone euro (estimé 0,03--0,08 sur CDS 2005--2012). Cette valeur est nettement inférieure au 0,24 réglementaire car le risque souverain en monnaie propre est fondamentalement différent du risque corporate.
\item \textbf{Empirique (\citealt{de-servigny2002})~:} $\rho = 0{,}06$ pour les obligations corporate (mélange 70\% IG / 30\% HY), reflétant la diversification par rating et secteur.
\end{enumerate}

La titrisation utilise $\rho = 0{,}30$ (corrélation implicite élevée due à l'effet de falaise~: un choc modéré n'affecte pas la tranche senior, mais un choc sévère déclenche des pertes brutales).

\subsection{Sensibilités macro et convention de signe}

Chaque classe possède un vecteur de 5 sensibilités macro $(w_{\text{GDP}}, w_{\text{unemp}}, w_{\text{IR}}, w_{\text{HPI}}, w_{\text{infl}})$ qui détermine sa réaction au facteur systématique $Z$. Ces sensibilités varient considérablement~:

\begin{itemize}
\item \textbf{Crédit immobilier}~: $w_{\text{HPI}} = 2{,}5$ (la plus forte sensibilité du portefeuille), $w_{\text{IR}} = 2{,}0$ --- le prix de l'immobilier et le taux sont les déterminants dominants du défaut hypothécaire \citep{mian2009}.
\item \textbf{Crédit consommation}~: $w_{\text{unemp}} = 2{,}0$ --- le chômage est le premier déterminant du défaut retail non collatéralisé (EBA 2023).
\item \textbf{Souverain}~: sensibilités quasi-nulles ($w_i \leq 0{,}20$) --- en zone euro, les obligations souveraines domestiques sont un actif de fuite vers la qualité dont les spreads \emph{compriment} en crise.
\item \textbf{Financement de projet}~: $w_{\text{IR}} = 0{,}0$ --- les SPV à levier fixe ne bénéficient pas directement des baisses de taux (le signal de faiblesse économique compense), conformément à \citet{blanc2014}.
\item \textbf{Obligations corporate}~: $w_{\text{IR}} = 0{,}3$ (contre 1,2 pour le crédit corporate) --- le canal de refinancement favorable domine le canal de service de la dette pour les émetteurs obligataires, conformément aux spreads \citet{gilchrist2012}.
\end{itemize}

Le chapitre~\ref{chap:conventions} analysera en profondeur comment la convention de signe de ces sensibilités (notamment IR et inflation) détermine l'allocation.


\section{Générateurs de positions}
\label{sec:multi-generateurs}

Quatre classes bénéficient d'une modélisation position par position (Niveau~1), avec des générateurs dédiés produisant des DataFrames de positions individuelles. Les dix autres classes utilisent une approche paramétrique basée sur le profil \texttt{AssetClassProfile}.

\subsection{Crédit corporate~: DGP v4.5}

Le générateur corporate (chapitre~\ref{chap:donnees}) produit 30\,000 clients avec un DGP à 3 blocs~: 40\% signal monotone, 35\% interactions, 25\% conjonctions. Ce générateur alimente directement les 3 suites PD (chapitre~\ref{chap:pd}).

\subsection{Actions cotées~: modèle de Merton}

Le générateur actions produit environ 300 positions sur 30 grandes capitalisations européennes (EURO STOXX 50/600). La PD est estimée par le modèle de \citet{merton1974}~:

\begin{equation}
\text{DD} = \frac{\ln(V/D) + (\mu - \sigma_A^2/2) \cdot T}{\sigma_A \sqrt{T}}, \quad
\text{PD} = \exp(-1{,}60 - 0{,}231 \times \text{DD})
\label{eq:merton-dd}
\end{equation}

\noindent où $V$ est la valeur des actifs, $D$ la dette, $\sigma_A$ la volatilité des actifs (dé-levragée via l'approximation d'Itô) et DD la distance au défaut. Le mapping exponentiel DD$\to$PD suit la calibration KMV \citep{crosbie2003}. La PD est bornée entre 3~bp (plancher CRR3) et 10\%.

Les actions sont comptabilisées en FVTPL (IFRS~9 §4.1.4)~: pas d'ECL, la perte est une perte de valeur \textit{mark-to-market}. La volatilité de marché est calibrée à 20\% (VSTOXX historique, \citealt{bekaert2014}). Le stress s'exerce via 4 canaux~: choc indiciel ($\beta$), sensibilité taux (duration 15~ans), rotation sectorielle et expansion de volatilité GJR.

\subsection{Obligations corporate~: PD par rating}

Le générateur obligations produit environ 500 positions avec un mélange 70\% \textit{investment grade} / 30\% \textit{high yield}. La PD est dérivée des matrices de transition de \citet{moody2023}~:

\begin{table}[htbp]
\centering
\caption{PD à 1 an par rating (Moody's 1983--2022)}
\label{tab:rating-pd}
\begin{tabular}{lcccccc}
\toprule
\textbf{Rating} & AAA & AA & A & BBB & BB & B--CCC \\
\midrule
PD (\%) & 0,01 & 0,02 & 0,05 & 0,20 & 1,00 & 5,00 \\
RW CRR3 ECRA & 20\% & 20\% & 50\% & 75\% & 100\% & 150\% \\
\bottomrule
\end{tabular}
\end{table}

La LGD est fonction de la séniorité~: secured (35\%), senior unsecured (45\%), subordonnée (70\%), conformément aux taux de recouvrement historiques \citep{altman2006}. Le traitement comptable est FVOCI~: l'ECL est calculée comme pour le coût amorti, mais les variations de juste valeur transitent par les OCI.

\subsection{Repos et SFT~: calibration ICMA}

Le générateur repos produit environ 400 positions sur 25 contreparties bancaires européennes. La calibration suit l'ICMA European Repo Survey (décembre 2024)~:

\begin{itemize}
\item \textbf{Collatéral}~: 80\% obligations souveraines, 10\% corporate IG, 7\% actions, 3\% autre.
\item \textbf{Ténors}~: 35\% O/N, 20\% 1 semaine, 25\% 1 mois, 15\% 3 mois, 5\% $>$ 3 mois.
\item \textbf{Décotes} (CRR3 Art.~224)~: 2\% souverain, 5\% corporate, 15\% actions, 10\% autre.
\end{itemize}

La LGD est très faible (moyenne $\sim$3--12\%) grâce au collatéral. Le stress suit le modèle de spirale de décotes de \citet{brunnermeier2009}~:

\begin{equation}
H_{\text{stress}} = H_{\text{base}} \times \left(1 + \gamma \cdot \max(0, -\Delta_{\text{GDP}})\right)
\label{eq:haircut-spiral}
\end{equation}

\noindent où $\gamma$ est l'élasticité de la spirale~: en récession ($\Delta_{\text{GDP}} < 0$), les décotes augmentent, forçant des appels de marge qui amplifient la contagion.

\subsection{Dérivés~: SA-CCR et CVA}

Le générateur dérivés produit environ 500 positions réparties sur 3 desks~: IRS ($\sim$75\%), FX ($\sim$15\%), CDS ($\sim$10\%). L'EAD est calculée selon l'approche standard SA-CCR (CRR3 Art.~274--280)~:

\begin{equation}
\text{EAD} = \alpha \times (\text{RC} + \text{PFE}), \quad \alpha = 1{,}4
\label{eq:sa-ccr}
\end{equation}

\noindent où RC est le coût de remplacement et PFE l'exposition future potentielle. L'EAD représente typiquement 15--30\% du notionnel, grâce aux conventions de \textit{netting} (ISDA Master) et aux accords de collatéral (CSA). La pondération RW utilise l'approche SA-CVA (CRR3 Art.~382--386), avec un RW moyen de 60\%.


\section{Bilan multi-actifs et agrégation ECL}
\label{sec:multi-bilan}

L'extension de 5 secteurs à 14 classes nécessite un mécanisme d'agrégation unifié. Chaque classe produit une ligne de bilan contenant 19 champs standardisés~: PD, LGD, EAD, RW, ECL, NII, et les métriques sectorielles. Le module \texttt{balance\_sheet\_ecl.py} orchestre la transformation macro$\to$ECL en trois étapes.

\subsection{Transformation macro $\to$ facteur systématique}

Le facteur systématique $Z$ agrège les 5 variables macro pondérées par les sensibilités de la classe~:

\begin{equation}
Z = \sum_{i=1}^{5} s_i \cdot w_i \cdot \frac{x_i - \mu_i}{\sigma_i}
\label{eq:macro-to-z}
\end{equation}

\noindent où $x_i$ est la valeur courante de la variable macro $i$, $\mu_i$ et $\sigma_i$ les paramètres du scénario de base (tableau~\ref{tab:scenario-base-z}), $w_i$ la sensibilité de la classe, et $s_i$ le signe conventionnel~:

\begin{equation}
s_i = \begin{cases}
+1 & \text{si } i \in \{\text{chômage}, \text{taux d'intérêt}\} \quad \text{(hausse = adverse)} \\
-1 & \text{si } i \in \{\text{PIB}, \text{HPI}, \text{inflation}\} \quad \text{(hausse = favorable)}
\end{cases}
\label{eq:signe-z}
\end{equation}

\begin{table}[htbp]
\centering
\caption{Paramètres de référence du scénario de base}
\label{tab:scenario-base-z}
\begin{tabular}{lccc}
\toprule
\textbf{Variable} & $\mu$ (base) & $\sigma$ (échelle) & \textbf{Convention} \\
\midrule
PIB (\%) & 1,2 & 1,8 & Favorable ($s=-1$) \\
Chômage (\%) & 7,5 & 1,5 & Adverse ($s=+1$) \\
Taux d'intérêt (\%) & 3,5 & 1,0 & Adverse ($s=+1$) \\
Prix immobiliers (\%) & 2,0 & 3,0 & Favorable ($s=-1$) \\
Inflation (\%) & 2,5 & 1,2 & Favorable ($s=-1$) \\
\bottomrule
\end{tabular}
\end{table}

La convention $s_{\text{IR}} = +1$ (taux adverse) suit \citet{duffie2007}~: le canal du service de la dette (\textit{debt service burden}) domine empiriquement le canal de Merton (\textit{risk-free drift}). Cette convention a été validée sur 96 exécutions du pipeline (8 conventions $\times$ 12 scénarios, cf.\ chapitre~\ref{chap:conventions}).

\subsection{PD conditionnelle de Vasicek}

Le facteur $Z$ est injecté dans la formule ASRF de \citet{vasicek2002}, transposée au CRR3 (Art.~153)~:

\begin{equation}
\text{PD}_{\text{cond}}(Z) = \Phi\!\left(\frac{\Phi^{-1}(\text{PD}_{\text{TTC}}) + \sqrt{\rho} \cdot Z}{\sqrt{1 - \rho}}\right)
\label{eq:vasicek-multi}
\end{equation}

\noindent où $\Phi$ est la fonction de répartition normale standard. Pour un choc adverse $Z = +2$ (97,7\textsuperscript{e} percentile)~:

\begin{itemize}
\item \textbf{Corporate} ($\text{PD}=4{,}5\%$, $\rho=0{,}18$)~: $\text{PD}_{\text{cond}} \approx 18\%$ ($\times 4{,}0$).
\item \textbf{Souverain} ($\text{PD}=0{,}05\%$, $\rho=0{,}06$)~: $\text{PD}_{\text{cond}} \approx 0{,}4\%$ ($\times 8{,}0$, mais depuis une base très faible).
\item \textbf{Consommation} ($\text{PD}=3{,}5\%$, $\rho=0{,}04$)~: $\text{PD}_{\text{cond}} \approx 6{,}8\%$ ($\times 1{,}9$, faible corrélation $\to$ moins amplifié).
\end{itemize}

L'asymétrie est structurelle~: les classes à forte corrélation d'actif ($\rho \geq 0{,}20$~: PE, interbancaire, repos) subissent une amplification plus violente, ce qui pénalise leur consommation de capital en stress.

\subsection{Staging et ECL par classe}

L'ECL combine le staging IFRS~9 (Stage~1 vs Stage~2) avec la PD conditionnelle~:

\begin{equation}
\text{ECL} = (1 - p_{S2}) \cdot \underbrace{\text{PD}_{\text{cond}} \cdot \text{LGD} \cdot \text{EAD}}_{\text{Stage 1 (12 mois)}} + p_{S2} \cdot \underbrace{[1-(1-\text{PD}_{\text{cond}})^T] \cdot \text{LGD} \cdot \text{EAD}}_{\text{Stage 2 (lifetime)}}
\label{eq:ecl-staging-multi}
\end{equation}

\noindent avec une transition logistique douce~:

\begin{equation}
p_{S2} = \frac{1}{1 + \exp\!\left(-5 \cdot \left(\frac{\text{PD}_{\text{cond}}}{\text{PD}_{\text{orig}} \cdot \theta_{\text{SICR}}} - 1\right)\right)}
\label{eq:transition-staging}
\end{equation}

Deux exemptions sont appliquées~:
\begin{itemize}
\item \textbf{Souverain}~: \texttt{exempt\_from\_staging = True} (IFRS~9 B5.5.25) --- ECL 12 mois uniquement, pas de Stage~2.
\item \textbf{Actions FVTPL}~: \texttt{exempt\_from\_staging = True} (IFRS~9 §4.1.4) --- pas d'ECL du tout, la perte est une perte de valeur \textit{mark-to-market}.
\end{itemize}


\section{Décorrélation du facteur systématique}
\label{sec:multi-decorrelation}

La formule naïve de l'équation~\ref{eq:macro-to-z} présente un problème de double comptage~: les 5 variables macro sont corrélées entre elles. La loi d'Okun lie PIB et chômage ($r \approx -0{,}70$), la règle de Taylor lie taux et inflation ($r \approx +0{,}60$). Sans correction, un choc de stagflation (PIB$\downarrow$, chômage$\uparrow$, taux$\uparrow$, inflation$\uparrow$) est compté 3--4 fois au lieu d'une.

\subsection{Correction par la matrice de corrélation signée}

La correction utilise la matrice de corrélation signée $R_s$, construite à partir de la matrice de covariance macro \texttt{MACRO\_COVARIANCE}~:

\begin{equation}
Z_{\text{corrigé}} = Z_{\text{brut}} \times \sqrt{\frac{\mathbf{w}^T \mathbf{w}}{\mathbf{w}^T R_s \mathbf{w}}}
\label{eq:decorrelation}
\end{equation}

\noindent où $\mathbf{w}$ est le vecteur des sensibilités pondérées par le signe ($w_i \cdot s_i$) et $R_s$ la matrice de corrélation signée (les corrélations sont inversées pour les paires de signes opposés). Le facteur $\sqrt{\mathbf{w}^T \mathbf{w} / \mathbf{w}^T R_s \mathbf{w}}$ est toujours $\leq 1$ quand les variables sont positivement co-mobiles en termes de stress, ce qui réduit l'amplitude du $Z$.

\subsection{Impact quantitatif}

La décorrélation a un effet significatif sur les scénarios de crise où plusieurs variables se déplacent simultanément~:

\begin{itemize}
\item \textbf{Corporate}~: EL réduite de 16--25\% dans les scénarios affectés (Stagflation, GFC, COVID).
\item \textbf{Crise souveraine}~: l'écart de performance ajustée au risque (APR) passe de 4,05\% à 1,37\% après décorrélation.
\item \textbf{Convention de signe}~: le résultat pré-décorrélation ($A = 96\%$, $B = 93\%$, dilemme) était un artefact du double comptage. Après correction, $B = 97{,}1\%$ et $A$ est éliminée.
\end{itemize}

Ce résultat illustre un principe général~: \emph{la correction de la multicolinéarité est un pré-requis avant toute analyse de convention de signe} (chapitre~\ref{chap:conventions}).


\section{Capacité de marché et impact de prix}
\label{sec:multi-market-cap}

L'optimiseur BL-CVaR (chapitre~\ref{chap:blcvar}) intègre un mécanisme d'impact de prix logarithmique (\citealt{kyle1985}, \citealt{almgren2001}) qui pénalise la concentration au-delà de la capacité d'absorption du marché. La capacité de chaque classe est calibrée sur des sources institutionnelles~:

\begin{table}[htbp]
\centering
\caption{Capacité de marché par classe d'actifs (zone euro)}
\label{tab:market-capacity}
\begin{tabular}{lrl}
\toprule
\textbf{Classe} & \textbf{Capacité (Mds~\texteuro{})} & \textbf{Source} \\
\midrule
Souverain & 11\,000 & ECB SDW \\
Repos / SFT & 8\,000 & ICMA Survey \#48 \\
Prêts corporate & 5\,000 & ECB MFI Statistics \\
Crédit immobilier & 5\,000 & ECB BSI \\
Obligations corporate & 3\,000 & EU IG+HY (Dealogic) \\
\textit{Covered bonds} & 2\,500 & ECBC Factbook \\
Interbancaire & 2\,000 & ECB Money Market \\
Actions cotées & 2\,000 & BCE Holdings Survey \\
Crédit consommation & 1\,500 & ECB BSI \\
Titrisation & 1\,500 & EU ABS (AFME) \\
\textit{Trade finance} & 1\,000 & ICC Trade Register \\
Dérivés / CVA & 1\,000 & BIS OTC (base EAD) \\
Financement de projet & 500 & Dealogic PF \\
Private equity & 300 & Preqin EU \\
\bottomrule
\end{tabular}
\end{table}

La compression de spread résultante est~:

\begin{equation}
\mu_{\text{eff},i} = \frac{\mu_{\text{base},i}}{1 + \text{compression}_i}, \quad
\text{compression}_i = -\ln\!\left(1 - \frac{w_i \cdot \text{EAD}_{\text{total}}}{\text{Capacité}_i}\right)
\label{eq:spread-compression}
\end{equation}

Pour le private equity, une allocation de 4\% ($\approx 272\text{M}\texteuro{}$) représente une part de marché de $272/300 \approx 91\%$, ce qui génère une compression de $-\ln(0{,}09) \approx 2{,}41$, soit une réduction du rendement effectif de $\sim$71\%. C'est le mécanisme endogène principal qui empêche la sur-allocation au PE~: au-delà de $\sim$4\%, le rendement net s'effondre.


\section{Résultats de la validation multi-actifs}
\label{sec:multi-resultats}

La validation économique du portefeuille 14 classes repose sur 246 tests automatisés, dont 21 tests de cohérence économique inter-scénarios.

\subsection{Cohérence RAROC}

Le RAROC est calculé de manière unifiée pour les 14 classes, avec des marges spécifiques par classe~:

\begin{table}[htbp]
\centering
\caption{Marges commerciales et ratio coût/revenu par classe}
\label{tab:margins-14}
\begin{tabular}{lcc}
\toprule
\textbf{Classe} & \textbf{Marge (bps)} & \textbf{CIR} \\
\midrule
Crédit consommation & 350 & 58\% \\
Financement de projet & 220 & 55\% \\
Titrisation & 150 & 50\% \\
Crédit immobilier & 130 & 42\% \\
\textit{Trade finance} & 90 & 30\% \\
Obligations corporate & 80 & 20\% \\
\textit{Covered bonds} & 40 & 15\% \\
Dérivés / CVA & 30 & 45\% \\
Interbancaire & 15 & 12\% \\
Souverain & 10 & 8\% \\
Repos / SFT & 5 & 8\% \\
\bottomrule
\end{tabular}
\end{table}

Les marges reflètent la tarification de chaque marché~: le crédit consommation à 350~bps (NIM élevé, non collatéralisé), contre 5~bps pour les repos (spread de financement minimal, volume pur). Les CIR sont cohérents avec les benchmarks EBA~: le retail non collatéralisé est le plus coûteux (58\%), l'interbancaire le moins (12\%).

\subsection{Tests de cohérence économique}

Le validateur économique vérifie 50 règles de Spearman (7 \textsc{hard} + 43 \textsc{soft}) et 11 comparaisons par paires sur les 12 scénarios~:

\begin{itemize}
\item \textbf{Pro-cyclicité du PE et des actions}~: l'allocation PE doit être corrélée positivement au PIB et négativement au chômage (Spearman, $p < 0{,}05$). Vérifié.
\item \textbf{Crédit immobilier $\times$ HPI}~: corrélation positive significative entre allocation mortgage et prix immobiliers. Vérifié.
\item \textbf{Ordonnancement des crises}~: PE Reprise $>$ PE GFC, PE Hypercroissance $>$ PE Stagflation, Equities Reprise $>$ Equities Stagflation. Vérifié.
\item \textbf{RAROC croissance $>$ RAROC crise}~: le RAROC global est plus élevé dans les scénarios favorables. Vérifié.
\end{itemize}

Résultat~: \textbf{246 tests PASS, 0 FAIL, 27 avertissements SOFT} (structurels, documentés). Les avertissements concernent principalement des effets de second ordre~: le PE en Hypercroissance est pénalisé par un taux d'intérêt de 8\% (\textit{debt service burden} des LBO), ce qui réduit son allocation malgré un PIB favorable.

\subsection{Confirmation du résultat sectoriel}

Le résultat central de la Partie~I se confirme à l'échelle 14 classes~: \textbf{l'allocation sectorielle reste le déterminant principal du RAROC}. Les classes à forte consommation de capital réglementaire (PE à 250\%, financement de projet à 130\%, titrisation à 42\% blended) sont structurellement contraintes par le ratio CET1, tandis que les classes à faible RW (souverain 0\%, repos 10\%, \textit{covered bonds} 10\%) sont sur-pondérées par les contraintes LCR et NSFR. L'optimiseur BL-CVaR (chapitre~\ref{chap:blcvar}) formalise cet arbitrage.


% ══════════════════════════════════════════════════════════
% CHAPITRE 9 — Optimiseur BL-CVaR
% ══════════════════════════════════════════════════════════
\chapter{Optimiseur BL-CVaR~: allocation multi-contraintes}
\label{chap:blcvar}

Le chapitre précédent a établi le cadre comptable et prudentiel des 14 classes d'actifs. Ce chapitre présente l'optimiseur qui détermine l'allocation optimale sous contraintes multiples~: CET1, LCR, NSFR, IRRBB et capacité de marché. L'architecture suit le principe \textbf{zéro paramètre arbitraire}~: chaque constante est soit dérivée du portefeuille de base, soit calibrée sur une source réglementaire ou de marché.

\section{Architecture en deux phases}
\label{sec:blcvar-archi}

L'optimiseur fonctionne en deux phases séquentielles~:

\begin{enumerate}
\item \textbf{Phase~1 --- BL-CVaR pur}~: maximisation du rendement ajusté au risque, sans contrainte réglementaire. Produit une allocation économiquement optimale.
\item \textbf{Phase~2 --- Ajustement réglementaire}~: projection séquentielle de l'allocation Phase~1 sur les contraintes CET1, LCR, NSFR et IRRBB. Chaque étape modifie minimalement l'allocation pour satisfaire la norme.
\end{enumerate}

Cette séparation est délibérée~: elle permet d'identifier la contribution de chaque contrainte réglementaire à l'écart entre l'allocation économique et l'allocation réglementaire.

\section{Phase~1~: BL-CVaR avec impact de prix}
\label{sec:blcvar-phase1}

\subsection{Objectif}

L'objectif de la Phase~1 est~:

\begin{equation}
\mathbf{w}^* = \arg\max_{\mathbf{w} \in \Delta^{N-1}} \left[ \mu_{\text{eff}}(\mathbf{w}) - \kappa \cdot \text{CVaR}_\alpha(\mathbf{w}) \right]
\label{eq:blcvar-obj}
\end{equation}

\noindent où $\Delta^{N-1}$ est le simplexe ($w_i \geq 0$, $\sum w_i = 1$), $\mu_{\text{eff}}$ le rendement effectif après compression de spread (équation~\ref{eq:spread-compression}), et $\text{CVaR}_\alpha$ la \textit{Conditional Value-at-Risk} à 95\%.

\subsection{Estimation Monte Carlo de la CVaR}

La CVaR est estimée par simulation de 5\,000 scénarios centrés~:

\begin{equation}
\text{CVaR}_\alpha(\mathbf{w}) = -\frac{1}{|\mathcal{T}|} \sum_{s \in \mathcal{T}} \mathbf{r}_s^T \mathbf{w}
\label{eq:cvar-mc}
\end{equation}

\noindent où $\mathcal{T} = \{s : \mathbf{r}_s^T \mathbf{w} \leq \text{VaR}_\alpha\}$ est l'ensemble des scénarios dans la queue de 5\% (\citealt{rockafellar2000}). Les scénarios sont générés à partir de la matrice de corrélation $14 \times 14$ construite par \texttt{\_build\_corr\_matrix()}, incluant les corrélations de fuite vers la qualité (souverain--actions $r = -0{,}35$, souverain--obligations corporate $r = -0{,}30$).

\subsection{Score BL-CVaR et softmax}

Le gradient de CVaR par classe est~:

\begin{equation}
\frac{\partial \text{CVaR}}{\partial w_i} \approx -\frac{1}{|\mathcal{T}|} \sum_{s \in \mathcal{T}} r_{s,i}
\label{eq:cvar-gradient}
\end{equation}

Le score BL-CVaR combine rendement et risque~:

\begin{equation}
\text{score}_i = \mu_{\text{eff},i} - \kappa \cdot \frac{\partial \text{CVaR}}{\partial w_i}
\label{eq:blcvar-score}
\end{equation}

L'allocation est obtenue par softmax avec température auto-calibrée~:

\begin{equation}
w_i^{\text{opt}} = \frac{\exp(\text{score}_i / \tau)}{\sum_j \exp(\text{score}_j / \tau)}, \quad
\tau = \frac{\ln(R_{\max})}{\text{score}_{\max} - \text{score}_{\min}}
\label{eq:softmax-blcvar}
\end{equation}

\noindent où $R_{\max} = 20$ borne le ratio de poids maximal entre la meilleure et la pire classe. Le poids minimum est $1/N^2$ ($N = 14 \Rightarrow w_{\min} = 0{,}51\%$).

\subsection{Kappa pro-cyclique}
\label{sec:kappa}

Le paramètre de pénalité $\kappa$ n'est pas fixé a priori mais dérivé du portefeuille de base~:

\begin{equation}
\kappa_{\text{base}} = \text{clip}\!\left(\frac{\text{CVaR}_{\text{base}}}{|\mu_{p,\text{base}}|}, \; 0{,}2, \; 2{,}0\right)
\label{eq:kappa}
\end{equation}

Ce choix rend $\kappa$ \textbf{pro-cyclique}~: en expansion ($\mu_p$ élevé), $\kappa$ baisse et l'optimiseur tolère plus de risque~; en crise ($\mu_p$ faible), $\kappa$ augmente et l'optimiseur se défend. Le clipping à $[0{,}2 ; 2{,}0]$ empêche les comportements extrêmes. En pratique, $\kappa$ varie de 0,61 (Reprise) à 2,00 (toutes les crises, borné).

\subsection{Volatilité cyclique}

Pour les classes PE et actions, la volatilité des scénarios Monte Carlo est modulée par le cycle macro~:

\begin{equation}
v_i = v_{\text{base},i} \times \frac{1 + 0{,}5 \cdot \text{stress}}{1 + 0{,}3 \cdot \text{prospérité}}
\label{eq:vol-cyclique}
\end{equation}

\noindent où $\text{stress} = \max(-\text{PIB}/0{,}02, 0)$ et $\text{prospérité} = \max(\text{PIB}/0{,}02 - 1, 0)$. L'asymétrie 0,5/0,3 reflète l'effet de levier de \citet{black1976}~: la volatilité augmente plus vite en baisse qu'elle ne diminue en hausse. En crise, le gradient CVaR du PE double, réduisant mécaniquement son allocation à $\sim$0,5\%.

\subsection{Convergence par accélération d'Anderson}

La compression de spread crée une boucle de rétroaction~: les poids modifient les rendements, qui modifient les poids. Le point fixe est trouvé par accélération d'Anderson (\citealt{walker2011}) avec profondeur de mémoire $m = 3$~:

\begin{equation}
F(\mathbf{w}) = \text{softmax}\!\left(\frac{\mu_{\text{eff}}(\mathbf{w}) - \kappa \cdot \nabla \text{CVaR}(\mathbf{w})}{\tau}\right)
\label{eq:anderson}
\end{equation}

L'algorithme converge typiquement en 4--6 itérations ($\|\text{résidu}\| < 10^{-6}$). En cas de non-convergence, un \textit{fallback} vers la méthode à 2 passes (itérations de Picard) est activé.


\section{Phase~2~: contraintes réglementaires}
\label{sec:blcvar-phase2}

La Phase~2 projette séquentiellement l'allocation Phase~1 sur quatre contraintes réglementaires. L'ordre est important~: CET1 (la plus contraignante) en premier, puis LCR, NSFR et IRRBB.

\subsection{CET1~: projection par gradient RW}
\label{sec:cet1-proj}

La contrainte CET1 impose~:

\begin{equation}
\frac{K_{\text{CET1}} - \text{ECL}_{\text{total}}}{\sum_i w_i \cdot \text{EAD}_i \cdot \text{RW}_i} \geq \theta_{\text{CET1}} = 13\%
\label{eq:cet1-constraint}
\end{equation}

\noindent où $K_{\text{CET1}} = \text{RWA}_{\text{budget}} \times \theta_{\text{CET1}} = 3{,}69\text{T} \times 13\% = 479{,}7\text{Mds}\texteuro{}$. Plutôt qu'une interpolation vers le portefeuille de base (qui écraserait les différences inter-scénarios), la Phase~2 utilise une \textbf{projection multiplicative par gradient RW}~:

\begin{equation}
w_i^{\text{CET1}} = w_i^{\text{Phase1}} \times \left(1 - \lambda \cdot (\text{RW}_i - \overline{\text{RW}})\right)
\label{eq:cet1-gradient}
\end{equation}

\noindent où $\lambda$ est trouvé par recherche binaire pour satisfaire l'équation~\ref{eq:cet1-constraint}. Les classes à fort RW (PE~250\%, projet~130\%) sont réduites, les classes à faible RW (souverain~0\%, repos~10\%) sont augmentées, proportionnellement à l'écart au RW moyen.

En pratique, CET1 est la contrainte la plus active~: elle lie dans \textbf{10 scénarios sur 11} (seule la Trappe à liquidité passe sans ajustement). La marge résiduelle est de +0,10\% en crise (très serrée) et +0,78--1,52\% en expansion.

\subsection{LCR~: ratio de liquidité à court terme}
\label{sec:lcr}

Le LCR (Bâle~III, CRR Art.~412) impose~:

\begin{equation}
\text{LCR} = \frac{\text{HQLA}}{\text{TNCO}} \geq 100\%
\label{eq:lcr}
\end{equation}

\noindent avec~:

\begin{equation}
\text{HQLA} = \sum_i w_i \cdot \text{EAD}_i \cdot \mathbb{1}_{\text{HQLA},i} \cdot (1 - H_i)
\label{eq:hqla}
\end{equation}

\noindent où $H_i$ est la décote LCR~: 0\% pour le L1 (souverain), 15\% pour le L2A (\textit{covered bonds}, obligations corporate IG), 50\% pour le L2B (interbancaire). Le TNCO est simplifié à 15\% de l'EAD total. Le L2 est plafonné à 40\% du HQLA total (Bâle~III §50).

L'ajustement est adaptatif~: le pas de transfert est $\min(|\text{déficit}| \times 0{,}10, \; 0{,}02)$. En pratique, le LCR ne lie que dans 1 scénario sur 11 (Hypercroissance, où la forte croissance réduit les classes HQLA au profit d'actifs plus rentables).

\subsection{NSFR~: financement stable}
\label{sec:nsfr}

Le NSFR (Bâle~III, CRR Art.~428) impose~:

\begin{equation}
\text{NSFR} = \frac{\text{ASF}}{\text{RSF}} \geq 100\%, \quad
\text{RSF} = \sum_i w_i \cdot \text{EAD}_i \cdot \text{RSF}_i
\label{eq:nsfr}
\end{equation}

\noindent où $\text{RSF}_i$ est le poids de financement stable requis~: 0\% pour le souverain et l'interbancaire, 10\% pour le \textit{trade finance} (auto-liquidant), 35\% pour le crédit immobilier, 85\% pour la titrisation et le financement de projet, 100\% pour le PE (illiquide). Le financement stable disponible (ASF) est simplifié à 90\% de l'EAD total.

L'ajustement transfère des classes à haut RSF vers les classes à bas RSF, proportionnellement à la capacité résiduelle de chaque classe. Le NSFR ne lie dans aucun des 11 scénarios testés.

\subsection{IRRBB~: risque de taux}
\label{sec:irrbb}

L'IRRBB (EBA GL/2022/14) contraint la variation de valeur économique (EVE) sous un choc de taux de 200~bps~:

\begin{equation}
\frac{\Delta \text{EVE}}{\text{CET1}} = \frac{\sum_i w_i \cdot \text{EAD}_i \cdot D_i^{\text{eff}} \cdot \Delta r}{\text{CET1}} \leq 15\%
\label{eq:irrbb}
\end{equation}

\noindent où $D_i^{\text{eff}} = D_i \times (1 - h_{\text{ALM}})$ est la duration effective après couverture ALM ($h_{\text{ALM}} = 0{,}85$, cf.\ chapitre~\ref{chap:metriques}). Le ratio EVE au portefeuille de base est d'environ 91\% de la limite --- serré mais conforme. L'IRRBB ne lie dans aucun scénario, grâce à la couverture ALM de 85\%.


\section{Zéro paramètre arbitraire}
\label{sec:blcvar-zero-params}

Le principe fondateur de l'optimiseur est l'élimination de tout paramètre discrétionnaire. Le tableau~\ref{tab:zero-params} recense les 22 constantes supprimées entre les versions 1 et 3 de l'optimiseur~:

\begin{table}[htbp]
\centering
\small
\caption{Paramètres éliminés par auto-calibration ou suppression}
\label{tab:zero-params}
\begin{tabular}{lll}
\toprule
\textbf{Paramètre} & \textbf{Ancien} & \textbf{Remplacé par} \\
\midrule
$\lambda_{\text{HHI}}$ & 0,15 (grid search) & Supprimé \\
$\lambda_{\text{depth}}$ & 1,0 & Compression log (eq.~\ref{eq:spread-compression}) \\
\texttt{depth\_factor} & Par classe & \texttt{market\_capacity\_eur} (BCE/BRI) \\
$\kappa$ & 1,0 fixe & Pro-cyclique (eq.~\ref{eq:kappa}) \\
$\tau$ (température) & 2,0 fixe & $\ln(20)/\text{score\_range}$ \\
$w_{\min}$ & 1\% & $1/N^2 = 0{,}51\%$ \\
PE band $[a, b]$ & [3\%, 15\%] & Émergent (market impact) \\
L2 $\lambda_{\text{reg}}$ & 0,01 & Diagnostique uniquement \\
\bottomrule
\end{tabular}
\end{table}

La philosophie sous-jacente est que tout paramètre arbitraire introduit un degré de liberté exploitable par l'opérateur~: en ajustant $\lambda_{\text{HHI}}$ ou le PE band, on peut obtenir n'importe quelle allocation souhaitée. En le supprimant, on force l'allocation à émerger des données de marché (capacités), des normes (CET1, LCR, NSFR) et du profil risque-rendement (CVaR, spread). Le résultat n'est plus une préférence de l'opérateur mais une conséquence de la structure du marché.


\section{Score de qualité d'allocation (AQS)}
\label{sec:aqs}

Pour évaluer les allocations sans recourir à l'optimiseur, nous définissons un \textbf{score de qualité d'allocation} (AQS) inspiré de la désirabilité composite de \citet{derringer1980}~:

\begin{equation}
\text{AQS} = \left(\prod_{k=1}^{9} d_k\right)^{1/9}
\label{eq:aqs}
\end{equation}

\noindent où chaque $d_k \in [0, 1]$ est une désirabilité individuelle~:

\begin{itemize}
\item \textbf{4 désirabilités binaires} (réglementaires)~: CET1 $\geq$ 11,2\%, LCR $\geq$ 100\%, NSFR $\geq$ 100\%, IRRBB conforme. Si l'une est violée, $d_k = 0$ et AQS $= 0$ (la moyenne géométrique annule le score).
\item \textbf{5 désirabilités relatives} (économiques)~: RAROC $\in [-10\%, 20\%]$, CVaR $\in [0{,}1\%, 20\%]$ (inverse), taux de profit $\in [-0{,}5\%, 1{,}0\%]$, HHI $\in [500, 3000]$ (inverse), Sharpe $\in [-5, 5]$.
\end{itemize}

Les bornes de normalisation sont calibrées pour que les allocations raisonnables obtiennent un AQS entre 0,80 et 1,00. Résultat~: sous la convention~B, l'AQS moyen sur les 12 scénarios est de \textbf{0,9404}, avec 12/12 scénarios à AQS $> 0$ (aucune violation réglementaire).


\section{Analyse de sensibilité paramétrique}
\label{sec:param-sensitivity}

Pour vérifier que les résultats ne dépendent pas d'un paramètre spécifique, nous réalisons une analyse de sensibilité de Morris et de Sobol sur les 28 paramètres de l'optimiseur (14 volatilités de marché + 14 poids typiques, perturbés $\pm 50\%$).

\subsection{Criblage de Morris}

L'analyse de Morris (\citealt{morris1991}) évalue l'effet élémentaire de chaque paramètre via $r = 20$ trajectoires. Les indicateurs $\mu^*$ (importance) et $\sigma$ (interactions/non-linéarité) identifient les paramètres influents~:

\begin{itemize}
\item \textbf{Volatilité PE} ($\mu^* = 0{,}047$)~: le paramètre le plus influent --- toute augmentation de la volatilité PE réduit significativement son allocation (via le gradient CVaR).
\item \textbf{Volatilité actions} ($\mu^* = 0{,}031$)~: deuxième paramètre, même mécanisme.
\item \textbf{Poids typique souverain} ($\mu^* = 0{,}028$)~: influence via la contrainte CET1 (RW = 0\% réduit le RWA moyen).
\end{itemize}

\subsection{Indices de Sobol}

L'analyse de Sobol (\citealt{saltelli2002}) décompose la variance de l'allocation en effets principaux ($S_1$) et effets totaux ($S_T$). Le schéma de \citet{saltelli2002} nécessite $N \times (D + 2) = 128 \times 30 = 3\,840$ évaluations de l'optimiseur. Les résultats confirment que la volatilité PE et actions concentrent $>$60\% de la variance totale de l'allocation, tandis que les poids typiques ont un effet de second ordre ($S_T < 0{,}15$).

Ce résultat valide le principe de zéro paramètre arbitraire~: les paramètres les plus influents (volatilités) sont des observables de marché (VSTOXX, Cambridge PE Index), pas des choix de l'opérateur.


% ══════════════════════════════════════════════════════════
% CHAPITRE 10 — Risque de spécification du facteur systématique
% ══════════════════════════════════════════════════════════
\chapter{Risque de spécification du facteur systématique}
\label{chap:conventions}

Ce chapitre montre que \textbf{l'allocation est plus sensible à la spécification du modèle macro-facteur qu'aux paramètres de l'optimiseur}. La convention de signe --- la forme la plus élémentaire de cette spécification --- produit un swing de 17,8~points de pourcentage sur le RAROC, contre 2--5~pp pour l'optimisation BL-CVaR. Ce résultat suggère que le risque de modèle dominant dans l'allocation de capital n'est pas l'optimisation mais la modélisation macro.

\section{Origine du problème}
\label{sec:conv-origine}

La transformation macro$\to Z$ (équation~\ref{eq:macro-to-z}) requiert un signe conventionnel $s_i$ pour chacune des 5 variables. Deux variables ont un signe non ambigu~:

\begin{itemize}
\item \textbf{PIB}~: une hausse est toujours favorable (moins de défauts). Les 16 paires miroir montrent un effet marginal de $+8{,}2$~pp en moyenne, favorable dans 15/16 cas.
\item \textbf{Chômage}~: une hausse est toujours adverse. Les 16 paires miroir montrent $-2{,}9$~pp, adverse dans 16/16 cas.
\end{itemize}

Trois variables sont \textbf{ambiguës}~:

\begin{itemize}
\item \textbf{Taux d'intérêt}~: le canal du service de la dette (hausse $=$ adverse, \citealt{duffie2007}) s'oppose au canal de Merton (\textit{risk-free drift}, hausse $=$ favorable).
\item \textbf{Prix immobiliers (HPI)}~: une hausse réduit les défauts hypothécaires (favorable) mais gonfle les bulles (adverse à terme).
\item \textbf{Inflation}~: une hausse réduit la dette réelle (favorable) mais érode les revenus (adverse).
\end{itemize}

\subsection{Espace des conventions}

Avec 3 variables ambiguës à 2 choix chacune, l'espace des conventions contient $2^3 = 8$ configurations. Chaque convention est identifiée par un triplet $(s_{\text{IR}}, s_{\text{infl}}, s_{\text{HPI}})$~:

\begin{table}[htbp]
\centering
\caption{Les 8 conventions de signe testées}
\label{tab:8-conventions}
\begin{tabular}{lcccc}
\toprule
\textbf{Convention} & \textbf{IR} & \textbf{Inflation} & \textbf{HPI} & \textbf{Statut} \\
\midrule
A & Favorable & Favorable & Favorable & Éliminée \\
\textbf{B (production)} & \textbf{Adverse} & \textbf{Favorable} & \textbf{Favorable} & \textbf{Retenue} \\
C (miroir Taylor) & Favorable & Adverse & Favorable & Viable \\
D & Adverse & Adverse & Favorable & Éliminée \\
E--H & $\cdot$ & $\cdot$ & Adverse & Éliminées (4) \\
\bottomrule
\end{tabular}
\end{table}


\section{Protocole expérimental}
\label{sec:conv-protocole}

Chaque convention est testée sur les 12 scénarios prédéfinis (chapitre~\ref{chap:metriques}), pour un total de $8 \times 12 = 96$ exécutions complètes du pipeline (DGP $\to$ PD $\to$ ECL $\to$ RAROC $\to$ BL-CVaR). L'implémentation utilise un \textit{monkey-patch} de la fonction \texttt{macro\_to\_z}~: les signes sont inversés dynamiquement sans modifier le code de production.

\subsection{Critères d'évaluation}

Cinq critères composites ($C_1$ à $C_5$) évaluent chaque convention~:

\begin{enumerate}
\item[$C_1$] \textbf{Cohérence adverse}~: aucun scénario adverse ne doit produire un RAROC supérieur au scénario central. \emph{(Binaire, éliminatoire.)}
\item[$C_2$] \textbf{Cohérence favorable}~: aucun scénario favorable ne doit produire un RAROC inférieur au scénario central. \emph{(Binaire, éliminatoire.)}
\item[$C_3$] \textbf{Monotonicité adverse}~: les scénarios adverses doivent être ordonnés par sévérité (GFC $<$ COVID $<$ Souveraine). \emph{(Rang de Spearman.)}
\item[$C_4$] \textbf{Diversification}~: au moins 10/14 classes doivent s'ajuster dans la direction attendue pour chaque scénario. \emph{(Taux d'accord.)}
\item[$C_5$] \textbf{AQS}~: le score de qualité d'allocation (équation~\ref{eq:aqs}) doit être positif pour les 12 scénarios. \emph{(AQS $> 0$.)}
\end{enumerate}

Le score composite est~:

\begin{equation}
C_{\text{composite}} = \frac{C_1 + C_2 + C_3 + C_4 + C_5}{5}
\label{eq:composite}
\end{equation}


\section{Résultats~: élimination en 3 étapes}
\label{sec:conv-resultats}

\subsection{Étape 1~: élimination par HPI adverse (4 conventions)}

Les 4 conventions avec HPI adverse (E--H) sont éliminées car elles produisent $C_2 = 33\%$~: le scénario Boom immobilier (HPI $= +8\%$) est traité comme adverse, ce qui pousse le RAROC en dessous du central --- un résultat absurde car la hausse des prix immobiliers est manifestement favorable au portefeuille hypothécaire.

\textbf{L'espace viable se réduit à $2^2 = 4$ conventions} (A, B, C, D).

\subsection{Étape 2~: élimination par incohérence adverse (2 conventions)}

\paragraph{Convention A éliminée.} IR favorable fait que la Stagflation (taux $= 7\%$, PIB $= -2\%$) produit un RAROC de $+7{,}9\%$ (le moteur interprète la hausse des taux comme un signal favorable, compensant la récession). Résultat~: $C_1 = 75\%$ (3/4 adverses cohérents, la Stagflation échoue). Le score composite est de 84\%.

\paragraph{Convention D éliminée.} IR adverse + inflation adverse sature le facteur $Z$ dans les crises combinées. La GFC produit un RAROC de $+3{,}0\%$ (au lieu de négatif)~: les deux canaux adverses se compensent partiellement de manière non physique. Score composite~: 81\%.

\subsection{Étape 3~: deux survivantes, miroir de Taylor}

Les conventions B et C survivent avec des scores proches~:

\begin{table}[htbp]
\centering
\caption{Comparaison des deux conventions survivantes}
\label{tab:bc-compare}
\begin{tabular}{lcc}
\toprule
\textbf{Critère} & \textbf{B (IR adv., infl.\ fav.)} & \textbf{C (IR fav., infl.\ adv.)} \\
\midrule
$C_1$ (adverse) & 100\% & 100\% \\
$C_2$ (favorable) & 100\% & 100\% \\
$C_3$ (monotonicité) & 100\% & 100\% \\
$C_4$ (diversification) & 92\% & 94\% \\
$C_5$ (AQS) & 86\% & 100\% \\
\midrule
$C_{\text{composite}}$ & \textbf{97,1\%} & \textbf{99,0\%} \\
AQS moyen & 0,9404 & 0,9283 \\
Anomalies & 0 & 0 \\
\bottomrule
\end{tabular}
\end{table}

L'écart de 1,9~pp ($97{,}1\%$ vs $99{,}0\%$) n'est \textbf{pas significatif}~: les deux conventions sont structurellement indistinguables.

\subsection{Le miroir de Taylor}

B et C sont des images miroir via la corrélation Taylor (IR/inflation, $r \approx +0{,}60$)~: dans un modèle mono-facteur ASRF, inverser simultanément les signes de deux variables fortement corrélées produit des $Z$ quasi-identiques. Les deux conventions ne divergent que lorsque la règle de Taylor est violée~:

\begin{itemize}
\item \textbf{Stagflation} (taux$\uparrow$ + inflation$\uparrow$ + PIB$\downarrow$)~: Taylor violée. B produit $-9{,}8\%$, C produit $-14{,}3\%$. Delta = 4,5~pp.
\item \textbf{Déflation importée} (taux$\downarrow$ + inflation$\downarrow$ + PIB$\uparrow$)~: Taylor violée. B produit $+14{,}6\%$, C produit $+19{,}0\%$. Delta = 4,4~pp.
\item \textbf{Central, Boom immobilier}~: Taylor respectée. Delta $\leq 0{,}2$~pp.
\end{itemize}

\subsection{Interaction IR $\times$ Inflation~: exactement nulle}

L'analyse factorielle confirme que l'interaction IR$\times$Inflation est \textbf{exactement 0,00~pp} sur les 8 groupes $\times$ 12 scénarios. Les effets sont parfaitement additifs, conséquence de la linéarité de \texttt{macro\_to\_z}. Il n'y a pas de «~miroir de Taylor~» au sens d'une interaction non-linéaire --- seule une co-directionnalité~: IR et inflation vont dans la même direction dans 10/12 scénarios.


\section{Swing maximal et sensibilité}
\label{sec:conv-swing}

Le swing maximal entre les 8 conventions est de \textbf{17,8~pp} (Stagflation, convention~A vs~D). Ce swing est de premier ordre~: il dépasse l'amplitude de tout paramètre d'optimisation (volatilités, poids typiques, seuils SICR).

\begin{table}[htbp]
\centering
\caption{Swing inter-conventions par scénario (pp de RAROC)}
\label{tab:swing-conv}
\begin{tabular}{lcccc}
\toprule
\textbf{Scénario} & \textbf{8 conv.} & \textbf{$C_2$=100\% (4)} & \textbf{B vs C} & \textbf{Risque} \\
\midrule
Central & 0,0 & 0,0 & 0,0 & Nul \\
GFC & 8,5 & 5,4 & 1,4 & Modéré \\
Souveraine & 7,0 & 6,7 & 3,6 & Élevé \\
\textbf{Stagflation} & \textbf{9,7} & \textbf{9,6} & \textbf{4,5} & \textbf{Maximum} \\
COVID & 8,5 & 8,2 & 5,0 & Élevé \\
Reprise & 4,0 & 3,0 & 0,8 & Faible \\
Hypercroissance & 7,4 & 1,0 & 0,2 & Nul \\
Boom immobilier & 3,0 & 2,7 & 0,0 & Nul \\
Trappe à liquidité & 7,5 & 7,3 & 4,2 & Élevé \\
Déflation importée & 5,1 & 4,8 & 4,4 & Élevé \\
Transition climatique & 7,7 & 7,6 & 2,9 & Modéré \\
\bottomrule
\end{tabular}
\end{table}

\paragraph{Observation clé.} L'Hypercroissance chute de 7,4~pp (8 conv.) à 1,0~pp (4 viables)~: le swing massif était causé par les conventions non-viables (HPI adverse). Après élimination, la convention n'est plus un risque de premier ordre pour ce scénario. En revanche, la Stagflation reste à 9,6~pp même parmi les 4 viables~: c'est le scénario le plus convention-sensible.


\section{Effet marginal par variable}
\label{sec:conv-marginal}

L'analyse par paires miroir (chaque variable basculée F$\leftrightarrow$A, les 4 autres fixées) quantifie la contribution marginale de chaque variable~:

\begin{table}[htbp]
\centering
\caption{Effet marginal par variable macro (16 paires miroir)}
\label{tab:marginal-conv}
\begin{tabular}{lccc}
\toprule
\textbf{Variable} & $\Delta$\textbf{(F--A) moyen} & \textbf{F $>$ A} & \textbf{Verdict} \\
\midrule
PIB & $+8{,}2$~pp & 15/16 & Non ambigu --- favorable \\
Chômage & $-2{,}9$~pp & 0/16 & Non ambigu --- adverse \\
Taux d'intérêt & $+6{,}9$~pp & 15/16 & Ambigu, favorable domine \\
Prix immobiliers & $+5{,}9$~pp & 16/16 & Ambigu, favorable domine \\
Inflation & $+1{,}2$~pp & 7/16 & Ambigu, quasi-indifférent \\
\bottomrule
\end{tabular}
\end{table}

L'inflation est le paramètre le moins discriminant~: favorable dans seulement 7/16 paires, avec un effet moyen de 1,2~pp. Les conventions 02 (inflation favorable) et 10 (inflation adverse) sont ex aequo au Tier~1 (96\%). Ce résultat explique pourquoi B et C sont indistinguables~: la seule variable qui les sépare (inflation) a un pouvoir discriminant quasi-nul.


\section{Décision et justification}
\label{sec:conv-decision}

La convention B (IR adverse, inflation favorable, HPI favorable) est maintenue en production, malgré un score composite inférieur de 1,9~pp à la convention C. La justification est triple~:

\begin{enumerate}
\item \textbf{Justification économique}~: le canal du service de la dette domine empiriquement le canal de Merton pour les banques commerciales \citep{duffie2007}. Le rapport EBA~2023 sur la transmission des taux confirme que les hausses de taux augmentent les défauts corporate à horizon 1--3 ans.
\item \textbf{Non-significativité de l'écart}~: $\Delta C_{\text{composite}} = 1{,}9$~pp, avec un test de Spearman $p < 0{,}02$ pour les deux conventions. L'écart est dans le bruit.
\item \textbf{Risque de régression}~: changer de convention invaliderait tous les tests de cohérence économique, les calibrations de seuil SICR et les narratifs CRO. Le gain marginal ne justifie pas le risque opérationnel.
\end{enumerate}

\begin{tcolorbox}[colback=bleutitre!5, colframe=bleutitre, title=\textbf{Résultat central}]
Dans un cadre à facteur systématique linéaire et signe global, le risque de spécification macro domine le risque d'optimisation d'un facteur 4 à 9 (swing de 17,8~pp vs.\ 2--5~pp). Ce résultat constitue une borne supérieure~: le swing se réduirait avec un signe classe-spécifique ou un modèle non-linéaire, mais la dominance de la spécification sur l'optimisation resterait. \textbf{Toute allocation présentée sans spécifier et justifier sa modélisation macro est incomplète.}
\end{tcolorbox}


\section{Implications pour la pratique bancaire}
\label{sec:conv-implications}

Ce résultat a des implications directes pour les praticiens~:

\begin{enumerate}
\item \textbf{Transparence}~: tout modèle IFRS~9 utilisant un facteur systématique doit documenter sa convention de signe. L'absence de documentation constitue un risque de modèle de premier ordre.
\item \textbf{Robustesse}~: les résultats d'allocation devraient être présentés avec un intervalle de convention (B et C), pas une estimation ponctuelle. L'intervalle de 4--5~pp sur la Stagflation est irréductible.
\item \textbf{Validation}~: le back-testing d'un modèle ECL doit inclure la convention de signe dans les facteurs testés, au même titre que les paramètres PD ou LGD.
\item \textbf{Limites de l'optimisation}~: consacrer des ressources à affiner l'optimiseur ($\sim$2~pp de gain) avant de trancher la convention ($\sim$18~pp de swing) est une erreur de priorité.
\end{enumerate}


\section{Portée du résultat et limites structurelles}
\label{sec:conv-portee}

Les sections précédentes ont établi que le swing de convention atteint 17,8~pp de RAROC, soit un facteur 4 à 9 au-dessus du swing d'optimisation. Ce résultat est correct dans le cadre du modèle, mais trois propriétés structurelles de ce cadre amplifient mécaniquement le swing. Les expliciter est nécessaire pour interpréter correctement la portée du résultat et sa transférabilité aux données réelles.


\subsection{Le swing est amplifié par la linéarité de \texttt{macro\_to\_z}}
\label{sec:conv-linearite}

La transformation macro$\to Z$ est une combinaison linéaire à coefficients signés~:
%
\begin{equation}
Z = \sum_{i=1}^{5} s_i \cdot w_i \cdot \delta_i
\label{eq:z-lineaire}
\end{equation}
%
où $s_i \in \{-1, +1\}$ est la convention de signe, $w_i$ la sensibilité sectorielle et $\delta_i$ le choc macro normalisé. Le signe $s_i$ est un \textbf{multiplicateur binaire}~: basculer $s_{\text{IR}}$ de $+1$ à $-1$ revient à inverser \emph{intégralement} la contribution des taux au facteur $Z$.

L'amplitude de cette inversion est directement proportionnelle au produit $w_{\text{IR}} \cdot |\delta_{\text{IR}}|$. Pour la Stagflation ($\delta_{\text{IR}} = +2{,}3\sigma$, $w_{\text{IR}} = 1{,}2$ en moyenne pondérée), l'inversion du signe IR produit un basculement de $\Delta Z = 2 \times 1{,}2 \times 2{,}3 = 5{,}5$ unités de $Z$, ce qui se traduit, via la non-linéarité de la fonction de perte ECL, par un swing de RAROC de $\sim$9,6~pp (entre conventions A et B). Pour le scénario Central ($\delta_{\text{IR}} = 0$), le basculement est nul ($\Delta Z = 0$), ce qui explique le swing de 0,0~pp observé dans le tableau~\ref{tab:swing-conv}.

Le swing maximal de 17,8~pp (Stagflation, A vs.\ D) cumule le basculement IR \emph{et} le basculement inflation, qui agissent dans la même direction quand la règle de Taylor est violée.

Un modèle non-linéaire --- par exemple $Z_{\text{IR}} = f(\Delta_{\text{IR}})$ avec $f$ convexe ou à seuil --- produirait un swing \textbf{plus faible}, parce que le basculement ne serait pas symétrique~: l'effet d'une hausse de taux de 300~bp ne serait pas l'image miroir de l'effet d'une baisse de 300~bp. Le swing de 17,8~pp est donc une \emph{borne supérieure} du risque de spécification, pas une estimation centrale.


\subsection{Le signe global vs.\ le signe par classe d'actifs}
\label{sec:conv-signe-global}

Le modèle ASRF mono-facteur impose un signe $s_i$ identique pour les 14 classes d'actifs. Or l'effet économique des taux d'intérêt diffère substantiellement selon la classe. Le tableau~\ref{tab:signe-classe} résume le canal dominant et le signe probable par classe~:

\begin{table}[htbp]
\centering
\caption{Canal dominant IR et signe probable par classe d'actifs}
\label{tab:signe-classe}
\begin{tabular}{lll}
\toprule
\textbf{Classe} & \textbf{Canal dominant IR} & \textbf{Signe probable} \\
\midrule
Crédit immobilier   & Service de la dette, LTV              & Adverse \\
Crédit consommation & Service de la dette                   & Adverse \\
Crédit corporate    & Service de la dette, refinancement    & Adverse \\
Souverain           & Flight-to-quality, spread compression & Ambigu \\
Trade finance       & Duration $\approx 0$, auto-liquidant  & Neutre \\
Repos / SFT         & Marge de collatéral                   & Faiblement favorable \\
Obligations corp.   & Refinancement, spread                 & Faiblement adverse \\
Private equity      & Compression multiple de sortie        & Adverse \\
Titrisation         & Spread / base, seuils de tranche      & Adverse \\
Financement projet  & Levier fixe, SPV                      & Neutre \\
Actions cotées      & Actualisation, multiples              & Adverse \\
Dérivés             & Mark-to-market, CVA                   & Ambigu \\
\bottomrule
\end{tabular}
\end{table}

Un signe global $s_{\text{IR}} = -1$ (adverse) est cohérent avec les classes dominées par le canal du service de la dette (crédit immobilier, consommation, corporate, PE), mais pas avec les classes pour lesquelles le canal est neutre ou ambigu (trade finance, repos, souverain). En forçant toutes les classes à réagir dans la même direction, le signe global \textbf{amplifie le swing} au niveau du portefeuille~: les effets se cumulent au lieu de se compenser partiellement.

Avec des signes classe-spécifiques, le swing agrégé serait mécaniquement plus faible, parce que les classes à signe favorable (repos, trade finance) compenseraient partiellement les classes à signe adverse (crédit, PE) au niveau du portefeuille. L'analyse par force brute deviendrait impraticable ($2^{14} = 16\,384$ combinaisons pour IR seul), mais le protocole de test (critères $C_1$ à $C_5$) resterait applicable sur un espace réduit par des contraintes économiques \textit{a priori} (le tableau ci-dessus éliminerait la majorité des configurations).

Nous ne refaisons pas l'analyse avec des signes classe-spécifiques, car ce serait circulaire sur données synthétiques~: le DGP impose un signe implicite par construction, et estimer un signe à partir de données que l'on a générées avec ce signe est tautologique. Le seul moyen non-circulaire serait d'utiliser des données réelles (cf.\ section~\ref{sec:conv-transferabilite}).


\subsection{Transférabilité aux données réelles}
\label{sec:conv-transferabilite}

Si une banque appliquait le protocole du chapitre~\ref{chap:conventions} sur son portefeuille réel, trois changements structurels se produiraient.

\paragraph{Le signe deviendrait une estimation empirique.}
Au lieu d'être un paramètre libre à 8 valeurs possibles, le signe de l'effet des taux serait estimé par régression (PD $\sim$ variables macro, panel 5--20 ans). L'espace des conventions se réduirait de 8 à 1 (l'estimation ponctuelle), avec un intervalle de confiance. Le swing résiduel serait l'incertitude d'estimation --- typiquement l'écart-type du coefficient de régression --- et non l'espace complet des conventions. Pour des séries de 15 ans à fréquence trimestrielle (60 observations), l'incertitude d'estimation sur le signe de l'effet IR serait de l'ordre de $\pm 0{,}3\sigma$, bien en deçà du basculement binaire $\pm 1$ du modèle actuel.

\paragraph{La forme fonctionnelle serait non-linéaire.}
Le modèle mono-facteur linéaire serait vraisemblablement remplacé par un modèle à seuils ou un GAM (\textit{Generalized Additive Model}). La relation macro$\to$défauts est empiriquement convexe pour les taux d'intérêt (les hausses de taux augmentent les défauts plus que les baisses ne les réduisent, \citealt{duffie2007}) et concave pour le PIB (les récessions augmentent les défauts plus que les expansions ne les réduisent). Ces asymétries réduiraient le swing en brisant la symétrie du basculement de signe.

\paragraph{Le résultat qualitatif resterait robuste.}
Le résultat quantitatif (17,8~pp de swing) est une borne supérieure, amplifiée par la linéarité, le signe global et les corrélations déclaratives. Mais le résultat \emph{qualitatif} --- le risque de spécification domine le risque d'optimisation --- resterait vraisemblablement vrai, parce que l'incertitude sur la \emph{forme fonctionnelle} de la relation macro$\to$défauts est structurellement plus grande que l'incertitude sur les paramètres de l'optimiseur. La première porte sur la structure du modèle (linéaire vs.\ seuils vs.\ GAM, signe global vs.\ classe-spécifique, corrélations stables vs.\ régime-dépendantes)~; la seconde porte sur des paramètres continus dans un espace convexe. L'erreur de structure domine l'erreur de paramètre --- un résultat classique en théorie de la décision \citep{mcneil_frey_embrechts2015}.


% ══════════════════════════════════════════════════════════
% CHAPITRE 11 — Architecture et génie logiciel
% ══════════════════════════════════════════════════════════
\chapter{Architecture logicielle et reproductibilité}
\label{chap:architecture}

Ce chapitre documente l'architecture du système \texttt{ifrs9\_cockpit}, un outil intégrant la totalité du pipeline décrit dans les chapitres précédents~: du DGP synthétique au dashboard CRO, en passant par les modèles PD, le moteur ECL, l'optimiseur BL-CVaR et les 10 moteurs de gouvernance. L'objectif est double~: reproductibilité des résultats et transparence pour l'audit.

\section{Vue d'ensemble}
\label{sec:archi-overview}

Le système est organisé en 7 packages Python~:

\begin{table}[htbp]
\centering
\caption{Packages principaux et leur rôle}
\label{tab:packages}
\begin{tabular}{llr}
\toprule
\textbf{Package} & \textbf{Rôle} & \textbf{Modules} \\
\midrule
\texttt{config/} & Configuration déclarative & 6 \\
\texttt{synthetic\_generator/} & DGP + 14 générateurs & 15 \\
\texttt{training/} & Entraînement PD offline & 3 \\
\texttt{engine/} & ECL, comparator, gouvernance & 18 \\
\texttt{analytics/} & CRO analyst, narratifs & 10 \\
\texttt{dashboard/} & Callbacks, charts, layout & 24 \\
\texttt{utils/} & Helpers, cache, compatibilité & 4 \\
\bottomrule
\end{tabular}
\end{table}

Le pipeline complet s'exécute en moins de 5 secondes ($\sim$1,1~s recalcul ECL, $\sim$3,5~s entraînement PD avec cache). Les modèles PD pré-entraînés sont sérialisés en fichiers \texttt{.joblib} pour éviter le ré-entraînement à chaque démarrage.

\section{Configuration déclarative}
\label{sec:archi-config}

Le package \texttt{config/} contient 6 sous-modules exportant 76 symboles via \texttt{\_\_init\_\_.py}. Le principe est qu'aucune constante métier ne doit être codée en dur dans le code~:

\begin{itemize}
\item \texttt{sectors.py} (630 lignes)~: \texttt{SectorConfig}, \texttt{SECTORS}, \texttt{AssetClassProfile} (22 champs), \texttt{ASSET\_CLASSES} (14 instances).
\item \texttt{scenarios.py} (353 lignes)~: \texttt{MacroScenario}, \texttt{PREDEFINED\_SCENARIOS} (12 scénarios), \texttt{MACRO\_COVARIANCE}.
\item \texttt{models.py} (333 lignes)~: configurations PD, LGD, EAD, IFRS~9, SICR, listes de features.
\item \texttt{basel.py} (162 lignes)~: \texttt{BaselConfig}, \texttt{RiskAppetite}, \texttt{PEClassification}, seuils CRO.
\item \texttt{dashboard.py} (95 lignes)~: palettes, couleurs, formats d'affichage.
\item \texttt{rules.py} (179 lignes)~: règles d'incohérence macro, validation de configuration.
\end{itemize}

Cette séparation permet à un auditeur de vérifier toutes les hypothèses du modèle en lisant un seul répertoire, sans parcourir le code métier.

\section{Générateurs synthétiques}
\label{sec:archi-generators}

Les 14 classes d'actifs sont alimentées par 14 générateurs suivant un patron uniforme~:

\begin{enumerate}
\item \texttt{generate\_*\_positions(n, total\_ead, seed)} $\to$ DataFrame Polars.
\item \texttt{compute\_*\_pd/lgd/rw(df)} $\to$ ajoute les colonnes réglementaires.
\item \texttt{stress\_*\_positions(df, macro)} $\to$ dictionnaire avec PD/LGD/RW stressés.
\item \texttt{aggregate\_*\_to\_balance\_row(df, profile)} $\to$ ligne de bilan (19 clés standard).
\end{enumerate}

Le \texttt{DatasetBundle} stocke les positions brutes de chaque classe, permettant un stress bottom-up avant agrégation au bilan. Les graines aléatoires sont décalées par classe (\texttt{base\_seed + 1000} pour les actions, $+1100$ pour les obligations, etc.) pour garantir l'indépendance statistique.

\section{Moteur de comparaison multi-actifs}
\label{sec:archi-comparator}

Le \texttt{PortfolioComparator} est assemblé par un patron \textit{mixin} en 5 fichiers~:

\begin{itemize}
\item \texttt{crr3.py}~: calcul des pondérations CRR3 par classe (LTV, SEC-SA, ECRA, \textit{slotting}).
\item \texttt{metrics.py}~: 9 méthodes dont \texttt{compute\_raroc\_multiclass} (RAROC unifié 14 classes).
\item \texttt{optimizer.py}~: 10 méthodes (BL-CVaR Phase~1, Phase~2, Anderson, softmax, etc.).
\item \texttt{sensitivity.py}~: Morris, Sobol, AQS.
\item \texttt{\_\_init\_\_.py}~: assemblage final.
\end{itemize}

\section{Dashboard CRO}
\label{sec:archi-dashboard}

Le dashboard est construit avec Dash (Plotly) et s'organise en trois couches~:

\begin{enumerate}
\item \textbf{Layout}~: barre latérale (sliders macro, sélection scénarios) + contenu principal (9 sections, 8 panneaux dépliables, 2 modales de détail).
\item \textbf{Callbacks}~: 25 callbacks répartis en 9 fichiers (scénarios, pipeline, UI, modales, dépliables, constructeurs de contenu, exports).
\item \textbf{Charts}~: 45 fonctions de tracé dans 11 fichiers, utilisant un layout de base commun (\texttt{\_base\_layout()}) et une palette de couleurs unifiée.
\end{enumerate}

Le dashboard supporte 4 formats d'export~: PDF (rapport CRO), Excel (données brutes), JSON (état pipeline) et PNG (graphiques individuels).

\section{Tests et validation}
\label{sec:archi-tests}

La suite de tests comprend 665 tests organisés en 3 niveaux~:

\begin{itemize}
\item \textbf{Niveau 1 --- rapides} ($\sim$1 minute)~: tests unitaires, validation de schémas, cohérence des configurations. Exécutés à chaque commit.
\item \textbf{Niveau 2 --- complets} ($\sim$52 secondes avec \texttt{-n 4})~: tests d'intégration, validation économique (246 tests dans \texttt{test\_comparator.py}), cohérence inter-scénarios.
\item \textbf{Niveau 3 --- lents} (21 tests, marqués \texttt{@slow})~: tests standalone, reproductibilité, Morris/Sobol (3\,840 évaluations).
\end{itemize}

Les fixtures sont partagées au niveau module (\texttt{orchestrator\_state}) et session (\texttt{global\_pipeline\_results}) pour éviter les recalculs redondants. Le parallélisme utilise \texttt{pytest-xdist} avec 4 workers (optimal sur 32 c\oe urs, au-delà le surcoût des fixtures annule le gain).

\section{Migration Polars}
\label{sec:archi-polars}

Le pipeline interne utilise Polars (types stricts, évaluation paresseuse, parallélisme natif). Les frontières Pandas sont maintenues aux interfaces ML (\texttt{sklearn}, \texttt{TabNet}) via le module \texttt{frame\_compat.py}~:

\begin{itemize}
\item \texttt{to\_pandas(df)} / \texttt{to\_polars(df)}~: conversion bidirectionnelle.
\item \texttt{ensure\_numpy(series)}~: extraction de vecteurs pour les fonctions numériques.
\end{itemize}

Les fonctions de tracé utilisent le support natif Polars de Plotly~6.x (aucune conversion nécessaire).

\section{Reproductibilité}
\label{sec:archi-repro}

La reproductibilité est garantie par~:

\begin{enumerate}
\item \textbf{Graine fixe}~: \texttt{RANDOM\_SEED = 123} pour la génération de données, \texttt{seed = 42} pour l'entraînement (indépendance données/modèle).
\item \textbf{Séparation train/test}~: les données consommateur sont séparées en \texttt{train.parquet} (50k) et \texttt{test.parquet} (15k). Le pipeline de production utilise le holdout.
\item \textbf{Sérialisation}~: \texttt{pd\_suite.joblib}, \texttt{governance\_suite.joblib}, suites consommateur et hypothécaire --- les modèles entraînés sont des artefacts versionnés.
\item \textbf{Git}~: deux branches actives (\texttt{pe-bc} pour la Partie~I, \texttt{14-couches} pour la Partie~II), commits documentés.
\end{enumerate}


% ══════════════════════════════════════════════════════════
% CONCLUSION
% ══════════════════════════════════════════════════════════
\chapter*{Conclusion}
\addcontentsline{toc}{chapter}{Conclusion}
\label{chap:conclusion}

\section*{Synthèse des résultats}

Ce mémoire a construit un cadre intégré d'allocation de capital pour une banque universelle européenne, depuis les données synthétiques jusqu'au dashboard CRO. Les résultats s'organisent en trois niveaux~:

\paragraph{Niveau 1~: l'arbitrage crédit/PE est scénario-dépendant.} Les pondérations CRR3 du private equity (190--400\%) pénalisent fortement son allocation en capital, mais ne la plafonnent pas en présence de capital excédentaire. Dans le modèle à 5 secteurs (Partie~I), le split PE varie de 18\,\% (scénarios neutres) à 62\,\% (crise souveraine), l'optimiseur utilisant le PE comme diversificateur quand le RAROC crédit est négatif. À l'échelle des 14 classes (Partie~II), l'impact de prix logarithmique (capacité de marché de 300~Mds\texteuro{}) contraint davantage l'allocation PE. Le levier crédit/PE est donc réel mais conditionnel au niveau de capital et à la capacité de marché.

\paragraph{Niveau 2~: l'arbitrage sectoriel reste le levier principal.} Le RAROC est déterminé de premier ordre par la répartition sectorielle (sensibilités macro, LGD sectorielles, concentration HHI). En Partie~I, l'allocation sectorielle et le split crédit/PE sont tous deux scénario-réactifs (écart-type moyen de 3{,}6\,\% et 7{,}4\,\% respectivement), avec un alignement RAROC--poids de 71\,\%. En Partie~II, ce résultat est confirmé à l'échelle de 14 classes d'actifs. Les 246 tests de cohérence économique et les 12 scénarios historiquement calibrés valident la robustesse.

\paragraph{Niveau 3~: le risque de spécification macro est de premier ordre.} L'analyse systématique de 96 configurations révèle que la modélisation du facteur systématique --- dont la convention de signe est la forme la plus élémentaire --- produit un swing de 17,8~pp de RAROC, contre 2--5~pp pour l'optimisation. Ce swing est une borne supérieure, amplifiée par la linéarité du modèle et le caractère global du signe. Avec des données empiriques ou un modèle non-linéaire, le swing résiduel serait plus étroit mais resterait vraisemblablement dominant. Ce résultat déplace la priorité~: avant d'affiner l'optimiseur, il faut sécuriser la spécification macro.

\section*{Contributions}

Les contributions de ce travail sont~:

\begin{enumerate}
\item Un \textbf{optimiseur BL-CVaR à zéro paramètre arbitraire}, où la diversification émerge de l'impact de prix (Kyle 1985) et non de pénalités artificielles. Les 22 paramètres discrétionnaires des versions précédentes ont été supprimés ou remplacés par des observables de marché.

\item Un \textbf{protocole de validation des conventions de signe}, réplicable par toute institution utilisant un modèle macro-facteur. Le protocole (8 conventions $\times$ 12 scénarios = 96 runs, 5 critères composites) identifie les conventions viables et quantifie le risque de convention résiduel.

\item Un \textbf{système de gouvernance à 10 moteurs}, combinant des approches complémentaires (conformal prediction, Sobol, HMM, GFlowNet, TDA, VRP, RMT, signatures, compliance gates) pour produire un narratif CRO auditable et multi-perspective.

\item Un \textbf{dashboard CRO interactif} avec 45 graphiques, 12 scénarios, 14 classes d'actifs et 4 formats d'export, conçu pour la prise de décision en comité de risque.
\end{enumerate}

\section*{Limites}

\paragraph{Données synthétiques.} Le DGP v4.5 produit des distributions réalistes mais les amplitudes absolues dépendent de la calibration. Un portefeuille réel produirait des niveaux de RAROC différents, mais les \emph{classements relatifs} (quel scénario est pire que quel autre, quelle classe consomme le plus de capital) seraient préservés.

\paragraph{Liquidité simplifiée.} L'impact de prix utilise un modèle logarithmique statique. Un modèle dynamique (\citealt{almgren2001}) avec carnets d'ordres et résilience temporelle affinerait les résultats pour les classes illiquides (PE, financement de projet).

\paragraph{Horizon statique.} Le modèle est à horizon unique (1 an). Un modèle dynamique avec origination, \textit{run-off} et roulement des positions capturerait les effets de stock vs flux, particulièrement importants pour le crédit immobilier (ténor 20 ans, duration effective 4 ans).

\paragraph{Corrélations déclaratives.} La matrice de covariance macro est déclarative (calibrée par expertise). Une estimation empirique sur données historiques (séries temporelles BCE/Eurostat) renforcerait la crédibilité, au prix d'une instabilité temporelle des corrélations.

\paragraph{Linéarité et binarité du facteur systématique.} Le modèle \texttt{macro\_to\_z} est linéaire en les variables macro, et la convention de signe est binaire ($\pm 1$). Cette architecture maximise la sensibilité à la convention~: toute inversion de signe produit un basculement symétrique du facteur $Z$. Un modèle à forme fonctionnelle libre (GAM, splines, seuils) ou à signe classe-spécifique produirait un swing résiduel plus faible. Le swing de 17,8~pp doit donc être interprété comme une borne supérieure du risque de spécification, non comme une estimation de l'incertitude en pratique.

\section*{Perspectives}

Trois extensions naturelles émergent~:

\begin{enumerate}
\item \textbf{Corrélations empiriques.} Remplacer la matrice de covariance déclarative par une estimation DCC-GARCH (\citealt{engle2002}) sur les séries macro BCE/Eurostat. Cela permettrait de tester la stabilité des conventions de signe à travers les régimes historiques (pré-crise, post-GFC, COVID, normalisation 2022--2025).

\item \textbf{Modèle dynamique multi-périodes.} Étendre le cadre à un horizon de 3--5 ans avec origination, run-off et roulement. Cela nécessite un modèle de volume (projection des encours) couplé au moteur ECL, avec des contraintes LCR/NSFR à chaque période.

\item \textbf{Calibration sur données réelles.} Appliquer le protocole de convention de signe à un portefeuille bancaire réel (sous réserve de confidentialité). L'enjeu est de vérifier si les conclusions qualitatives (sectoriel $>$ actif, convention $>$ optimisation) survivent aux données empiriques.
\end{enumerate}

\vspace{1cm}

En définitive, ce mémoire montre que la question «~quelle est l'allocation optimale~?~» est mal posée sans spécifier la convention de signe. L'allocation n'est pas une propriété du portefeuille~: c'est une propriété conjointe du portefeuille \emph{et} de la convention. Reconnaître cette contingence n'affaiblit pas l'analyse --- elle la rend honnête.


\bibliographystyle{plainnat}
\bibliography{bibliographie}

\end{document}
